%% This is a LaTeX document.  Hey, Emacs, -*- latex -*- , get it?
\documentclass[12pt,a4paper]{article}
\usepackage[a4paper,hmargin=2cm,vmargin=3cm]{geometry}
\usepackage[shorthands=off,french]{babel}
\usepackage[utf8]{inputenc}
\usepackage[T1]{fontenc}
\usepackage{lmodern}
\usepackage{newtxtext}
% A tribute to the worthy AMS:
\usepackage{amsmath}
\usepackage{amsfonts}
\usepackage{amssymb}
\usepackage{amsthm}
%
\usepackage{mathrsfs}
\usepackage{wasysym}
\usepackage{url}
\usepackage{mathpartir}
\usepackage{flagderiv}
%
\usepackage{graphics}
\usepackage[usenames,dvipsnames]{xcolor}
\usepackage{tikz}
\usetikzlibrary{arrows}
%
\theoremstyle{definition}
\newtheorem{comcnt}{Tout}[section]
\newcommand\thingy{%
\refstepcounter{comcnt}\medskip\noindent\textbf{\thecomcnt.} }
\newcommand\exercice{%
\refstepcounter{comcnt}\bigskip\noindent\textbf{Exercice~\thecomcnt.}}
%
\newcommand{\dbllangle}{\mathopen{\langle\!\langle}}
\newcommand{\dblrangle}{\mathclose{\rangle\!\rangle}}
\newcommand{\dottedlimp}{\mathbin{\dot\Rightarrow}}
\newcommand{\dottedland}{\mathbin{\dot\land}}
\newcommand{\dottedlor}{\mathbin{\dot\lor}}
\newcommand{\dottedtop}{\mathord{\dot\top}}
\newcommand{\dottedbot}{\mathord{\dot\bot}}
\newcommand{\dottedneg}{\mathop{\dot\neg}}
%
\DeclareUnicodeCharacter{00A0}{~}
%
\newif\ifcorrige
\corrigetrue
\newenvironment{corrige}%
{\ifcorrige\relax\else\setbox0=\vbox\bgroup\fi%
\smallbreak\footnotesize\noindent{\underbar{\textit{Corrigé.}}\quad}}
{{\hbox{}\nobreak\hfill\checkmark}%
\ifcorrige\relax\else\egroup\fi\par}
%
%
%
\begin{document}
\ifcorrige
\title{Logique et Fondements de l'Informatique\\Exercices corrigés}
\else
\title{Logique et Fondements de l'Informatique\\Exercices}
\fi
\author{David A. Madore}
\maketitle

\centerline{\textbf{CSC-3TC34-TP / INF110}}

\vskip2cm

{\footnotesize
\immediate\write18{sh ./vc > vcline.tex}
\begin{center}
Git: \input{vcline.tex}
\\
(Recopier la ligne ci-dessus dans tout commentaire sur ce document)
\end{center}
\immediate\write18{echo ' (stale)' >> vcline.tex}
\par}

\pretolerance=8000
\tolerance=50000


%
%
%

\section{Calculabilité}


\setcounter{comcnt}{-1}\exercice\ (${\star}$)\par\nobreak

Pour chacune des fonctions partielles suivantes (des entiers naturels
vers les entiers naturels), dire si elle est totale ou non, calculable
ou non :
\begin{itemize}
\item[(a)] La fonction constante égale à $0$.
\item[(b)] La fonction valant $0$ en chaque $n\neq 0$ et $1$ en $0$.
\item[(c)] La fonction valant $0$ en chaque $n\neq 0$ et non définie
  en $0$.
\item[(d)] La fonction $n \mapsto A(n,n,n)$, où $(m,n,k) \mapsto
  A(m,n,k)$ est la fonction d'Ackermann
  (cf. exercice \ref{exercise-ackermann-indicator-is-p-r} pour sa
  définition).
\item[(e)] La fonction $\langle e,x\rangle \mapsto \varphi_e(x)$.
\item[(f)] La fonction qui à $e$ associe $1$ si
  $\varphi_e(0)\downarrow$ et $0$ si $\varphi_e(0)\uparrow$.
\item[(g)] La fonction qui à $e$ associe $1$ si
  $\varphi_e(0)\downarrow$ et non définie si $\varphi_e(0)\uparrow$.
\item[(h)] La fonction qui à $e$ associe $0$ si $\varphi_e(0)\uparrow$
  et non définie si $\varphi_e(0)\downarrow$.
\end{itemize}

\begin{corrige}
(a) est totale (définie partout) ; elle est de plus calculable, parce
  qu'il est évident d'écrire un programme (en gros « \texttt{return
    0;} ») qui renvoie toujours $0$.  (Avec la définition et
  numérotation utilisée en cours, on peut d'ailleurs même dire
  précisément que pour $e = \dbllangle 1,1,0\dblrangle = 23$ la
  fonction constante égale à $0$ est la fonction $\varphi_e$.)

(b) est totale (définie partout) ; elle est de plus calculable, parce
  qu'il est évident d'écrire un programme qui teste si son argument
  vaut $0$, renvoie $1$ si c'est le cas et $0$ sinon (en gros
  « \texttt{if (n==0) return 1; else return 0;} »).

(c) n'est pas totale puisque sa définition même dit qu'elle n'est pas
  définie en $0$.  Elle est (partielle) calculable, parce qu'il est
  évident d'écrire un programme qui teste si son argument vaut $0$,
  renvoie $0$ si ce n'est pas le cas, et fait une boucle infinie si
  c'est le cas (en gros « \texttt{if (n!=0) return 0; else while
    (1);} »).

(d) La fonction d'Ackermann est totale (une récurrence sur $k$, et
  pour $k$ fixé sur $n$, montre que $A(m,n,k)$ a bien un sens), donc
  la fonction $n \mapsto A(n,n,n)$ est aussi totale.  De plus, la
  fonction d'Ackermann est calculable par un algorithme récursif
  imitant directement sa définition, et on a vu en cours qu'on pouvait
  utiliser la récursion.  Donc la fonction d'Ackermann et \textit{a
    fortiori} la fonction $n \mapsto A(n,n,n)$ sont calculables.

(e) La fonction $\langle e,x\rangle \mapsto \varphi_e(x)$ est
  calculable par l'existence d'un algorithme universel.  Elle n'est
  pas totale car il existe des $(e,x)$ tels que $\varphi_e(x)\uparrow$
  (par exemple un programme qui fait une boucle infinie).

(f) La fonction est totale par définition (en tout $e$, elle vaut soit
  $0$ soit $1$, mais elle est toujours définie).  Il s'agit ici de la
  fonction indicatrice de $H := \{e \in \mathbb{N} :
  \varphi_e(0)\downarrow\}$.  Ce $H$ est indécidable par une des
  variations du problème de l'arrêt (ou par le théorème de Rice),
  c'est-à-dire justement que sa fonction indicatrice n'est pas
  calculable.  On a donc affaire à une fonction totale non calculable.

(g) La fonction n'est pas totale car il existe des $e$ tels que
  $\varphi_e(0)\uparrow$ (par exemple un programme qui fait une boucle
  infinie).  Elle est (partielle) calculable car il est possible
  d'écrire un programme qui prend $e$ en entrée, simule l'exécution de
  $e$ sur l'entrée $0$ (en utilisant la machine universelle, comme
  en (e)), et renvoie $1$ toujours si l'exécution termine (tandis que
  si l'exécution ne termine pas, elle n'a, bien sûr, pas d'autre choix
  que de ne pas terminer).  De façon équivalente, on peut dire que
  c'est la fonction « semi-indicatrice » du $H$ défini en (f), et ce
  $H$ est semi-calculable comme variante du problème de l'arrêt.

(h) La fonction n'est pas totale car il existe des $e$ tels que
  $\varphi_e(0)\downarrow$.  Montrons par l'absurde qu'elle n'est pas
  calculable : s'il existait $h$ tel que $\varphi_h(e)$ soit la
  fonction dont on parle (c'est-à-dire $0$ si $\varphi_e(0)\uparrow$
  et non définie si $\varphi_e(0)\downarrow$), on pourrait décider le
  problème de l'arrêt de la manière suivante : donné un programme $e$,
  on lance en parallèle l'exécution de $\varphi_e(0)$ et de
  $\varphi_h(e)$, et forcément exactement l'un des deux va terminer
  (si $\varphi_e(0)\downarrow$ c'est qu'elle termine, et si
  $\varphi_e(0)\uparrow$ c'est que $\varphi_h(e)\downarrow$ par
  définition), donc on a juste à voir lequel termine ; or on sait que
  le problème de l'arrêt n'est pas décidable, donc notre hypothèse est
  absurde (ce $h$ n'existe pas).  De façon équivalente, on peut dire
  que notre fonction est (à un décalage de $1$ près) la fonction
  « semi-indicatrice » du complémentaire de $H$ défini en (f), et ce
  $H$ est semi-calculable non calculable donc son complémentaire n'est
  pas calculable.
\end{corrige}

%

\exercice\ (${\star}{\star}$)\par\nobreak

On considère la fonction $f\colon \mathbb{N} \to \mathbb{N}$ qui à $n
\in \mathbb{N}$ associe le $n$-ième chiffre de l'écriture décimale de
$\sqrt{2} \approx 1.41421356237309504880\ldots$, c'est-à-dire $f(0) =
1$, $f(1) = 4$, $f(2) = 1$, $f(3) = 4$, etc.

La fonction $f$ est-elle calculable ?  Est-elle primitive récursive ?
On expliquera précisément pourquoi.

\begin{corrige}
On peut calculer $f(n)$ selon l'algorithme suivant : calculer $N =
10^n$, puis pour $i$ allant de $0$ à $2N$, tester si $i^2 \leq 2 N^2 <
(i+1)^2$ : lorsque c'est le cas (et ce sera le cas pour exactement un
$i$ dans l'intervalle), renvoyer le reste $i\% 10$ de la division
euclidienne de $i$ par $10$.

Cet algorithme est correct car l'inégalité $i^2 \leq 2 N^2 < (i+1)^2$
testé équivaut à $\frac{i}{N} \leq \sqrt{2} < \frac{i+1}{N}$, ce qui
se produit pour exactement un $i$, à savoir $\lfloor \sqrt{2}\times
10^n \rfloor$ (on peut arrêter la boucle à $2N$ car $\sqrt{2} < 2$),
et que le dernier chiffre décimal $i\% 10$ de ce $i$ est le $n$-ième
chiffre de l'écriture décimale de $\sqrt{2}$.

D'autre part, comme on a donné un algorithme explicite, $f$ est
calculable.  Mieux : comme la boucle utilisée est bornée \textit{a
  priori}, $f$ est primitive récursive.
\end{corrige}

%

\exercice\ (${\star}$)\par\nobreak

Supposons que $A \subseteq B \subseteq \mathbb{N}$.  \textbf{(1)} Si
$B$ est décidable, peut-on conclure que $A$ est décidable ?
\textbf{(2)} Si $A$ est décidable, peut-on conclure que $B$ est
décidable ?

\begin{corrige}
La réponse est non dans les deux cas : pour le voir appelons $H := \{e
\in \mathbb{N} : \varphi_e(0)\downarrow\}$ (disons) : il est
indécidable par une des variations du problème de l'arrêt (ou par le
théorème de Rice).  Le fait que $H \subseteq \mathbb{N}$ avec
$\mathbb{N}$ décidable réfute (1), et le fait que $\varnothing
\subseteq H$ avec $\varnothing$ décidable réfute (2).
\end{corrige}

%

\exercice\label{exercise-computable-image-is-semidecidable}\ (${\star}{\star}$)\par\nobreak

\textbf{(1)} Soit $f\colon \mathbb{N} \to \mathbb{N}$ totale
calculable.  Montrer que l'image $f(\mathbb{N})$ (c'est-à-dire $\{f(i)
: i\in\mathbb{N}\}$) est semi-décidable.

\textbf{(2)} Soit $f\colon \mathbb{N} \to \mathbb{N}$ totale
calculable et strictement croissante.  Montrer que l'image
$f(\mathbb{N})$ (c'est-à-dire $\{f(i) : i\in\mathbb{N}\}$) est
décidable.

\begin{corrige}
\textbf{(1)} L'algorithme évident suivant semi-décide $\{f(i) :
i\in\mathbb{N}\}$ : donné $m \in \mathbb{N}$ l'entier à tester, faire
une boucle infinie sur $i$ parcourant les entiers naturels et pour
chacun, tester si $f(i) = m$ : si c'est le cas, terminer et répondre
« oui », sinon, continuer la boucle.

\textbf{(2)} L'algorithme évident suivant décide $\{f(i) :
i\in\mathbb{N}\}$ : donné $m \in \mathbb{N}$ l'entier à tester, faire
une boucle pour $i$ parcourant les entiers naturels, et pour chacun,
tester si $f(i) = m$ : si c'est le cas, terminer et répondre « oui »,
tandis que si $f(i) > m$, terminer et répondre « non », sinon,
continuer la boucle.  La boucle termine en temps fini car $f(i) \geq
i$ (inégalité claire pour une fonction $\mathbb{N} \to \mathbb{N}$
strictement croissante) et notamment la boucle s'arrêtera au pire
lorsque $i$ vaut $m+1$.  (Du coup, si on préfère, on peut réécrire la
boucle potentiellement infinie comme une boucle pour $i$ allant de $0$
à $m$.)
\end{corrige}

%

\exercice\ (${\star}$)\par\nobreak

Montrer que l'ensemble des $e\in \mathbb{N}$ tels que
$\varphi^{(1)}_e(0) = \varphi^{(1)}_e(1)$ (rappel : ceci signifie que
\emph{soit} $\varphi^{(1)}_e(0) \downarrow$ et $\varphi^{(1)}_e(1)
\downarrow$ et $\varphi^{(1)}_e(0) = \varphi^{(1)}_e(1)$, \emph{soit}
$\varphi^{(1)}_e(0) \uparrow$ et $\varphi^{(1)}_e(1) \uparrow$) n'est
pas décidable.

\begin{corrige}
L'ensemble $F$ des fonctions partielles calculables $f\colon
\mathbb{N} \dasharrow \mathbb{N}$ telles que $f(0) = f(1)$ n'est ni
vide (la fonction totale constante de valeur $0$ est dans $F$) ni
plein (la fonction totale identité n'est pas dans $F$).  D'après le
théorème de Rice, l'ensemble des $e$ tels que $\varphi^{(1)}_e \in F$
est indécidable : c'est exactement ce qui était demandé.
\end{corrige}

%

\exercice\label{exercise-image-computable-partial-function}\ (${\star}{\star}{\star}$)\par\nobreak

\textbf{(1)} Soit $B \subseteq \mathbb{N}$ semi-décidable et non-vide.
Montrer qu'il existe $f\colon \mathbb{N} \to \mathbb{N}$ totale
calculable telle que $f(\mathbb{N}) = B$.

(\emph{Indication :} si $m_0 \in B$ et si $B$ est semi-décidé par le
$e$-ième programme, i.e., $B = \{m : \varphi_e(m)\downarrow\}$, on
définira $\tilde f\colon \mathbb{N}^2 \to \mathbb{N}$ par $\tilde
f(n,m) = m$ si $T(n,e,\dbllangle m\dblrangle)$, où $T(n,e,v)$ est
comme dans le théorème de la forme normale de Kleene\footnote{Rappel :
  c'est-à-dire que $T(n,e,\dbllangle \underline{x}\dblrangle)$
  signifie : « $n$ est le code d'un arbre de calcul de
  $\varphi_e(\underline{x})$ termine » (le résultat
  $\varphi_e(\underline{x})$ du calcul étant alors noté $U(n)$).}, et
$\tilde f(n,m) = m_0$ sinon.  Alternativement, si on préfère raisonner
sur les machines de Turing : si $B$ est semi-décidé par la machine de
Turing $\mathscr{M}$, on définit $\tilde f(n,m) = m$ si $\mathscr{M}$
termine sur l'entrée $m$ en $\leq n$ étapes d'exécution, et $\tilde
f(n,m) = m_0$ sinon.)

\textbf{(2)} Soit $f\colon \mathbb{N} \dasharrow \mathbb{N}$ partielle
calculable.  Montrer que l'image $f(\mathbb{N})$ (c'est-à-dire $\{f(i)
: i\in\mathbb{N} \text{~et~} f(i){\downarrow}\}$) est semi-décidable.
(\emph{Indication :} chercher à formaliser l'idée de lancer les
calculs des différents $f(i)$ « en parallèle ».)

\begin{corrige}
\textbf{(1)} La fonction $\tilde f \colon \mathbb{N}^2 \to \mathbb{N}$
définie dans l'indication est calculable (et d'ailleurs même primitive
récursive) : si on a pris la définition avec $T$ le fait que $T$ soit
p.r. fait partie du théorème de la forme normale ; si on préfère les
machines de Turing, c'est le fait qu'on peut simuler l'exécution de
$\mathscr{M}$ pour $n$ étapes (de façon p.r.).  Et on voit qu'on a
$\tilde f(n,m) \in B$ dans tous les cas : donc $\tilde f(\mathbb{N}^2)
\subseteq B$.  Mais réciproquement, si $m \in B$, alors
$\varphi_e(m)\downarrow$ (si on préfère les machines de Turing,
$\mathscr{M}$ termine sur l'entrée $m$), et ceci dit précisément qu'il
existe $n$ tel que $\tilde f(n,m) = m$, donc $m \in \tilde
f(\mathbb{N}^2)$ ; bref, $B \subseteq \tilde f(\mathbb{N}^2)$.  On a
donc $\tilde f(\mathbb{N}^2) = B$ par double inclusion.  Quitte à
remplacer $\tilde f \colon \mathbb{N}^2 \to \mathbb{N}, (n,m) \mapsto
\tilde f(n,m)$ par $f \colon \mathbb{N} \to \mathbb{N}, \langle
n,m\rangle \mapsto \tilde f(n,m)$, on a $f(\mathbb{N}) = B$.

\textbf{(2)} Ici on ne peut pas appliquer bêtement l'algorithme exposé
dans l'exercice \ref{exercise-computable-image-is-semidecidable}
question (1) car si le calcul de $f(i)$ ne termine pas, il bloquera
tous les suivants.  Il faut donc mener le calcul des $f(i)$ « en
parallèle ».  On va procéder par énumération des couples $(n,i)$ et
lancer le calcul de $f(i)$ sur $n$ étapes.

Plus précisément : considérons l'algorithme suivant : il prend en
entrée un entier $m$ dont il s'agit de semi-décider s'il appartient à
$f(\mathbb{N})$.  L'algorithme fait une boucle infinie sur $p$
parcourant les entiers naturels : chaque $p$ est d'abord décodé comme
le code $\langle n,i\rangle$ d'un couple d'entiers naturels (ceci est
bien sûr calculable).  On teste si l'exécution de $f(i)$ termine en
$\leq n$ étapes (ou, si on préfère le théorème de la forme de normale,
on teste si $T(n,e,\dbllangle i\dblrangle)$, où $e$ est un code de la
fonction $f = \varphi^{(1)}_e$) : si oui, et si la valeur $f(i)$
calculée est égale à l'entier $m$ considéré, on termine en renvoyant
« oui », sinon on continue la boucle.

Cet algorithme semi-décide bien $f(\mathbb{N})$ : en effet, dire que
$m \in f(\mathbb{N})$, équivaut à l'existence de $i$ tel que
$f(i){\downarrow} = m$, c'est-à-dire à l'existence de $n,i$ tel que
l'algorithme renverra « oui » en testant $\langle n,i\rangle$.

(\emph{Variante :} plutôt qu'utiliser le codage des couples $\langle
n,i\rangle$, on peut aussi faire ainsi : on parcourt les entiers
naturels $p$ en une boucle infini et pour chacun on effectue deux
boucles bornées pour $0\leq n\leq p$ et $0\leq i\leq p$ : peu
importent les bornes précises, l'important est que pour $p$ assez
grand on va finir par tester le couple $(n,i)$.)
\end{corrige}

%

\exercice\label{exercise-indices-total-functions}\ (${\star}{\star}{\star}$)\par\nobreak

Soit
\[
T := \{e \in \mathbb{N} : \varphi^{(1)}_e\text{~est~totale}\}
\]
l'ensemble des codes des fonctions générales récursives totales
(c'est-à-dire telles que $\forall
n\in\mathbb{N}.\,(\varphi^{(1)}_e(n)\downarrow)$).  On se propose de
montrer que ni $T$ ni son complémentaire $\complement T$ ne sont
semi-décidables.

\textbf{(1)} Montrer en guise d'échauffement que $T$ n'est pas
décidable.

\textbf{(2)} Soit $H := \{d \in \mathbb{N} :
\varphi^{(1)}_d(0)\downarrow\}$ (variante du problème de l'arrêt).
Rappeler brièvement pourquoi $H$ est semi-décidable mais non
décidable, et pourquoi son complémentaire $\complement H$ n'est pas
semi-décidable.

\textbf{(3)} Montrer qu'il existe une fonction $\rho \colon \mathbb{N}
\to \mathbb{N}$ (totale) calculable (d'ailleurs même p.r.) telle que
$\varphi^{(1)}_d(0)\downarrow$ si et seulement si
$\varphi^{(1)}_{\rho(d)}$ est totale (\emph{indication :} on pourra
par exemple construire un programme $e$ qui ignore son argument et qui
simule $d$ sur l'entrée $0$).  Reformuler cette affirmation comme une
réduction.  En déduire que le complémentaire $\complement T$ de $T$
n'est pas semi-décidable.

\textbf{(4)} Montrer qu'il existe une fonction $\sigma \colon
\mathbb{N} \to \mathbb{N}$ (totale) calculable (d'ailleurs même p.r.)
telle que $\varphi^{(1)}_d(0)\downarrow$ si et seulement si
$\varphi^{(1)}_{\sigma(d)}$ \emph{n'est pas} totale
(\emph{indication :} on pourra par exemple construire un programme $e$
qui lance $d$ sur l'entrée $0$ pour un nombre d'étapes donné en
argument, et fait une boucle infinie si cette exécution termine avant
le temps imparti).  Reformuler cette affirmation comme une réduction.
En déduire que $T$ n'est pas semi-décidable.

\begin{corrige}
On notera « $\varphi$ » pour « $\varphi^{(1)}$ » de manière à alléger
les notations.

\textbf{(1)} L'ensemble des fonctions calculables $\mathbb{N}
\dasharrow \mathbb{N}$ qui sont en fait totales ($\mathbb{N} \to
\mathbb{N}$) n'est ni vide (la fonction totale constante de valeur $0$
est dedans) ni plein (la fonction nulle part définie n'est pas
dedans).  D'après le théorème de Rice, l'ensemble $T$ des $e$ tels que
$\varphi_e$ soit totale est indécidable : c'est exactement ce qui
était demandé.

\textbf{(2)} Toujours d'après le théorème de Rice, ou comme variante
du problème de l'arrêt (qui s'y ramène par le théorème s-m-n),
l'ensemble $H$ n'est pas décidable.  Il est cependant semi-décidable
par universalité (on peut lancer l'exécution de $e$ sur l'entrée $0$
et, si elle termine, renvoyer « oui »).  On en déduit que $\complement
H$ n'est pas semi-décidable (car si $H$ et $\complement H$ étaient
semi-décidables, $H$ serait décidable, ce qu'il n'est pas).

\textbf{(3)} Considérons la fonction $\rho$ qui prend en entrée un
programme $d$ (supposé d'un argument) et renvoie le programme $e =:
\rho(d)$ (toujours d'un argument) qui ignore son argument et exécute
$d$ sur l'entrée $0$ : essentiellement par le théorème s-m-n, cette
fonction $\rho$ est totale calculable (d'ailleurs même p.r.).  Par
définition, on a $\varphi_{\rho(d)}(n) = \varphi_d(0)$ (rappelons que
ceci signifie que chacun est défini ssi l'autre l'est et, le cas
échéant, que ces valeurs sont égales).  Notamment, si
$\varphi_d(0)\downarrow$, alors $\varphi_{\rho(d)}$ est totale (et
constante !), tandis que si $\varphi_d(0)\uparrow$, alors
$\varphi_{\rho(d)}$ n'est nulle part définie (donc certainement pas
totale).

Bref, on a construit $\rho\colon\mathbb{N}\to\mathbb{N}$ totale
calculable telle que $d \in H$ si et seulement si $\rho(d) \in T$, ou,
ce qui revient au même, $d \in \complement H$ si et seulement si
$\rho(d) \in \complement T$.  En termes de réductions, ceci signifie
$H \mathrel{\leq_\mathrm{m}} T$, ou, ce qui revient au même,
$\complement H \mathrel{\leq_\mathrm{m}} \complement T$ (le symbole
« $\mathrel{\leq_\mathrm{m}}$ » désignant la réduction many-to-one).
Comme $\complement H$ n'est pas semi-décidable, $\complement T$ ne
l'est pas non plus.

\emph{Remarque :} On n'est pas obligé d'utiliser le terme de
« réduction many-to-one » pour argumenter que $\complement T$ n'est
pas semi-décidable : on peut simplement dire « supposant par l'absurde
que $\complement T$ soit semi-décidable, on pourrait semi-décider
$\complement H$ de la façon suivante : donné $d$, on calcule
$\rho(d)$, on semi-décide si $\rho(d) \in \complement T$ et, si c'est
le cas, on termine en renvoyant “oui” ; or ce n'est pas possible, d'où
une contradiction ».

\textbf{(4)} Considérons la fonction $\sigma$ qui prend en entrée un
programme $d$ (supposé d'un argument) et renvoie le programme $e :=
\sigma(d)$ (toujours d'un argument) défini ainsi : le programme $e$
prend en entrée un nombre $n$ et exécute le programme $d$ sur l'entrée
$0$ pendant $\leq n$ étapes (mettons que ce soient des machines de
Turing, sinon remplacer cet argument par une recherche d'arbre de
calcul parmi les entiers naturels de $0$ à $n$) : si cette exécution a
terminé en $\leq n$ étapes, alors $d$ effectue une boucle infinie,
sinon $d$ termine (et renvoie, disons, $1729$).

Il n'y a pas de difficulté à coder ce programme $e$ (on rappelle
qu'exécuter un programme donné sur $\leq n$ étapes est calculable,
d'ailleurs même primitif récursif), et de plus la fonction $\sigma$
transformant $d$ en $e$ est elle-même calculable (et d'ailleurs elle
aussi primitive récursive).

Par définition de $e := \sigma(d)$, la fonction $\varphi_e$ est :
\begin{itemize}
\item soit définie pour tout $n$ (et de valeur $1729$), ce qui se
  produit exatement lorsque l'exécution de $d$ ne termine jamais,
  i.e. $\varphi_d(0) \uparrow$,
\item soit définie jusqu'en un certain $n$ et non définie après, ce
  qui se produit exactement lorsque l'exécution de $d$ termine en un
  certain nombre d'étapes, i.e. $\varphi_d(0) \downarrow$.
\end{itemize}
En particulier, si $\varphi_d(0)\uparrow$, alors $\varphi_{\sigma(d)}$
est totale (et constante !), tandis que si $\varphi_d(0)\downarrow$,
alors $\varphi_{\sigma(d)}$ n'est pas totale.

Bref, on a construit $\sigma\colon\mathbb{N}\to\mathbb{N}$ totale
calculable telle que $d \in H$ si et seulement si $\sigma(d) \not\in
T$, ou, ce qui revient au même, $d \in \complement H$ si et seulement
si $\sigma(d) \in T$.  En termes de réductions, ceci signifie
$\complement H \mathrel{\leq_\mathrm{m}} T$.  Comme $\complement H$
n'est pas semi-décidable, $T$ ne l'est pas non plus.

\emph{Remarque :} Comme dans la question précédente, on n'est pas
obligé d'utiliser le terme de « réduction many-to-one » pour
argumenter que $T$ n'est pas semi-décidable : on peut simplement dire
« supposant par l'absurde que $T$ soit semi-décidable, on pourrait
semi-décider $\complement H$ de la façon suivante : donné $d$, on
calcule $\sigma(d)$, on semi-décide si $\sigma(d) \in T$ et, si c'est
le cas, on termine en renvoyant “oui” ; or ce n'est pas possible, d'où
une contradiction ».
\end{corrige}

%

\exercice\label{exercise-computable-graph}\ (${\star}{\star}{\star}$)\par\nobreak

Soit $f\colon \mathbb{N} \to \mathbb{N}$ une fonction totale : montrer
qu'il y a équivalence entre les affirmations suivantes :
\begin{enumerate}
\item la fonction $f$ est calculable,
\item le graphe $\Gamma_f := \{(i,f(i)) : i\in\mathbb{N}\} =
  \{(i,q)\in\mathbb{N}^2 : q=f(i)\}$ de $f$ est décidable,
\item le graphe $\Gamma_f$ de $f$ est semi-décidable.
\end{enumerate}

(Montrer que (3) implique (1) est le plus difficile : on pourra
commencer par s'entraîner en montrant que (2) implique (1).  Pour
montrer que (3) implique (1), on pourra chercher une façon de tester
en parallèle un nombre croissant de valeurs de $q$ de manière à
s'arrêter si l'une quelconque convient.  On peut s'inspirer de
l'exercice \ref{exercise-image-computable-partial-function}
question (2).)

\begin{corrige}
Montrons que (1) implique (2) : si on dispose d'un algorithme
$\mathscr{F}$ capable de calculer $f(i)$ en fonction de $i$, alors il
est facile d'écrire un algorithme $\mathscr{D}$ capable de décider si
$q=f(i)$ (il suffit de calculer $f(i)$ avec l'algorithme $\mathscr{F}$
supposé exister, de comparer avec la valeur de $q$ fournie, et de
renvoyer vrai/$1$ si elles sont égales, et faux/$0$ sinon),
c'est-à-dire que l'algorithme $\mathscr{D}$ décide $\Gamma_f$.

Le fait que (2) implique (3) est évident car tout ensemble décidable
est en particulier semi-décidable.

Montrons que (2) implique (1) même si ce ne sera au final pas utile :
supposons qu'on ait un algorithme $\mathscr{D}$ qui décide $\Gamma_f$
(i.e., donnés $(i,q)$, termine toujours en temps fini, en répondant
« oui » si $q=f(i)$ et « non » si $q\neq f(i)$), et on cherche à
écrire un algorithme $\mathscr{F}$ qui calcule $f(i)$.  Pour cela,
donné un $i$, il suffit de lancer l'algorithme $\mathscr{D}$
successivement sur les valeurs $(i,0)$ puis $(i,1)$ puis $(i,2)$ et
ainsi de suite (c'est-à-dire faire une boucle infinie sur $q$
parcourant les entiers naturels et lancer $\mathscr{D}$ sur chaque
couple $(i,q)$) jusqu'à trouver un $q$ pour lequel $\mathscr{D}$
réponde vrai : on termine alors en renvoyant la valeur $q$ qu'on a
trouvée, qui vérifie $q=f(i)$ par définition de $\mathscr{D}$.
L'algorithme $\mathscr{F}$ qu'on vient de décrire termine toujours car
$f$ était supposée totale, donc il existe bien un $q$ pour lequel
$\mathscr{D}$ répondra « oui ».

Reste à montrer que (3) implique (1) : supposons maintenant qu'on ait
un algorithme $\mathscr{S}$ qui « semi-décide » $\Gamma_f$ (i.e.,
donnés $(i,q)$, termine en temps fini et répond « oui » si $q=f(i)$,
et ne termine pas sinon), et on cherche à écrire un algorithme qui,
donné $i$ en entrée, calcule $f(i)$.  Notre algorithme
(appelons-le $\mathscr{F}$) fait une boucle infinie sur $p$ parcourant
les entiers naturels : chaque $p$ est d'abord décodé comme le code
$\langle n,q\rangle$ d'un couple d'entiers naturels.  On teste si
l'exécution de $\mathscr{S}$ sur l'entrée $(i,q)$ termine en $\leq n$
étapes (ce qui est bien faisable algorithmiquement) : si oui, on
renvoie la valeur $q$ ; sinon, on continue la boucle.

Cet algorithme $\mathscr{F}$ termine toujours : en effet, pour chaque
$i$ donné, il existe $q$ tel que $(i,q) \in \Gamma_f$, à savoir $q =
f(i)$ ; et alors l'algorithme $\mathscr{S}$ doit terminer sur l'entrée
$(i,q)$, c'est-à-dire que pour $n$ assez grand, il termine en $\leq n$
étapes, donc $\mathscr{F}$ terminera lorsqu'il arrivera à $p = \langle
n,q\rangle$, et il renverra bien $q$ comme annoncé.  On a donc montré
que $f$ était calculable puisqu'on a exhibé un algorithme qui la
calcule.

(Comme dans
l'exercice \ref{exercise-image-computable-partial-function}, on peut
utiliser le $T$ de la forme normale de Kleene au lieu de parler
d'« étapes » d'exécution d'une machine de Turing.  Aussi, plutôt
qu'utiliser le codage des couples $\langle n,i\rangle$, on peut
préférer faire ainsi : on parcourt les entiers naturels $p$ en une
boucle infini et pour chacun on effectue deux boucles bornées pour
$0\leq n\leq p$ et $0\leq q\leq p$ : peu importent les bornes
précises, l'important est que pour $p$ assez grand on va finir par
tester le couple $(n,q)$.)
\end{corrige}

%

\exercice\label{exercise-recognizing-p-r-functions}\ (${\star}{\star}{\star}$)\par\nobreak

Si $e \mapsto \psi^{(1)}_e$ est la numérotation standard des fonctions
primitives récursives en une variable (= d'arité $1$) et $e \mapsto
\varphi^{(1)}_e$ celle des fonctions générales récursives en une
variable, on considère les ensembles
\[
M := \{e \in \mathbb{N} : \psi^{(1)}_e\text{~définie}\}
\]
\[
N := \{e \in \mathbb{N} : \exists
e'\in\mathbb{N}.(\psi^{(1)}_{e'}\text{~définie~et~}\varphi^{(1)}_e =
\psi^{(1)}_{e'})\}
\]
Expliquer informellement ce que signifient ces deux ensembles (en
insistant sur le rapport entre eux), dire s'il y a une inclusion de
l'un dans l'autre, et dire si l'un ou l'autre est décidable.

\begin{corrige}
L'ensemble $M$ est l'ensemble des codes valables de fonction
primitives récursives, c'est-à-dire de codes légitimes dans le langage
primitif récursif ; l'ensemble $N$ qui est $\{e \in \mathbb{N} :
\varphi^{(1)}_e \text{~est~p.r.}\}$ est l'ensemble des codes de
fonctions générales récursives qui s'avèrent être primitives
récursives (même si ce n'est pas forcément manifeste sur le
programme).  Si on préfère, $M$ est l'ensemble des \emph{intentions}
primitives récursives, alors que $N$ est l'ensemble des intentions
dont l'\emph{extension} est primitive récursive ; \textit{grosso
  modo}, l'appartenance à $M$ se lit sur le code de la fonction, celle
à $N$ se lit sur les valeurs de la fonction.

Manifestement, $M \subseteq N$, car si $\psi^{(1)}_e$ est définie, on
a $\varphi^{(1)}_e = \psi^{(1)}_e$ (la définition des fonctions
générales récursives \emph{étend} celle des fonctions p.r.).
L'inclusion dans l'autre sens ne vaut pas : on peut calculer une
fonction non p.r., jeter le résultat, et renvoyer $0$, ce qui fournit
un code $e$ tel que $\varphi^{(1)}_e$ est primitive récursive (donc
$e\in N$) et pourtant $\psi^{(1)}_e$ n'est pas définie (donc $e\not\in
M$).

L'ensemble $M$ est décidable : on peut décider de façon algorithmique
si $e$ est un numéro valable de fonction primitive récursive (i.e., si
$\psi^{(1)}_e$ est définie), il s'agit pour cela simplement de
« décoder » $e$ et de vérifier qu'il suit les conventions utilisées
pour numéroter les fonctions primitives récursives (pour être très
précis, le décodage termine parce que le code d'une liste est
supérieur à tout élément de cette liste).

L'ensemble $N$ n'est pas décidable : si $F$ désigne l'ensemble des
fonctions p.r. $\mathbb{N} \to \mathbb{N}$ (c'est-à-dire l'image de
$M$ par $e \mapsto \psi^{(1)}_e$), alors $N$ est $\{e\in\mathbb{N} :
\varphi^{(1)}_e \in F\}$, et comme $F$ n'est ni vide ni l'ensemble de
toutes les fonctions générales récursives, le théorème de Rice dit
exactement que $N$ est indécidable.
\end{corrige}

%

\exercice\label{exercise-diagonalization-0-1-p-r-functions}\ (${\star}{\star}{\star}{\star}$)\par\nobreak

On considère la fonction $f\colon \mathbb{N}^2 \to \mathbb{N}$ qui à
$(e,x)$ associe $1$ si $\psi^{(1)}_e(x) = 0$, et $0$ sinon (y compris
si $\psi^{(1)}_e$ n'est pas définie) ; ici, $e \mapsto \psi^{(1)}_e$
est la numérotation standard des fonctions primitives récursives en
une variable (= d'arité $1$).

La fonction $f$ est-elle calculable ?  Est-elle primitive récursive ?
On expliquera précisément pourquoi.  (On s'inspirera de résultats vus
en cours.)  Cela changerait-il si on inversait les valeurs $0$ et $1$
dans $f$ ?

\begin{corrige}
La fonction $f$ est calculable.  En effet,
\begin{itemize}
\item on peut décider de façon algorithmique si $e$ est un numéro
  valable de fonction primitive récursive (i.e., si $\psi^{(1)}_e$ est
  définie), il s'agit pour cela simplement de « décoder » $e$ et de
  vérifier qu'il suit les conventions utilisées pour numéroter les
  fonctions primitives récursives (pour être très précis, le décodage
  termine parce que le code d'une liste est supérieur à tout élément
  de cette liste) ;
\item lorsque c'est le cas, on peut calculer $\psi^{(1)}_e(x)$ car
  quand elle est définie elle coïncide avec $\varphi^{(1)}_e(x)$
  (numérotation des fonctions générales récursives), dont on sait
  qu'il est calculable (universalité) ;
\item calculer $f$ ne pose ensuite aucune difficulté.
\end{itemize}

Montrons que $f$ n'est pas primitive récursive (on a vu en cours que
$(e,x) \mapsto \psi^{(1)}_e(x)$ ne l'est pas, mais cela ne suffit
pas : on pourrait imaginer que le fait qu'il soit égal à $0$ soit plus
facile à tester).  Pour cela, supposons par l'absurde que $f$ soit
primitive récursive.  Par le théorème de récursion de Kleene, il
existe $e$ tel que $\psi^{(1)}_e(x) = f(e,x)$.  Or la définition même
de $f$ fait que $f(e,x) \neq \psi^{(1)}_e(x)$ dans tous les cas : ceci
est une contradiction.  Donc $f$ n'est pas primitive récursive.

Cela ne change bien sûr rien d'échanger $0$ et $1$, c'est-à-dire de
remplacer $f$ par $1 - f$ (l'une est récursive, resp. p.r., ssi
l'autre l'est), mais la démonstration ne se serait pas appliquée telle
quelle.
\end{corrige}

%

\exercice\ (${\star}{\star}{\star}{\star}$)\par\nobreak

Soit
\[
Z := \{e \in \mathbb{N} : \exists n \in \mathbb{N}.\, (\psi^{(1)}_e(n) = 0)\}
\]
l'ensemble des codes $e$ des fonctions p.r. $\mathbb{N} \to
\mathbb{N}$ qui prennent (au moins une fois) la valeur $0$ (ici, $e
\mapsto \psi^{(1)}_e$ est la numérotation standard des fonctions
primitives récursives en une variable).

Montrer que $Z$ est semi-décidable.  Montrer qu'il n'est pas
décidable.

\begin{corrige}
Comme dans le début du corrigé de
l'exercice \ref{exercise-diagonalization-0-1-p-r-functions}, on
explique qu'on peut décider si $\psi^{(1)}_e(n)\downarrow$ (il s'agit
juste de vérifier si $e$ est un code valable de fonction p.r.) et, une
fois ce point vérifié, si $\psi^{(1)}_e(n) = 0$ (on peut calculer
$\psi^{(1)}_e(n) = \varphi^{(1)}_e(n)$ par universalité des fonctions
générales récursives).

Dès lors, pour semi-décider si $e \in Z$, il suffit de faire une
boucle infinie pour $n$ parcourant les entiers naturels, décider si
$\psi^{(1)}_e(n) = 0$ pour chacun, et si l'un d'eux est effectivement
nul, terminer et renvoyer « oui », sinon on continue la boucle.  Ceci
montre que $Z$ est semi-décidable.

Montrons qu'il n'est pas décidable : pour cela, on va ramener le
problème de l'arrêt à $Z$.  C'est en fait essentiellement ce que fait
le théorème de la forme normale de Kleene : en effet, considérons
$(p,x)$ dont il s'agit de décider si $\varphi^{(1)}_p(x)\downarrow$ :
d'après le théorème de la forme normale, ceci se produit si et
seulement si il existe un (entier codant un) arbre de calcul $n$
attestant que $\varphi^{(1)}_p(x)\downarrow$, ce qu'on écrit
$T(n,p,\dbllangle x\dblrangle)$, où $T$ est un prédicat p.r.,
c'est-à-dire qu'il s'écrit $t(n,p,\dbllangle x\dblrangle) = 0$ pour
une certaine fonction p.r. $t$ (qui teste si $n$ code un arbre de
calcul valable pour $\varphi^{(1)}_p(x)$ et renvoie $0$ si c'est le
cas, $1$ sinon).  On a ainsi $\varphi^{(1)}_p(x)\downarrow$ ssi
$\exists n\in\mathbb{N}.\, (t(n,p,\dbllangle x\dblrangle) = 0)$.
Maintenant, d'après le théorème s-m-n, on peut calculer de façon
p.r. en $p$ et $x$ le code $\rho(p,x)$ d'une fonction p.r. telle que
$\psi^{(1)}_{\rho(p,x)}(n) = t(n,p,\dbllangle x\dblrangle)$, et
d'après ce qui a été dit juste avant, on a $\rho(p,x) \in Z$,
c'est-à-dire $\exists n\in\mathbb{N}.\, (t(n,p,\dbllangle x\dblrangle)
= 0)$, se produit si et seulement si $\varphi^{(1)}_p(x)\downarrow$,
c'est-à-dire $(p,x) \in \mathscr{H}$ (où $\mathscr{H} := \{(p,x) \in
\mathbb{N}^2 : \varphi^{(1)}_p(x)\downarrow\}$ désigne le problème de
l'arrêt).  Ceci constitue une réduction \textit{many-to-one} de
$\mathscr{H}$ à $Z$, donc $Z$ ne peut pas être décidable : en effet,
si $Z$ était décidable, pour tester si $(p,x) \in \mathscr{H}$ il
suffirait de tester si $\rho(p,x) \in Z$, donc le problème de l'arrêt
serait décidable, ce qui n'est pas le cas.

(De nouveau, si on n'aime pas le théorème de la forme normale de
Kleene, on peut faire ça avec des étapes de machine de Turing :
appeler $t(n,p,x)$ la fonction qui renvoie $0$ si la machine de Turing
codée par $p$ termine en $\leq n$ étapes à partir de la configuration
initiale codée par $x$, et $1$ sinon : le reste du raisonnement est
essentiellement identique.)
\end{corrige}

%

%% \exercice\ (${\star}{\star}$)\par\nobreak

%% Montrer qu'il existe une machine de Turing qui, quand on la lance sur
%% la configuration vierge (c'est-à-dire un ruban vierge et dans
%% l'état $1$), termine après avoir écrit son propre programme sur sa
%% bande\footnote{Par exemple avec la convention suivante : les
%% instructions du programme $\delta \colon \{1,\ldots,m\} \times \{0,1\}
%% \to \{0,\ldots,m\} \times \{0,1\} \times \{\texttt{L},\texttt{R}\}$
%% sont écrites de la gauche vers la droite dans l'ordre $\delta(1,0)$,
%% $\delta(1,1)$, $\delta(2,0)$, $\delta(2,1)$, etc., chacune étant
%% écrite sous forme du nouvel état, du nouveau symbole, et de la
%% direction codée par $\texttt{L}\mapsto 0, \texttt{R}\mapsto 1$, tous
%% les trois en unaire séparés par des $0$.}

%

\exercice\ (${\star}{\star}{\star}$)\par\nobreak

On rappelle que le mot « configuration », dans le contexte de
l'exécution d'une machine de Turing, désigne la donnée de l'état
interne de la machine, de la position de la tête de lecture, et de la
totalité de la bande.  (Et la « configuration vierge » est la
configuration où l'état est $1$, la tête est à la position $0$, et la
bande est entièrement remplie de $0$.)

On considère l'ensemble $\mathscr{F}$ des machines de Turing $M$ dont
l'exécution, à partir de la configuration vierge $C_0$, conduit à un
nombre fini de configurations distinctes (i.e., si on appelle
$C^{(n)}$ la configuration atteinte au bout de $n$ étapes d'exécution
en démarrant sur $C_0$, on demande que l'ensemble $\{C^{(n)} : n\in
\mathbb{N}\}$ soit fini).

\textbf{(1)} Montrer que $\mathscr{F}$ est semi-décidable.
(\emph{Indication :} on pourra commencer par remarquer, en le
justifiant, que « passer par un nombre fini de configurations
distinctes » équivaut à « terminer ou revenir à une configuration déjà
atteinte ».)

\textbf{(2)} Montrer que $\mathscr{F}$ n'est pas décidable.
(\emph{Indication :} si on savait décider $\mathscr{F}$ on saurait
décider le problème de l'arrêt.)

\begin{corrige}
\textbf{(1)} Commençons par remarquer que « passer par un nombre fini
de configurations distinctes » équivaut à « terminer ou revenir à une
configuration déjà atteinte ».  En effet, dans un sens, si l'exécution
termine (i.e., termine en temps fini), il est évident qu'elle n'a
parcouru qu'un nombre fini de configurations distinctes ; mais si elle
revient à une configuration déjà atteinte, la machine boucle
indéfiniment à partir de cet état puisque l'exécution est déterministe
(la configuration contient toute l'information nécessaire à
l'exécution de la machine $M$) : si $C^{(i)} = C^{(j)}$ avec $i<j$
alors $C^{(i+k)} = C^{(j+k)}$ pour tout $k$, et donc toute
configuration atteinte est une de $C^{(0)}, \ldots, C^{(j)}$) Dans
l'autre sens, si la machine ne passe que par un nombre fini de
configurations distinctes et ne s'arrête pas, par le principe des
tiroirs, il y aura forcément une configuration atteinte plusieurs
fois, c'est-à-dire $C^{(i)} = C^{(j)}$ avec $i<j$.  Ceci montre
l'équivalence affirmée.

Pour semi-décider $\mathscr{F}$, il suffit de lancer l'exécution à
partir de $C_0$, en enregistrant chaque configuration atteinte (on
rappelle qu'une configuration est une donnée finie puisqu'il n'y a, à
un instant donné, qu'un nombre fini de $1$ sur la bande), et la
comparer à toutes les configurations précédemment atteintes.  Si on
repasse dans une configuration déjà atteinte, on termine et répond
« oui », sinon on continue l'exécution.  D'après ce qui vient d'être
dit, ceci semi-décide $\mathscr{F}$.

\textbf{(2)} Supposons par l'absurde qu'on soit capable de décider
$\mathscr{F}$ et montrons que ceci permettrait de décider le problème
de l'arrêt à partir de la configuration vierge (dont on a vu en cours
qu'il est indécidable).  En effet, donné $M$, on commence par tester
(grâce à notre hypothèse) si $M \in \mathscr{F}$ : si ce n'est pas le
cas, on sait déjà que $M$ ne terminera pas et on répond « non » ; si
c'est le cas, on sait que l'exécution de $M$ à partir de la
configuration vierge conduira soit à l'arrêt soit à retomber sur une
configuration déjà atteinte : il suffit de simuler cette exécution en
enregistrant chaque configuration atteinte, et, si on tombe sur une
configuration déjà atteinte on répond « non », tandis que si on
s'arrête on répond « oui ».  Cet algorithme termine toujours et décide
le problème de l'arrêt, ce qui est impossible : c'est donc que
$\mathscr{F}$ n'était pas décidable.
\end{corrige}

%

\exercice\ (${\star}{\star}{\star}{\star}$)\par\nobreak

Soit $f \colon\mathbb{N} \to \mathbb{N}$ une fonction calculable par
une machine de Turing en \emph{complexité d'espace} primitive
récursive : cela signifie qu'il existe $p \colon\mathbb{N} \to
\mathbb{N}$ primitive récursive et une machine de Turing $\mathscr{M}$
telle que si on lui présente $n \in \mathbb{N}$ en entrée (écrit avec
les conventions usuelles, c'est-à-dire en unaire sur la bande), la
machine s'arrête en temps fini en ayant écrit $f(n)$ sur la bande, et
de plus le nombre de cases de la bande que la tête de lecture a
parcourues est $\leq p(n)$ (c'est-à-dire que $\mathscr{M}$ a utilisé
$\leq p(n)$ cellules mémoire pour faire le calcul).

On veut montrer que $f$ elle-même est primitive récursive
(c'est-à-dire qu'une fonction calculable en complexité d'espace
p.r. est elle-même p.r., de la même manière qu'une fonction calculable
en complexité en temps p.r. est elle-même p.r.).

\textbf{(1)} Si $\mathscr{M}$ a $\leq m$ états, montrer que le calcul
de $f(n)$ par $\mathscr{M}$ ne peut faire intervenir qu'au plus
$m\times p(n)\times 2^{p(n)+n}$ configurations différentes (on rappelle
qu'une \emph{configuration} est la donnée d'un état, d'une position de
la tête, et de la valeur de chaque cellule du ruban).

\textbf{(2)} En déduire que le calcul de $f(n)$ par $\mathscr{M}$
termine en au plus $m\times p(n)\times 2^{p(n)+n}$ étapes
(\textit{indication :} sinon le calcul bouclerait indéfiniment, et on
a supposé que ce n'était pas le cas).

\textbf{(3)} En déduire que $f$ est primitive récursive.

\begin{corrige}
\textbf{(1)} Une configuration de l'exécution de $\mathscr{M}$ est la
donnée d'un état parmi au plus $m$, d'une position de la tête parmi au
plus $p(n)$ (puisque la tête visite au plus ce nombre de cellules), et
de la valeur de chaque cellule du ruban ; or au plus $p(n)+n$ cellules
peuvent contenir un $1$ (à n'importe quel moment de l'exécution), car
la tête n'en visite qu'au plus $p(n)$ et au plus $n$ portent un $1$
initialement : il y a donc au plus $2^{p(n)+n}$ configurations
possibles du ruban, et au plus $m\times p(n)\times 2^{p(n)+n}$
configurations de l'ensemble de la machine.

\textbf{(2)} Si $\mathscr{M}$ retombe sur une configuration exacte
qu'elle a déjà exécutée, elle exécutera de nouveau exactement les
mêmes instructions et ne retombera indéfiniment sur cette
configuration sans jamais finir.  Comme on a supposé que le calcul
terminait, c'est qu'il doit parcourir des configurations toutes
distinctes, donc termine en au plus $m\times p(n)\times 2^{p(n)+n}$
étapes.

\textbf{(3)} On vient de montrer que le calcul de $f$ par
$\mathscr{M}$ fait en temps au plus $m\times p(n)\times 2^{p(n)+n}$.
Mais ceci est une fonction p.r. de $n$ (car $p$ l'est, et que $k
\mapsto 2^k$ l'est, et que $m$ est une constante ici).  Donc $f$ est
calculée en complexité en temps p.r., donc elle est elle-même p.r.
\end{corrige}

%

\exercice\label{exercise-arrays-for-p-r-functions}\ (${\star}{\star}{\star}$)\par\nobreak

On s'intéresse à des tableaux d'entiers naturels, indicés par les
entiers naturels, dont toutes les valeurs valent $0$ sauf un nombre
fini (i.e., des fonctions $\tau \colon \mathbb{N} \to \mathbb{N}$ dont
le support $\{i \in \mathbb{N} : \tau(i) \neq 0\}$ est fini).  Un tel
tableau $\tau$ sera Gödel-codé comme un entier naturel sous la forme
(disons) de la liste $\dbllangle \langle i_1,\tau(i_1)\rangle, \ldots,
\langle i_k,\tau(i_k)\rangle \dblrangle$ où $i_1 < i_2 < \cdots < i_k$
sont les indices tels que la valeur $\tau(i)$ dans le tableau à cet
indice soit $\neq 0$.

\textbf{(1)} Sans rentrer dans énormément de détails, expliquer
pourquoi la fonction $(\tau,i) \mapsto \tau(i)$ (« lecture du tableau
$\tau$ à l'indice $i$ ») et $(\tau,i,v) \mapsto \tau'$ où $\tau'(j) =
\tau(j)$ sauf $\tau'(i) = v$ (« modification du tableau $\tau$ à
l'indice $i$ pour y mettre la valeur $v$ ») sont primitives
récursives.

\textbf{(2)} Sans rentrer dans énormément de détails, en déduire
pourquoi, du coup, un algorithme primitif récursif peut utiliser un
tableau dans une boucle où elle lit et écrit des valeurs arbitraires
du tableau.

\begin{corrige}
\textbf{(1)} La fonction de lecture $(\tau,i) \mapsto \tau(i)$
consiste à parcourir tous les couples de la liste $\dbllangle \langle
i_1,\tau(i_1)\rangle, \ldots, \langle i_k,\tau(i_k)\rangle \dblrangle$
qui représente le tableau et, pour chacune, comparer $i$ à la première
composante et, s'il y a égalité, renvoyer la seconde composante, sinon
renvoyer $0$.  La boucle est bornée a priori car elle parcourt une
liste connue (dont la longueur est majorée par le numéro qui la code).
Comme le décodage des couples (et donc des listes) est primitif
récursif, tout ceci est primitif récursif.

La fonction d'écriture $(\tau,i,v) \mapsto \tau'$ consiste à parcourir
la liste qui représente le tableau, et si on a $i = i_r$, modifier la
seconde composante du couple correspondant, tandis que si on a $i_r <
i < i_{r+1}$ (ou $i < i_0$ ou $i > i_k$) on insère un nouveau couple :
comme l'encodage et le décodage des couples (et notamment l'insertion
d'un élément dans une liste) sont primitifs récursifs, tout ceci est
primitif récursif.

\textbf{(2)} Le tableau est codé sous forme d'entier naturel comme on
l'a dit, donc il devient une simple variable de boucle, sur laquelle
on peut effectuer des lectures et modifications par les fonctions
qu'on a expliquées (et qui sont primitives récursives).  Le fait de
disposer d'une variable dans une boucle (bornée !) pour un algorithme
primitif récursif est bien permis (essentiellement par la récursion
primitive, qui permet précisément la modification d'une variable à
chaque tour de boucle).
\end{corrige}

%

\exercice\label{exercise-ackermann-indicator-is-p-r}\ (${\star}{\star}{\star}{\star}{\star}$)\par\nobreak

On rappelle la définition de la fonction d'Ackermann $A\colon
\mathbb{N}^3 \to \mathbb{N}$ :
\[
\begin{aligned}
A(m,n,0) &= m+n \\
A(m,0,1) &= 0 \\
A(m,0,k) &= 1\text{~si~}k\geq 2 \\
A(m,n+1,k+1) &= A(m,\,A(m,n,k+1),\,k)
\end{aligned}
\]
On a vu en cours que cette fonction est calculable mais non primitive
récursive.  On admettra sans discussion que $A(m,n,k)$ est croissante
en chaque variable dès que $m\geq 2$ et $n\geq 2$.  On pourra aussi
utiliser sans discussion les faits suivants :
\[
\begin{aligned}
A(m,n,1) &= mn \\
A(m,n,2) &= m^n \\
A(0,n,k) &= ((n+1)\%2)\text{~si~}k\geq 3 \\
A(1,n,k) &= 1\text{~si~}k\geq 2 \\
A(m,1,k) &= m\text{~si~}k\geq 1 \\
A(2,2,k) &= 4 \\
\end{aligned}
\]

On considérera aussi la \emph{fonction indicatrice du graphe} de $A$,
c'est-à-dire la fonction $B\colon \mathbb{N}^4 \to \mathbb{N}$ :
\[
\begin{array}{ll}
B(m,n,k,v) = 1 &\text{~si~}v = A(m,n,k) \\
B(m,n,k,v) = 0 &\text{~sinon} \\
\end{array}
\]

\textbf{(1)} Écrire un algorithme qui calcule $A(m,n,k)$ à partir de
$m$, $n$, $k$ et d'un \emph{majorant} $b$ de $A(m,n,k)$ selon le
principe suivant : si $m\geq 2$ et $n\geq 2$, pour chaque $\ell$
allant de $0$ à $k$ et chaque $i$ allant de $0$ à $b$ (bien noter :
c'est $b$ et pas $n$ ici), on calcule $A(m,i,\ell)$, et on la stocke
dans la case $(i,\ell)$ d'un tableau, à condition que les valeurs déjà
calculées et contenues dans le tableau permettent de la calculer (pour
les petites valeurs $m\leq 1$ ou $n\leq 1$ on utilise les formules
données ci-dessus ; sinon, on essaye d'utiliser la formule de
récurrence en consultant le tableau).  Expliquer pourquoi la valeur
$A(m,n,k)$ est bien calculée par cet algorithme.

\textbf{(2)} Expliquer pourquoi l'algorithme qu'on a écrit en (1) est
primitif récursif (on pourra prendre connaissance des conclusions de
l'exercice \ref{exercise-arrays-for-p-r-functions}).

\textbf{(3)} En déduire que la fonction $B$ est primitive récursive.

(Autrement dit, on ne peut pas calculer $A$ par un algorithme p.r.,
mais on peut \emph{vérifier} sa valeur, si elle est donnée en entrée,
par un tel algorithme.)

\begin{corrige}
\textbf{(1)} On commence par remarquer que les « petites valeurs »
$m\leq 1$ ou $n\leq 1$ de la fonction d'Ackermann, se calculent
facilement par des formules de l'énoncé.

L'algorithme calculant $A(m,n,k)$ est alors le suivant.  Si $m\leq 1$
ou $n\leq 1$ on peut facilement calculer la valeur comme on vient de
l'expliquer, donc on se place dans le cas $m \geq 2$ et $n \geq 2$.
(Notons aussi que toute valeur de la fonction d'Ackermann pour $m\geq
1$ est non nulle, ce qui nous permet d'utiliser $0$ pour représenter
« non calculé » dans le tableau.)

On initialise un tableau $\tau$ de deux indices, $i,\ell$,
initialement rempli de $0$, qui servira à stocker les valeurs de la
fonction d'Ackermann $A(m,i,\ell)$.  Pour chaque $\ell$ allant de $0$
à $k$ et chaque $i$ allant de $0$ à $b$, on calcule $A(m,i,\ell)$ de
la manière suivante : si $i\leq 1$ ou $\ell=0$ on utilise la formule
évoquée ci-dessus et stocke la valeur dans le tableau ; sinon, on
consulte le tableau en $(i-1,\ell)$ et, si cette valeur $u$ est
définie (c'est-à-dire non nulle), on consulte le tableau en
$(u,\ell-1)$ et, si cette valeur $w$ est définie, on la stocke dans le
tableau en $(i,\ell)$.

Autrement dit :
\begin{itemize}
\item si $m\leq 1$ ou $n\leq 1$, calculer facilement $A(m,n,k)$ et
  renvoyer sa valeur ; sinon :
\item initialiser un tableau $\tau$ (d'indices $i$ allant de $0$ à $b$
  et $\ell$ allant de $0$ à $k$, et initialement rempli de $0$),
\item pour $\ell$ allant de $0$ à $k$,
\begin{itemize}
\item pour $i$ allant de $0$ à $b$,
\begin{itemize}
\item si $i\leq 1$ ou $\ell=0$, calculer facilement $w := A(m,i,\ell)$
  et stocker $\tau(i,\ell) \leftarrow w$,
\item sinon, consulter $u := \tau(i-1,\ell)$,
\item si $u \neq 0$, consulter $w := \tau(u,\ell-1)$,
\item si $w \neq 0$, stocker $\tau(i,\ell) \leftarrow w$.
\end{itemize}
\end{itemize}
\item finalement : renvoyer $\tau(k,n)$ si elle est $>0$, sinon
  « échec ».
\end{itemize}

En Python, avec de petites variations :

{\tt\noindent
def ackermann\_small(m,n,k):\\
\strut\quad \# Returns A(m,n,k) value if m<=1 or n<=1 or k<=2\\
\strut\quad if k==0: return m+n\\
\strut\quad if k==1: return m*n\\
\strut\quad if k==2: return m**n\\
\strut\quad if n==0: return 1\\
\strut\quad if n==1: return m\\
\strut\quad if m==0: return (n+1)\%2\\
\strut\quad if m==1: return 1\\
\\
def ackermann\_bounded(m,n,k,b):\\
\strut\quad \# Returns A(m,n,k) (at least) if its value is <=b\\
\strut\quad if m<=1 or n<=1: return ackermann\_small(m,n,k)\\
\strut\quad tab = \{\}\\
\strut\quad for l in range(k+1):\\
\strut\quad \strut\quad for i in range(b+1):\\
\strut\quad \strut\quad \strut\quad if l==0 or i<=1:\\
\strut\quad \strut\quad \strut\quad \strut\quad w = ackermann\_small(m,i,l)\\
\strut\quad \strut\quad \strut\quad \strut\quad tab[(i,l)] = w\\
\strut\quad \strut\quad \strut\quad else:\\
\strut\quad \strut\quad \strut\quad \strut\quad if (i-1,l) in tab:\\
\strut\quad \strut\quad \strut\quad \strut\quad \strut\quad u = tab[(i-1,l)]\\
\strut\quad \strut\quad \strut\quad \strut\quad \strut\quad if (u,l-1) in tab:\\
\strut\quad \strut\quad \strut\quad \strut\quad \strut\quad \strut\quad w = tab[(u,l-1)]\\
\strut\quad \strut\quad \strut\quad \strut\quad \strut\quad \strut\quad tab[(i,l)] = w\\
\strut\quad if (n,k) in tab:\\
\strut\quad \strut\quad return tab[(n,k)]
}

L'algorithme repose sur la formule $A(m,i,\ell) =
A(m,A(m,i-1,\ell),\ell-1)$ (qui est une simple réécriture de la
troisième ligne de la définition), où on a appelé $u = A(m,i-1,\ell)$
et $w = A(m,u,\ell-1)$ : cete formule montre que si les valeurs
$(i-1,\ell)$ et $(u,\ell-1)$ sont trouvées dans le tableau, la valeur
$A(m,i,\ell)$ sera correctement calculée.

Or on montre par récurrence sur $\ell$ et $i$ que toutes les valeurs
pour lesquelles $A(m,i,\ell) \leq b$ (ou bien $i\leq 1$ ou $\ell=0$)
seront effectivement calculées par l'algorithme : en effet, si
$A(m,i,\ell) \leq b$, alors par la croissance en la deuxième variable
de la fonction d'Ackermann, $u := A(m,i-1,\ell)$ est lui-même $\leq b$
(ou alors $i=1$), donc aura été correctement calculé avant
$A(m,i,\ell)$, et $A(m,u,\ell-1)$ aura été calculé et stocké dans le
tableau puisque la boucle sur $\ell$ est extérieure à celle sur $i$ et
que la valeur $u$ est dans les bornes de la boucle sur $i$.

En particulier, si $b$ est un majorant de la valeur $A(m,n,k)$ qu'on
cherchait, alors l'algorithme renvoie $A(m,n,k)$.

\textbf{(2)} L'algorithme qu'on a décrit ci-dessus ne fait aucun appel
récursif et n'utilise que des boucles bornées (deux boucles
imbriquées).  L'utilisation d'un tableau est justifiée par
l'exercice \ref{exercise-arrays-for-p-r-functions}.  On a donc bien
défini une fonction primitive récursive.

\textbf{(3)} L'algorithme qu'on a décrit calcule de façon primitive
récursive en $m,n,k,b$ la valeur $A(m,n,k)$ si tant est que celle-ci
est $\leq b$.  En particulier, pour calculer $B(m,n,k,v)$ il suffit
d'appliquer cet algorithme à $m,n,k,v$ (c'est-à-dire avec $v$ lui-même
comme borne) et, s'il renvoie une valeur $w$, tester si $v=w$ : si
c'est le cas on renvoie « oui » (enfin, $1$), sinon, ou si
l'algorithme n'a pas réussi à caluler $A(m,n,k)$ on renvoie « non »
(enfin, $0$).  La fonction $B$ est donc primitive récursive.

(On dit parfois abusivement que la fonction d'Ackermann a un graphe
primitif récursif pour dire que la fonction indicatrice de son graphe
est primitive récursive.  On comparera à
l'exercice \ref{exercise-computable-graph} d'après lequel une
fonction dont la fonction indicatrice du graphe est calculable, i.e.,
générale récursive, est elle-même calculable, i.e., générale
récursive.)
\end{corrige}

%

\exercice\label{exercise-computably-inseparable-sets}\ (${\star}{\star}$)\par\nobreak

On dira que deux parties $L,M$ de $\mathbb{N}$ disjointes
(c'est-à-dire $L\cap M = \varnothing$) sont \textbf{calculablement
  séparables} lorsqu'il existe un $E \subseteq \mathbb{N}$ décidable
tel que $L \subseteq E$ et $M \subseteq \complement E$ (où
$\complement E$ désigne le complémentaire de $E$) ; dans le cas
contraire, on les dit \textbf{calculablement inséparables}.

\textbf{(1)} Expliquer pourquoi deux ensembles $L,M \subseteq
\mathbb{N}$ disjoints sont calculablement séparables si et seulement
s'il existe un algorithme qui, prenant en entrée un élément $x$
de $\mathbb{N}$ :
\begin{itemize}
\item termine toujours en temps fini,
\item répond « vrai » si $x\in L$ et « faux » si $x \in M$ (rien n'est
  imposé si $x\not\in L\cup M$).
\end{itemize}

\textbf{(2)} Expliquer pourquoi deux ensembles \emph{décidables}
disjoints sont toujours calculablement séparables.

On cherche maintenant à montrer qu'il existe deux ensembles $L,M
\subseteq \mathbb{N}$ \emph{semi-décidables} disjoints et
calculablement \emph{in}séparables.

Pour cela, on appelle $L := \{\langle e,x\rangle :
\varphi_e(x){\downarrow} = 1\}$ l'ensemble des codes des couples
$\langle e,x\rangle$ formés d'un programme (=algorithme) $e$ et d'une
entrée $x$, tels que l'exécution du programme $e$ sur l'entrée $x$
termine en temps fini et renvoie la valeur $1$ ; et $M := \{\langle
e,x\rangle : \varphi_e(x){\downarrow} = 2\}$ l'ensemble défini de la
même manière mais avec la valeur $2$.

\textbf{(3)} Pourquoi $L$ et $M$ sont-ils disjoints ?

\textbf{(4)} Pourquoi $L$ et $M$ sont-ils semi-décidables ?

\textbf{(5)} En imitant la démonstration du théorème de Turing sur
l'indécidabilité du problème de l'arrêt, ou bien en utilisant le
théorème de récursion de Kleene, montrer qu'il n'existe aucun
algorithme qui, prenant en entrée le code d'un couple $\langle
e,x\rangle$, termine toujours en temps fini et répond « vrai » si
$\langle e,x\rangle\in L$ et « faux » si $\langle e,x\rangle \in M$
(\emph{indication :} si un tel algorithme existait, on pourrait s'en
servir pour faire le contraire de ce qu'il prédit).

\textbf{(6)} Conclure.

\begin{corrige}
\textbf{(1)} Si $E$ est décidable tel que $L \subseteq E$ et $M
\subseteq \complement E$, alors un algorithme qui décide $E$
(c'est-à-dire, quand on lui fournit l'entrée $x$, répond « vrai » si
$x\in E$, et « faux » si $x \not\in E$) répond bien aux critères
demandés.  Réciproquement, donné un algorithme qui répond aux critères
demandés, si $E$ est l'ensemble des $x$ sur lesquels il répond
« vrai », alors $E$ est bien décidable (on peut toujours modifier
l'algorithme si nécessaire pour qu'il ne réponde que « vrai » ou
« faux »), et on a $L \subseteq E$ et $M \subseteq \complement E$.

\textbf{(2)} Si $L,M$ sont décidables disjoints, on peut poser $E =
L$, qui est décidable et vérifie à la fois $L \subseteq E$
(trivialement) et $M \subseteq \complement E$ (c'est une reformulation
du fait que $M$ est disjoint de $E=L$).

\textbf{(3)} Comme $L$ est l'ensemble des codes des couples $\langle
e,x\rangle$ tels que $\varphi_e(x) = 1$ et $M$ l'ensemble des codes
des couples $\langle e,x\rangle$ tels que $\varphi_e(x) = 2$, aucun
élément ne peut appartenir aux deux, c'est-à-dire qu'ils sont
disjoints.

\textbf{(4)} Pour semi-décider si le code d'un couple $\langle
e,x\rangle$ appartient à $L$, il suffit de lancer l'exécution du
programme $e$ sur l'entrée $x$ et, si elle termine en retournant $1$,
renvoyer « vrai », tandis que si elle termine en renvoyant n'importe
quelle autre valeur, faire une boucle infinie (bien sûr, si le
programme $e$ ne termine jamais sur l'entrée $x$, on ne termine pas
non plus).  Ceci montre que $L$ est semi-décidable.  Le même
raisonnement s'applique pour $M$.

\textbf{(5)} Supposons par l'absurde qu'il existe un algorithme $g$
comme annoncé (i.e., qui prend $\langle e,x\rangle$ en entrée, termine
toujours, et renvoie « vrai » si $\langle e,x\rangle\in L$ et « faux »
si $\langle e,x\rangle \in M$).  Définissons un nouvel algorithme qui,
donné un entier $e$, effectue les calculs suivants : (1º) interroger
l'algorithme $g$ supposé exister en lui fournissant le code du couple
$\langle e,e\rangle$ comme entrée, et ensuite (2º) si $g$ répond vrai,
renvoyer la valeur $2$, tandis que si $g$ répond n'importe quoi
d'autre, renvoyer la valeur $1$.  L'algorithme qui vient d'être décrit
aurait un certain numéro, disons, $c$, et la description de
l'algorithme fait qu'il termine toujours, que la valeur $\varphi_c(e)$
qu'il renvoie vaut toujours soit $1$ soit $2$, et qu'elle vaut $2$ si
$\langle e,e\rangle \in L$ (c'est-à-dire si $\varphi_e(e) = 1$) et $1$
si $\langle e,e\rangle \in M$ (c'est-à-dire si $\varphi_e(e) = 2$).
En particulier, en prenant $e=c$, on voit que $\varphi_c(c)$ doit
valoir $1$ ou $2$, doit valoir $2$ si $\varphi_c(c) = 1$ et $1$ si
$\varphi_c(c) = 2$, ce qui est une contradiction.

\emph{Variante :} La preuve ci-dessus a été rédigée en explicitant
l'argument diagonal.  On peut aussi, si on préfère, utiliser le
théorème de récursion de Kleene.  L'argument est alors le suivant.
Supposons par l'absurde qu'il existe un algorithme $g$ comme annoncé
(i.e., qui prend $\langle e,x\rangle$ en entrée, termine toujours, et
renvoie « vrai » si $\langle e,x\rangle\in L$ et « faux » si $\langle
e,x\rangle \in M$).  Définissons un nouvel algorithme qui, donné un
couple $(e,x)$, effectue les calculs suivants : (1º) interroger
l'algorithme $g$ supposé exister en lui fournissant le code $\langle
e,x\rangle$ comme entrée, et ensuite (2º) si $g$ répond vrai, renvoyer
la valeur $2$, tandis que si $g$ répond n'importe quoi d'autre,
renvoyer la valeur $1$.  On obtient ainsi une fonction $h$ calculable
totale $\mathbb{N}^2 \to \{1,2\}$ telle que $h(e,x) = 2$ lorsque
$\langle e,x\rangle\in L$ et $h(e,x) = 1$ lorsque $\langle
e,x\rangle\in M$.  Le théorème de récursion de Kleene assure qu'il
existe $e$ tel que $\varphi_e(x) = h(e,x)$ pour tout $x$, et
notamment, quelle que soit $x$ la valeur $\varphi_e(x)$ et définie et
vaut soit $1$ soit $2$, et elle vaut $2$ si $\langle e,x\rangle \in L$
(c'est-à-dire si $\varphi_e(x) = 1$) et $1$ si $\langle e,x\rangle \in
M$ (c'est-à-dire si $\varphi_e(x) = 2$).  Ceci est une contradiction.

\textbf{(6)} La question (5) montre (compte tenu de la question (1))
que $L$ et $M$ ne sont pas calculablement séparables, i.e.., sont
calculablement inséparables, tandis que (3) et (4) montrent que
$L$ et $M$ sont disjoints et semi-décidables.  On a donc bien montré
l'existence d'ensembles semi-décidables disjoints et calculablement
inséparables.
\end{corrige}

%

\exercice\label{exercise-good-bad-choice-lemma}\ (${\star}{\star}{\star}$)\par\nobreak

Dans cet exercice, on suppose qu'on doit faire un choix entre deux
options qu'on appelera « $X$ » et « $Y$ ».  L'un de ces choix est le
« bon » choix et l'autre est le « mauvais » choix, mais on ignore
lequel est lequel.

On dispose d'un programme $p$ ayant la propriété garantie suivante :
si on fournit en entrée à $p$ un programme $q$ (ne prendant, lui,
aucune entrée) qui termine et renvoie le bon choix, alors $p$ lui
aussi termine et renvoie le bon choix.  (En revanche, si $q$ fait
autre chose que renvoyer le bon choix, que ce soit parce qu'il ne
termine pas, parce qu'il renvoie le mauvais choix, ou qu'il renvoie
autre chose que $X$ ou $Y$, alors $p$ appelé sur $q$ peut faire
n'importe quoi, y compris ne pas terminer, renvoyer le mauvais choix,
ou renvoyer autre chose que $X$ ou $Y$.)

\textbf{(1)} Expliquer comment construire un programme $q$ tel que $p$
appelé sur l'entrée $q$ \emph{ne peut pas renvoyer le mauvais choix}
(il se peut qu'il ne termine pas, ou qu'il renvoie autre chose que
$X$ ou $Y$, mais il ne peut pas terminer et renvoyer la mauvais
choix).

(\emph{Indication :} utiliser l'astuce de Quine, c'est-à-dire le
théorème de récursion de Kleene, en s'inspirant d'une démonstration de
l'indécidabilité du problème de l'arrêt ou du théorème de Rice.)

\textbf{(2)} Expliquer pourquoi il n'y a pas moyen, avec le $p$ qu'on
nous a fourni, mais sans connaître le bon choix, de faire un programme
qui termine à coup sûr et renvoie le bon choix.

(\emph{Indication :} raisonner par symétrie en proposant un $p$ qui
marchera quel que soit le bon choix.)

\begin{corrige}
\textbf{(1)} On construit le programme $q$ suivant : il invoque le
programme $p$ sur $q$ lui-même et ensuite, si $p$ termine et renvoie
$X$ ou $Y$, il échange $X$ et $Y$ (autrement dit, si $p$ appelé sur
$q$ renvoie $X$, alors il renvoie $Y$, et si $p$ appelé sur $q$
renvoie $Y$, alors il renvoie $X$), tandis que dans tout autre cas il
fait une boucle infinie.  La construction de ce $q$, et notamment le
fait que $q$ fasse appel à lui-même dans sa définition, est justifiée
par l'astuce de Quine (formellement : si on note $\varphi_p(q)$ le
résultat du programme $p$ invoqué sur $q$, on appelle $h$ la fonction
qui à $q$ associe $X$ si $\varphi_p(q) = Y$ et $Y$ si $\varphi_p(q) =
X$ et $\uparrow$ dans tout autre cas, alors $h$ est calculable, donc
par le théorème de récursion de Kleene, il existe $q$ tel que le
résultat $\varphi_q(0)$ de $q$ soit $h(q)$).

Alors le résultat de $p$ invoqué sur $q$ ne peut pas être le mauvais
choix, car si c'était le mauvais choix, $q$ tout seul renverrait le
bon choix (par construction de $q$), donc $p$ invoqué sur $q$
renverrait le bon choix (par la garantie fournie sur $p$), ce qui est
une contradiction.

\textbf{(2)} Considérons le programme $p$ qui prend en entrée un
programme $q$, l'exécute (sans argument) et renvoie son résultat.
Cette opération est bien algorithmique d'après l'existence d'une
machine universelle.  Or le programme ainsi défini répond à la
garantie exigée de $p$ : si $q$ termine en renvoyant le bon choix,
alors $p$ appelé avec $q$ en argument termine aussi en renvoyant le
bon choix.  Comme ce programme ne dépend pas de l'identité du bon
choix (qui peut encore être $X$ ou $Y$), il ne permet pas de trancher
entre les deux : il est donc impossible de s'en servir pour faire un
programme qui termine à coup sûr et renvoie le bon choix.  (Plus
exactement, si on avait un tel moyen de faire, ce moyen devrait
s'appliquer encore quand on échange $X$ et $Y$, ce qui contredit le
fait qu'il renvoie toujours le bon choix.)
\end{corrige}

%

\exercice\ (${\star}{\star}{\star}$)\par\nobreak

\textbf{(1)} Montrer qu'il existe un programme $p$ tel que
$\varphi_p(q) = \varphi_q(q)$ pour tout $q$ et que $\varphi_p(p) = 42$
(autrement dit, quand on lui fournit en entrée un autre programme $q$,
le programme $p$ exécute $q$ sur $q$, mais par ailleurs $p$ exécuté
sur lui-même doit renvoyer $42$).  On expliquera soigneusement
l'utilisation du théorème de récursion de Kleene ici.

\textbf{(2)} On construit $p'$ de façon analogue à la question
précédente, mais en remplaçant $42$ par $1729$.  Que donne le résultat
de $p$ exécuté sur $p'$ ?  Et de $p'$ exécuté sur $p$ ?  Que dire des
fonctions partielles $\varphi_p$ et $\varphi_{p'}$ ?

\begin{corrige}
\textbf{(1)} Le programme $p$ fonctionne ainsi : il teste si le
programme $q$ qu'on lui a passé en entrée est $p$ lui-même (cette
autoréférence est permise par l'astuce de Quine) et, si oui,
renvoie $42$, sinon, il invoque la fonction universelle $(e,i) \mapsto
\varphi_e(i)$ pour calculer $\varphi_q(q)$.

De façon plus soigneuse et formelle : soit $h(p,q)$ valant $42$ si
$q=p$ et $\varphi_q(q)$ sinon ; cette fonction partielle $h \colon
\mathbb{N}^2 \dasharrow \mathbb{N}$ est manifestement récursive (étant
calculée par l'algorithme : « comparer les deux arguments, renvoyer
$42$ s'ils sont égaux, et sinon invoquer la fonction universelle sur
le second argument deux fois »).  Par le théorème de récursion de
Kleene, il existe $p$ tel que $\varphi_p(q) = h(p,q)$ : on a donc bien
$\varphi_p(p) = 42$ et $\varphi_p(q) = \varphi_q(q)$ pour tout
autre $q$ (mais pour $q=p$ cette relation est triviale donc elle vaut
encore).

\textbf{(2)} Le programme $p$ exécuté sur $p'$ renvoie $\varphi_p(p')
= \varphi_{p'}(p') = 1729$ (la première égalité découlant de la
définition de $p$ et la seconde de celle de $p'$).  Symétriquement, le
programme $p'$ exécuté sur $p$ renvoie $\varphi_{p'}(p) = \varphi_p(p)
= 42$ (la première égalité découlant de la définition de $p'$ et la
seconde de celle de $p$).

Les fonctions partielles $e \mapsto \varphi_p(e)$ et $e \mapsto
\varphi_{p'}(e)$ sont \emph{égales} puisque toutes les deux coïncident
avec $e \mapsto \varphi_e(e)$ en tout point d'après la définition de
$p$ et $p'$.

\emph{Remarque :} Ce qui rend cet exercice légèrement « paradoxal »,
c'est qu'on a construit deux programmes $p,p'$ qui calculent la
\emph{même} fonction (comme on vient de l'expliquer), et pourtant, $p$
appelé sur lui-même renvoie $42$ tandis que $p'$ appelé sur lui-même
renvoie $1729$.  Ce n'est pourtant pas si mystérieux quand on se
rappelle que l'intention d'un programme (son code) n'est pas la même
chose que l'extension (ce qu'il calcule).
\end{corrige}

%

\exercice\ (${\star}$)\par\nobreak

Soit $h\colon \mathbb{N} \dasharrow \mathbb{N}$ une fonction partielle
calculable.  Montrer qu'il existe une fonction primitive récursive $H$
qui à $e \in \mathbb{N}$ associe un $e'$ tel que $\varphi_{e'} =
h\circ \varphi_e$ (c'est-à-dire que $\varphi_{e'}(n) =
h(\varphi_e(n))$, à condition que $m := \varphi_e(n)$ soit défini et
que $h(m)$ le soit).  On rédigera très soigneusement.

\begin{corrige}
\emph{Première approche possible :} En se rappelant la manière dont on
a numéroté les fonctions générales récursives, $e' = \dbllangle 3, 1,
d, e\dblrangle$ convient, où $d$ est un code de $h$ (c'est-à-dire $h =
\varphi_{d}$), or la fonction $e \mapsto \dbllangle 3, 1, d,
e\dblrangle$ est primitive récursive.

\emph{Deuxième approche possible :} On considère l'algorithme qui,
donné $e$ et $n$, calcule d'abord $\varphi_e(n)$ (ceci est faisable
grâce à l'existence d'un interpréteur universel), puis lui applique
$h$ (ceci est faisable car $h$ est calculable), et renvoie le
résultat.  Ceci calcule la fonction $(e,n) \mapsto h(\varphi_e(n))$,
qui est donc calculable : appelons $p$ un code de celle-ci comme
fonction générale récursive.  Par le théorème s-m-n (et en notant
comme dans le cours $s_{1,1}$ la fonction de substitution d'une
variable dans un programme de deux variables) on a
$\varphi_{s_{1,1}(p,e)}(n) = \varphi_p(e,n) = h(\varphi_e(n)) =
(h\circ\varphi_e)(n)$, donc la fonction $H \colon n \mapsto
s_{1,1}(p,e)$ convient, et elle est bien primitive récursive.

\emph{Remarque :} La première approche a l'avantage de fonctionner
encore pour les fonctions primitives récursives (i.e., si $h$ est
supposée primitive récursive et qu'on remplace $\varphi_e$ par la
numérotation correspondante $\psi_e$ des fonctions primitives
récursives.  Elle a l'inconvénient de dépendre des détails du codage
des fonctions générales récursives.  La seconde approche montre que
cette dépendance est illusoire.
\end{corrige}

%

\exercice\ (${\star}$)\par\nobreak

Soit $e \in \mathbb{N}$.  Montrer qu'il existe une fonction primitive
récursive $j \colon \mathbb{N} \to \mathbb{N}$ telle que
$\varphi_{j(n)}(k) = \varphi_e(n)$ pour tous $n,k \in \mathbb{N}$.  On
rédigera très soigneusement.

\begin{corrige}
\emph{Première rédaction possible :} En se rappelant la manière dont
on a numéroté les fonctions générales récursives, $j(n) = \dbllangle
3, 1, e, \dbllangle 1, 1, n \dblrangle\dblrangle$ convient (car
$\varphi_{\dbllangle 1, 1, n \dblrangle}(k) = n$ et
$\varphi_{\dbllangle 3, 1, e, c \dblrangle}(k) =
\varphi_e(\varphi_c(k))$), or la fonction $n \mapsto \dbllangle 3, 1,
e, \dbllangle 1, 1, n \dblrangle\dblrangle$ est primitive récursive.

\emph{Deuxième rédaction possible :} On considère l'algorithme qui,
donné $n$ et $k$, ignore $k$ et calcule $\varphi_e(n)$ (ce qui est
certainement possible car $\varphi_e$ est, par définition,
calculable).  Ceci montre que la fonction $(n,k) \mapsto \varphi_e(n)$
est calculable : appelons $p$ un code de celle-ci comme fonction
générale récursive.  Par le théorème s-m-n (et en notant comme dans le
cours $s_{1,1}$ la fonction de substitution d'une variable dans un
programme de deux variables) on a $\varphi_{s_{1,1}(p,n)}(k) =
\varphi_p(n,k) = \varphi_e(n)$, donc la fonction $j \colon n \mapsto
s_{1,1}(p,n)$ convient, et elle est bien primitive récursive.
\end{corrige}


%
%
%

\section{\texorpdfstring{$\lambda$}{Lambda}-calcul non typé}


\exercice\ (${\star}$)\par\nobreak

Pour chacun des termes suivants du $\lambda$-calcul non typé, dire
s'il est en forme normale, ou en donner la forme normale s'il y en a
une.
\textbf{(a)} $(\lambda x.x)(\lambda x.x)$
\hskip 1emplus1emminus1em
\textbf{(b)} $(\lambda x.xx)(\lambda x.x)$
\hskip 1emplus1emminus1em
\textbf{(c)} $(\lambda x.xx)(\lambda x.xx)$
\hskip 1emplus1emminus1em
\textbf{(d)} $(\lambda xx.x)(\lambda xx.x)$
\hskip 1emplus1emminus1em
\textbf{(e)} $(\lambda xy.x)(\lambda xy.x)$
\hskip 1emplus1emminus1em
\textbf{(f)} $(\lambda xy.xy)y$
\hskip 1emplus1emminus1em
\textbf{(g)} $(\lambda xy.xy)(\lambda xy.xy)$

\begin{corrige}
\textbf{(a)} $(\lambda x.x)(\lambda x.x) \rightarrow_\beta \lambda x.x$
\hskip 1emplus1emminus1em
\textbf{(b)} $(\lambda x.xx)(\lambda x.x) \rightarrow_\beta (\lambda
x.x)(\lambda x.x) \rightarrow_\beta \lambda x.x$
\hskip 1emplus1emminus1em
\textbf{(c)} $(\lambda x.xx)(\lambda x.xx) \rightarrow_\beta (\lambda
x.xx)(\lambda x.xx) \rightarrow_\beta \cdots$ la seule
$\beta$-réduction possible boucle donc il n'y a pas de forme normale.
\hskip 1emplus1emminus1em
\textbf{(d)} On renomme d'abord les variables liées en se rappelant
que chaque variable est liée par le $\lambda$ le plus \emph{intérieur}
sur son nom : $(\lambda xx.x)(\lambda xx.x) = (\lambda x.\lambda
x.x)(\lambda x.\lambda x.x) \mathrel{\equiv_\alpha} (\lambda x.\lambda
y.y)(\lambda u.\lambda v.v) \rightarrow_\beta \lambda y.y
\mathrel{\equiv_\alpha} \lambda x.x$
\hskip 1emplus1emminus1em
\textbf{(e)} $(\lambda xy.x)(\lambda xy.x) = (\lambda x.\lambda
y.x)(\lambda x.\lambda y.x) \rightarrow_\beta \lambda y.\lambda
x.\lambda y.x \mathrel{\equiv_\alpha} \lambda y.\lambda x.\lambda z.x
= \lambda yxz.x$
\hskip 1emplus1emminus1em
\textbf{(f)} $(\lambda xy.xy)y = (\lambda x.\lambda y.xy)y$ ici pour
faire la $\beta$-réduction on doit d'abord renommer la variable liée
par le second $\lambda$ pour éviter qu'elle capture le $y$ libre :
$(\lambda x.\lambda y.xy)y \mathrel{\equiv_\alpha} (\lambda x.\lambda
z.xz)y \rightarrow_\beta \lambda z.yz$ (le piège serait de répondre
$\lambda y.yy$ ici !)
\hskip 1emplus1emminus1em
\textbf{(g)} $(\lambda xy.xy)(\lambda xy.xy) = (\lambda x.\lambda
y.xy)(\lambda x.\lambda y.xy) \rightarrow_\beta \lambda y.(\lambda
x.\lambda y.xy)y \mathrel{\equiv_\alpha} \lambda y.(\lambda x.\lambda
z.xz)y \rightarrow_\beta \lambda y.\lambda z.yz = \lambda yz.yz
\mathrel{\equiv_\alpha} \lambda xy.xy$
\end{corrige}

%

\exercice\ (${\star}{\star}$)\par\nobreak

\textbf{(1)} Considérons le terme $T_2 := (\lambda x.xxx)(\lambda
x.xxx)$ du $\lambda$-calcul non typé.  Étudier le graphe des
$\beta$-réductions dessus, c'est-à-dire tous les termes obtenus par
$\beta$-réduction à partir de $T_2$, et les $\beta$-réductions entre
eux.

\textbf{(2)} Que se passe-t-il pour $V := (\lambda x.x(xx))(\lambda
x.x(xx))$ ?  Sans entrer dans les détails, on donnera quelques chemins
de $\beta$-réduction, notamment celui suivi par la réduction
extérieure gauche.

\textbf{(3)} Étudier de façon analogue le comportement du terme $R
:= (\lambda x.\lambda v.xxv) (\lambda x.\lambda v.xxv)$ sous l'effet
de la $\beta$-réduction.

\begin{corrige}
\textbf{(1)} La $\beta$-réduction du seul redex de $T_2$ s'écrit
$(\lambda x.xxx)(\lambda x.xxx) \rightarrow (\lambda x.xxx)(\lambda
x.xxx)(\lambda x.xxx)$.  Appelons $T_3$ le terme en question, et plus
généralement $T_n$ le terme $(\lambda x.xxx)\cdots (\lambda x.xxx)$
avec $n-1$ applications sur $(\lambda x.xxx)$ de $(\lambda x.xxx)$
(donc $n$ fois ce sous-terme au total) ; on se rappellera bien que les
parenthèses sont vers la \emph{gauche}, c'est-à-dire que $T_4$ est
$((T_1 T_1)T_1)T_1$ par exemple.  Il y a un \emph{unique} redex dans
$T_n$ (bien qu'il y ait $n$ lambdas, un seul est appliqué), à savoir
celui des deux $T_1$ les plus à gauche (ou les plus profondément
imbriqués) : la seule $\beta$-réduction possible consiste à remplacer
ce redex $T_1 T_1$ (soit $T_2$) par son réduit $(T_1 T_1) T_1$
(soit $T_3$), ce qui donne $T_{n+1}$.  Le graphe des
$\beta$-réductions est donc $T_2 \rightarrow T_3 \rightarrow T_4
\rightarrow \cdots$ avec une unique $\beta$-réduction possible à
chaque fois.  Le terme n'est pas faiblement (ni à plus forte raison
fortement) normalisable.

\textbf{(2)} La $\beta$-réduction du seul redex donne $(\lambda
x.x(xx))(\lambda x.x(xx)) \rightarrow (\lambda x.x(xx))((\lambda
x.x(xx))(\lambda x.x(xx)))$.  Notant $U := \lambda x.x(xx)$ pour y
voir plus clair, on a donc $V := UU \rightarrow U(UU)$.  Maintenant on
a deux possibilités de redex à réduire : le redex extérieur $U(UU)$
formé par l'expression tout entière, et le redex intérieur $UU$.  Le
redex extérieur $U(UU)$ se réduit en $(UU)((UU)(UU))$ (qui a
maintenant trois redex), tandis que la réduction du redex intérieur
$UU$ dans $U(UU)$ donne $U(U(UU))$.  Il est alors facile de construire
toutes sortes de chemins de $\beta$-réductions en se rappelant que
tout $UX$ est un redex, avec pour réduit $X(XX)$.  On peut notamment
distinguer la réduction intérieure gauche (ou droite, ici elles
coïncident)
\[
UU \rightarrow U(UU) \rightarrow U(U(UU)) \rightarrow U(U(U(UU)))
\rightarrow U(U(U(U(UU)))) \cdots
\]
et la réduction extérieure gauche (on utilise $V := UU$ pour plus de
clarté ; mais on gardera bien à l'esprit que $V$, contrairement à $U$,
n'est pas une abstraction donc ne forme pas un redex quand on
l'applique, par contre c'est lui-même un redex qui se réduit en $UV$)
\[
\begin{aligned}
UU := V &\rightarrow UV \rightarrow V(VV) \rightarrow (UV)(VV)
\rightarrow (V(VV))(VV)\\
&\rightarrow ((UV)(VV))(VV) \rightarrow ((VV)(VV))(VV)
\rightarrow \cdots
\end{aligned}
\]
ou encore la réduction extérieure droite
\[
\begin{aligned}
UU := V &\rightarrow UV \rightarrow V(VV) \rightarrow V(V(UV))
\rightarrow V(V(V(VV)))\\
&\rightarrow V(V(V(V(UV)))) \rightarrow V(V(V(V(V(VV)))))
\rightarrow \cdots
\end{aligned}
\]
Aucun de ces chemins ne termine (on a vu en cours que si la réduction
extérieure gauche ne termine pas, aucun chemin de $\beta$-réductions
ne termine, mais ici c'est clair car une $\beta$-réduction ne peut de
toute façon qu'augmenter le nombre de $U$ dans l'expression).

(Seule la notion de réduction extérieure gauche a été définie en
cours ; on peut néanmoins définir sans difficulté les quatre
réductions extérieure gauche, intérieure gauche, extérieure droite et
intérieure droite : le redex extérieur gauche est celui dont le bord
gauche est le plus à gauche, le redex intérieur gauche est celui dont
le bord droit est le plus à gauche, le redex extérieur droite est
celui dont le bord droit est le plus à droite, et le redex intérieur
droite est celui dont le bord gauche est le plus à droite.  Peu
importent ces définitions, cependant, ici le but est simplement
d'illustrer quelques chemins possibles.)

(\emph{Remarque :} Je n'ai pas réfléchi à trouver une caractérisation
de tous les termes en lesquels $UU$ peut se réduire, mais on peut les
décrire comme des arbres d'application avec $U$ aux feuilles, et la
$\beta$-réduction se voir alors comme une transformation simple sur
les arbres.)

\textbf{(3)} La $\beta$-réduction du seul redex de $R$ s'écrit
$(\lambda x.\lambda v.xxv) (\lambda x.\lambda v.xxv) \rightarrow
\lambda v.(\lambda x.\lambda v.xxv)(\lambda x.\lambda v.xxv)v =
\lambda v. R v$.  On peut alors continuer ainsi : $R \rightarrow
\lambda v. Rv \rightarrow \lambda v. (\lambda v. Rv)v \rightarrow
\lambda v. (\lambda v. (\lambda v. Rv)v)v \rightarrow \cdots$.  Même
si ces écritures sont correctes (rappelons que chaque variable est
liée par le $\lambda$ le plus \emph{intérieur} sur son nom), il est
considérablement plus clair de renommer les variables liées, par
exemple ainsi :
\[
R \rightarrow
\lambda v_1. Rv_1 \rightarrow \lambda v_1. (\lambda v_2. Rv_2)v_1 \rightarrow
\lambda v_1. (\lambda v_2. (\lambda v_3. Rv_3)v_2)v_1 \rightarrow \cdots
\]
(on prendra garde à ne pas confondre le troisième terme de cette
suite, par exemple, avec $\lambda v_1. \lambda v_2. Rv_2v_1$ qui
désigne $\lambda v_1. \lambda v_2. (Rv_2v_1)$ et qui peut s'écrire
$\lambda v_1 v_2. Rv_2v_1$ : ce n'est pas du tout la même chose !).
\end{corrige}

%

\exercice\label{exercise-traduction-entiers-de-church}\ (${\star}{\star}{\star}$)\par\nobreak

On considère la traduction évidente des termes du $\lambda$-calcul
en langage Python et/ou en Scheme définie de la manière suivante :
\begin{itemize}
\item une variable se traduit en elle-même (i.e., en
  l'identificateur de ce nom),
\item une application $(P Q)$ du $\lambda$-calcul se traduit par
  $\mathtt{P}(\mathtt{Q})$ pour le Python et par $(\mathtt{P}\ 
  \mathtt{Q})$ pour le Scheme (dans les deux cas, c'est la notation
  pour l'application d'une fonction à un terme), où
  $\mathtt{P},\mathtt{Q}$ sont les traductions de $P,Q$
  respectivement,
\item une abstraction $\lambda v. E$ du $\lambda$-calcul se traduit
  par $\texttt{(lambda $\mathtt{v}$: $\mathtt{E}$)}$ en Python et
  $\texttt{(lambda ($\mathtt{v}$) $\mathtt{E}$)}$ en Scheme (dans les
  deux cas, c'est la notation pour la création d'une fonction
  anonyme), où $\mathtt{E}$ est la traduction de $E$ et $\mathtt{v}$
  l'identificateur ayant pour nom celui de la variable $v$.
\end{itemize}

\textbf{(a)} Traduire les entiers de Church $\overline{0},
\overline{1}, \overline{2}, \overline{3}$ en Python et en Scheme.

\textbf{(b)} Écrire une fonction dans chacun de ces langages prenant
en entrée (la conversion d')un entier de Church et renvoyant l'entier
natif (c'est-à-dire au sens usuel du langage) correspondant.  On
pourra pour cela utiliser la fonction successeur qui s'écrit
$\texttt{(lambda n: n+1)}$ en Python et $\texttt{(lambda (n) (+ n
  1))}$ en Scheme.

\textbf{(c)} Traduire les fonctions $\lambda mnfx.nf(mfx)$, $\lambda
mnf.n(mf)$ et $\lambda mn.nm$ qui représentent $(m,n)\mapsto m+n$,
$(m,n)\mapsto mn$ et $(m,n)\mapsto m^n$ sur les entiers de Church en
Python et en Scheme, et vérifier leur bon fonctionnement sur quelques
exemples (en utilisant la fonction écrite en (b) pour décoder le
résultat).

\textbf{(d)} Traduire le terme non-normalisable $(\lambda x.xx)
(\lambda x.xx)$ en Python et Scheme : que se passe-t-il quand on le
fait exécuter à un interpréteur de ces langages ?  Expliquer
brièvement cette différence.

\textbf{(e)} Proposer une tentative de traduction des termes du
$\lambda$-calcul en OCaml ou Haskell : reprendre les questions
précédentes en indiquant ce qui change pour ces langages.

\begin{corrige}
\textbf{(a)} En Python : $\overline{0}$ devient \texttt{lambda f:
  lambda x: x}, $\overline{1}$ devient \texttt{lambda f: lambda x:
  f(x)}, $\overline{2}$ devient \texttt{lambda f: lambda x: f(f(x))}
(chacun sur une ligne) et $\overline{3}$ devient \texttt{lambda f:
  lambda x: f(f(f(x)))} (chacun sur une ligne).  En Scheme :
$\overline{0}$ devient \texttt{(lambda (f) (lambda (x) x))},
$\overline{1}$ devient \texttt{(lambda (f) (lambda (x) (f x)))},
$\overline{2}$ devient \texttt{(lambda (f) (lambda (x) (f (f x))))} et
$\overline{3}$ devient \texttt{(lambda (f) (lambda (x) (f (f (f
  x)))))} (espacement indifférent mais les parenthèses sont
critiques).

\textbf{(b)} Pour convertir un entier de Church en entier natif, il
suffit d'itérer la fonction successeur la nombre de fois représenté
par l'entier de Church, ce que l'entier de Church permet justement de
faire, en l'appliquant au final à $0$.  En Python, cela donne :
\texttt{def fromchurch(ch): return (ch (lambda n: n+1))(0)} (ou
\texttt{fromchurch = lambda ch: (ch (lambda n: n+1))(0)} mais dans
tous les cas sur une seule ligne) ; en Scheme : \texttt{(define
  (fromchurch ch) ((ch (lambda (n) (+ n 1))) 0))} ce qui est du sucre
syntaxique pour \texttt{(define fromchurch (lambda (ch) ((ch (lambda
  (n) (+ n 1))) 0)))}.

\textbf{(c)} Voici un exemple de code vérifiant que $2+3=5$, que
$2\times 3=6$ et que $2^3=8$ sur les entiers de Church, d'abord en
Python :

\noindent\texttt{%
churchzero = (lambda f: lambda x: x)\\
churchone = (lambda f: lambda x: f(x))\\
churchtwo = (lambda f: lambda x: f(f(x)))\\
churchthree = (lambda f: lambda x: f(f(f(x))))\\
fromchurch = lambda ch: (ch (lambda n: n+1))(0)\\
churchadd = lambda m: lambda n: lambda f: lambda x: (n(f))((m(f))(x))\\
churchmul = lambda m: lambda n: lambda f: n(m(f))\\
churchpow = lambda m: lambda n: n(m)\\
\# Check 2+3 == 5:\\
fromchurch((churchadd(churchtwo))(churchthree))\\
\# Check 2*3 == 6:\\
fromchurch((churchmul(churchtwo))(churchthree))\\
\# Check 2\textasciicircum 3 == 8:\\
fromchurch((churchpow(churchtwo))(churchthree))
}

\noindent …puis en Scheme :

\noindent\texttt{%
(define churchzero (lambda (f) (lambda (x) x)))\\
(define churchone (lambda (f) (lambda (x) (f x))))\\
(define churchtwo (lambda (f) (lambda (x) (f (f x)))))\\
(define churchthree (lambda (f) (lambda (x) (f (f (f x))))))\\
(define fromchurch (lambda (ch) ((ch (lambda (n) (+ n 1))) 0)))\\
(define churchadd (lambda (m) (lambda (n) (lambda (f) (lambda (x)\\
\ \ ((n f) ((m f) x)))))))\\
(define churchmul (lambda (m) (lambda (n) (lambda (f) (n (m f))))))\\
(define churchpow (lambda (m) (lambda (n) (n m))))\\
;; Check 2+3 == 5:\\
(fromchurch ((churchadd churchtwo) churchthree))\\
;; Check 2*3 == 6:\\
(fromchurch ((churchmul churchtwo) churchthree))\\
;; Check 2\textasciicircum 3 == 8:\\
(fromchurch ((churchpow churchtwo) churchthree))
}

Dans les deux cas, les valeurs retournées sont successivement $5$, $6$
et $8$.

Noter que dans les deux langages la syntaxe est rendue lourdingue par
le fait que (conformément aux conventions du $\lambda$-calcul dont on
a mécaniquement traduit des termes) on ne crée que des fonctions
d'\emph{un} argument, ce qui oblige les fonctions d'opération à
prendre les arguments sous forme « curryfiée ».  Il serait bien plus
naturel d'écrire par exemple \texttt{churchpow = lambda m,n: n(m)} en
Python et \texttt{(define churchpow (lambda (m n) (n m)))} en Scheme
pour définir directement une fonction de \emph{deux} arguments, qu'on
peut ensuite utiliser comme \texttt{churchpow(churchtwo,churchthree)}
et \texttt{(churchpow churchtwo churchthree)} respectivement.

\texttt{(d)} En Python : $\texttt{(lambda x: x(x))(lambda x: x(x))}$ ;
en Scheme : \texttt{((lambda (x) (x x)) (lambda (x) (x x)))}.  Le
premier termine rapidement avec un débordement de pile (au moins dans
la version actuelle Python 3.11), le second, quel que soit
l'interpréteur Scheme (au moins tous ceux que j'ai pu tester), boucle
indéfiniment (mais sans consommation de pile supplémentaire ni d'autre
forme de mémoire).

La raison de cette différence est que Scheme effectue (et la
spécification du langage impose) une \emph{récursion terminale
propre} : lorsque le code d'une fonction $f$ termine par l'appel à une
autre fonction $g$ (en renvoyant sa valeur), le contrôle de
l'exécution est simplement passé de $f$ à $g$ sans empilement
d'adresse de retour (qui n'a pas lieu d'être puisque la valeur de
retour de $f$ sera justement celle de $g$) ; en Python, en revanche,
la récursion terminale n'est pas traitée spécialement, donc chaque
appel à \texttt{x(x)} est empilé et jamais dépilé et la pile déborde
rapidement.

\texttt{(e)} La traduction du $\lambda$-calcul en OCaml ou Haskell est
évidente en utilisant $\texttt{fun $\mathtt{v}$ -> $\mathtt{E}$}$
(noté $\texttt{\textbackslash $\mathtt{v}$ -> $\mathtt{E}$}$ en
Haskell) pour traduire $\lambda v.E$ (et toujours
$(\mathtt{P}\ \mathtt{Q})$ pour $(P Q)$).  Néanmoins, il n'est pas
évident qu'on puisse toujours écrire les termes qu'on souhaite, parce
qu'ils ne seront pas forcément typables.

En OCaml, le test sur les entiers de Church donne :

\noindent\texttt{%
let churchzero = fun f -> fun x -> x\\
let churchone = fun f -> fun x -> f x\\
let churchtwo = fun f -> fun x -> f (f x)\\
let churchthree = fun f -> fun x -> f (f (f x))\\
let fromchurch = fun ch -> ch (fun n->(n+1)) 0\\
let churchadd = fun m -> fun n -> fun f -> fun x -> (n f)(m f x)\\
let churchmul = fun m -> fun n -> fun f -> n (m f)\\
let churchpow = fun m -> fun n -> n m\\
;;\\
(* Check 2+3 == 5: *)\\
fromchurch(churchadd churchtwo churchthree) ;;\\
(* Check 2*3 == 6: *)\\
fromchurch(churchmul churchtwo churchthree) ;;\\
(* Check 2\textasciicircum 3 == 8: *)\\
fromchurch(churchpow churchtwo churchthree) ;;
}

\noindent …et en Haskell :

\noindent\texttt{%
let churchzero = \textbackslash f -> \textbackslash x -> x\\
let churchone = \textbackslash f -> \textbackslash x -> f x\\
let churchtwo = \textbackslash f -> \textbackslash x -> f (f x)\\
let churchthree = \textbackslash f -> \textbackslash x -> f (f (f x))\\
let fromchurch = \textbackslash ch -> ch (\textbackslash n->(n+1)) 0\\
let churchadd = \textbackslash m -> \textbackslash n -> \textbackslash f -> \textbackslash x -> (n f)(m f x)\\
let churchmul = \textbackslash m -> \textbackslash n -> \textbackslash f -> n (m f)\\
let churchpow = \textbackslash m -> \textbackslash n -> n m\\
-- Check 2+3 == 5:\\
fromchurch(churchadd churchtwo churchthree)\\
-- Check 2*3 == 6:\\
fromchurch(churchmul churchtwo churchthree)\\
-- Check 2\textasciicircum 3 == 8:\\
fromchurch(churchpow churchtwo churchthree)
}

Il se trouve que sur ces exemples simples le typage n'empêche pas la
construction, mais si on essayait de faire la fonction $n \mapsto
n^n$, par exemple, la fonction \texttt{fun n -> churchpow n n} ne type
pas.  De même, le terme non normalisant de la question (d), qui se
traduirait \texttt{(fun x -> x x)(fun x -> x x)} en OCaml, et
\texttt{(\textbackslash x -> x x)(\textbackslash x -> x x)} en
Haskell, est refusé par le système de typage : ni le OCaml ni le
Haskell ne permet de traduire tous les termes du $\lambda$-calcul non
typé, précisément parce qu'ils sont typés.

(\textbf{Attention :} si le Python comme le Scheme permettent de
\emph{traduire} tous les termes du $\lambda$-calcul non typé, le
comportement de l'évaluateur, dans les deux cas, ne correspond pas
forcément à une stratégie évidente de $\beta$-réduction du
$\lambda$-calcul.  Notamment, le terme $(\lambda uz.z)((\lambda
x.xx)(\lambda x.xx))$, bien que faiblement normalisable en
$\lambda$-calcul, conduira une fois traduit à une boucle dans ces deux
langages, parce que l'évaluateur commence par évaluer les arguments
d'une fonction avant d'appliquer la fonction ; inversement, le terme
$\lambda u.(\lambda x.xx)(\lambda x.xx)$, bien qu'il ne soit même pas
faiblement normalisable en $\lambda$-calcul, est accepté sans broncher
par ces deux langages car le corps d'une fonction n'est évalué qu'à
l'application de la fonction.  Cependant, les calculs de fonctions
primitives récursives sur les entiers de Church ne font intervenir que
des termes fortement normalisants sur lesquels ces difficultés ne se
posent pas.)
\end{corrige}

%

\exercice\label{exercise-coding-of-pairs-untyped}\ (${\star}{\star}$)\par\nobreak

On s'intéresse à une façon d'implémenter les couples en
$\lambda$-calcul non-typé : $\Pi := \lambda xyf.fxy$ (servant à faire
un couple) et $\pi_1 := \lambda p.p(\lambda xy.x)$ et $\pi_2 :=
\lambda p.p(\lambda xy.y)$ (servant à en extraire la première et la
seconde composantes).

\textbf{(1)} Montrer que, pour tous termes $X,Y$, le terme $\pi_1(\Pi
XY)$ se $\beta$-réduit en $X$ et $\pi_2(\Pi XY)$ se $\beta$-réduit en
$Y$.

\textbf{(2)} Expliquer intuitivement comment fonctionnent $\Pi$,
$\pi_1$, $\pi_2$ : comment est représentée le couple $(x,y)$ par $\Pi$
(c'est-à-dire $\Pi xy$) ?

\textbf{(3)} Écrire les fonctions $\Pi$, $\pi_1$, $\pi_2$ (on pourra
les appeler par exemple \texttt{pairing}, \texttt{proj1},
\texttt{proj2}) dans un langage de programmation fonctionnel (on
pourra prendre connaissance de l'énoncé de
l'exercice \ref{exercise-traduction-entiers-de-church}), et vérifier
leur bon fonctionnement.  (Mieux vaut, ici, choisir un langage
fonctionnel non typé, c'est-à-dire dynamiquement typé, pour mieux
refléter le $\lambda$-calcul non typé et éviter d'éventuels tracas
liés au typage.  Si le langage a des couples natifs, on pourra écrire
des conversions des couples natifs dans le codage défini ici, et vice
versa.)  Si on a des notions de compilation : sous quelle forme est
stockée l'information du couple dans la représentation faite
par $\Pi$ ?

\begin{corrige}
\textbf{(1)} Effectuons par exemple la $\beta$-réduction extérieure
gauche (mais on rappelle que le théorème de Church-Rosser affirme que
la normalisation est confluente : tout chemin de $\beta$-réduction
peut rejoindre tout autre chemin, notamment si on arrive à une forme
normale ce sera la même) : $\pi_1(\Pi XY) = (\lambda p.p(\lambda
xy.x))((\lambda xyf.fxy)XY) \rightarrow ((\lambda xyf.fxy)XY)(\lambda
xy.x) \rightarrow \rightarrow (\lambda f.fXY)(\lambda xy.x)
\rightarrow (\lambda xy.x)XY \rightarrow \rightarrow X$.  Le résultat
est le même, \textit{mutatis mutandis}, pour $\pi_2$, à savoir :
$\pi_2(\Pi XY) = (\lambda p.p(\lambda xy.y))((\lambda xyf.fxy)XY)
\rightarrow ((\lambda xyf.fxy)XY)(\lambda xy.y) \rightarrow
\rightarrow (\lambda f.fXY)(\lambda xy.y) \rightarrow (\lambda xy.y)XY
\rightarrow \rightarrow Y$.

\textbf{(2)} Le couple $(x,y)$ est codé par $\Pi$ en le terme $\Pi xy$
c'est-à-dire (à $\beta$-réduction près) la fonction $\lambda f.fxy$
qui prend une fonction $f$ et l'applique (de façon « currifiée ») aux
deux composantes du couple.  (Autrement dit, pour appliquer une
fonction au couple, on applique la représentation du couple à la
fonction !)  Pour décoder le couple, il s'agit simplement d'utiliser
pour $f$ la fonction $\lambda xy.x$ qui renvoie son premier argument
lorsqu'on veut récupérer celui-ci, et c'est ce que fait $\pi_1$, ou la
fonction $\lambda xy.x$ qui renvoie son premier argument lorsqu'on
veut récupérer celui-ci, et c'est ce que fait $\pi_2$.

\textbf{(3)} Voici une implémentation en Scheme, dans laquelle on a
pris la liberté d'utiliser des fonctions de plusieurs variables (le
Scheme permet de définir des fonctions de plusieurs variables sans
passer les arguments un par un de façon « curryfiée » : la notation
est \texttt{(f x y)} pour appeler une telle fonction \texttt{f} sur
deux arguments \texttt{x} et \texttt{y}, et \texttt{(lambda (x y)
  ...)} pour en définir une ; par ailleurs, les fonction
\texttt{cons}, \texttt{car} et \texttt{cdr} du Scheme sont les
fonctions servant nativement à créer et projeter des paires, i.e., ce
sont les équivalents natifs des fonctions \texttt{pairing},
\texttt{proj1} et \texttt{proj2} qu'on définit ici) :

\noindent\texttt{%
(define pairing (lambda (x y) (lambda (f) (f x y))))\\
(define proj1 (lambda (p) (p (lambda (x y) x))))\\
(define proj2 (lambda (p) (p (lambda (x y) y))))\\
(define fromnative (lambda (z) (pairing (car z) (cdr z))))\\
(define tonative (lambda (p) (p cons)))
}

On peut ensuite faire différents tests, par exemple \texttt{(proj1
  (pairing 42 "coucou"))} renvoie \texttt{42}, comme \texttt{(proj1
  (fromnative (cons 42 "coucou")))} ; et \texttt{(tonative (pairing 42
  "coucou"))} renvoie la paire native, notée \texttt{(42 . "coucou")}
en Scheme.

Voici maintenant le code équivalent en OCaml : il n'était pas évident
\textit{a priori} que le codage puisse être implémenté dans ce langage
(i.e., qu'il soit typable), mais il s'avère qu'il l'est, modulo la
subtilité qui sera expiquée ci-dessous :

\noindent\texttt{%
let pairing = fun x -> fun y -> fun f -> f x y\\
let proj1 = fun p -> p (fun x -> fun y -> x)\\
let proj2 = fun p -> p (fun x -> fun y -> y)\\
(* Conversion from and to native pairs *)\\
let fromnative = fun (x,y) -> pairing x y\\
let tonative = fun p -> p (fun x -> fun y -> (x,y))\\
;;
}

On peut alors tester que \texttt{proj1 (pairing 42 "coucou")} renvoie
\texttt{42}, comme \texttt{proj1 (fromnative (42,"coucou"))} ; et
\texttt{tonative (pairing 42 "coucou")} renvoie la paire native, notée
\texttt{(42, "coucou")} en OCaml.

\textbf{Subtilité :} Même si cette version OCaml fonctionne bien, une
petite variation apparemment anodine, à savoir écrire \texttt{let
  tonative = fun p -> (proj1 p, proj2 p)} (notons que l'équivalent
Scheme, \texttt{(define tonative (lambda (p) (cons (proj1 p) (proj2
  p))))}, fonctionne parfaitement) pose des problèmes de typage : avec
cette nouvelle définition, \texttt{tonative (pairing 42 1729)}
fonctionne toujours, mais \texttt{tonative (pairing 42 "coucou")}
conduit à une erreur de typage.  Sans entrer dans les détails de cette
erreur, le problème est que comme le type du résultat de la fonction
\texttt{pairing} appliquée à deux types $a$ et $b$ est $(\forall c) (a
\to b \to c) \to c$ (par exemple, celui du couple \texttt{pairing 42
  "coucou"} est $(\forall c) (\texttt{int} \to \texttt{string} \to c)
\to c$), c'est-à-dire une fonction polymorphe acceptant $f : a \to b
\to c$ pour n'importe quel type $c$ et renvoyant ce même type $c$ ; la
fonction \texttt{tonative} « devrait » donc avoir pour type $(\forall
a,b) ((\forall c) ((a \to b \to c) \to c)) \to a \times b$ (ce qui
\emph{n'est pas} la même chose que $(\forall a,b,c) ((a \to b \to c)
\to c) \to a \times b$), i.e., elle devrait recevoir une fonction
polymorphe en argument ; mais le système de typage de OCaml, basé sur
l'algorithme de Hindley-Milner, ne peut exprimer que des fonctions
polymorphes, pas des fonctions attendant une fonction polymorphe en
argument, si bien que OCaml ne peut pas typer correctement la fonction
\texttt{tonative} (et selon l'expression précise utilisée, on obtient
des approximations plus ou moins bonnes du « vrai » type qu'on vient
de dire).  Voir l'exercice \ref{exercise-system-f}(2) pour un cadre
donnant un sens précis aux types intervenant dans ce paragraphe.

Dans un quelconque de ces langages, si on implémente la fonction
\texttt{pairing} comme on vient de le dire, la valeur de $\mathtt{x}$
et $\mathtt{y}$ dans $\texttt{pairing}\ \mathtt{x}\ \mathtt{y}$ est
stockée dans la \emph{clôture} de la fonction renvoyée, c'est-à-dire
les liaisons locales (de $\mathtt{x}$ et $\mathtt{y}$ à leurs valeurs
respectives) qui ont été faites lors de la création de la fonction et
qui peuvent, ainsi que le montre cet exemple, survivre bien au-delà de
la portée de définition de la fonction dans un langage fonctionnel.
\end{corrige}


%
%
%

\section{Correspondance de Curry-Howard et calcul propositionnel intuitionniste}


\exercice\ (${\star}{\star}$)\par\nobreak

Pour chacune des preuves suivantes écrites informellement en langage
naturel dans le calcul propositionnel, écrire le $\lambda$-terme de
preuve (c'est-à-dire le terme du $\lambda$-calcul propositionnel
simplement typé étendu d'un type $\bot$ ayant pour type la proposition
prouvée) qui lui correspond.  Ces raisonnements sont-ils
intuitionnistement valables ?  Qu'en conclut-on ?

\textbf{(a)} On va prouver $\neg\neg(A\Rightarrow B) \Rightarrow
\neg\neg A\Rightarrow\neg\neg B$.  Pour cela, supposons
$\neg\neg(A\Rightarrow B)$ et $\neg\neg A$ et $\neg B$ et on veut
arriver à une contradiction.  Supposons $A\Rightarrow B$.  Alors si on
a $A$, on a $B$, ce qui contredit $\neg B$ ; donc $\neg A$ : mais ceci
contredit $\neg\neg A$.  Donc $\neg(A\Rightarrow B)$.  Mais ceci
contredit $\neg\neg(A\Rightarrow B)$, comme annoncé.

\textbf{(b)} On va prouver $(A\Rightarrow\neg\neg B) \Rightarrow
\neg\neg(A\Rightarrow B)$.  Supposons $A\Rightarrow\neg\neg B$ ainsi
que $\neg(A\Rightarrow B)$ et on veut arriver à une contradiction.  Si
on a $B$, alors certainement $A\Rightarrow B$, ce qui contredit
$\neg(A\Rightarrow B)$ : ceci montre $\neg B$.  Si on a $A$, alors
notre hypothèse $A\Rightarrow\neg\neg B$ nous donne $\neg\neg B$, d'où
une contradiction, et notamment $B$.  On a donc prouvé $A\Rightarrow
B$, d'où la contradiction recherchée.

\textbf{(c)} On va prouver $(\neg\neg A\Rightarrow\neg\neg B)
\Rightarrow (A\Rightarrow\neg\neg B)$.  Mais on sait que $A$ implique
$\neg\neg A$, donc $\neg\neg A \Rightarrow C$ implique $A\Rightarrow
C$ (quel que soit $C$, et notamment pour $\neg\neg B$).

\begin{corrige}
\textbf{(a)} On commence par se placer avec $f_1:\neg\neg(A\Rightarrow
B)$ et $x_1:\neg\neg A$ et $k:\neg B$ dans le contexte.  Appelons
encore $f_0$ l'hypothèse $A\Rightarrow B$.  Le raisonnement « si on a
$A$, on a $B$, ce qui contredit $\neg B$ » se formalise par le
$\lambda$-terme $\lambda(x_0:A). k(f_0 x_0)$ de type $\neg A$ ; le
« mais ceci contredit $\neg\neg A$ » s'obtient en lui appliquant
$x_1$.  Le $\lambda$-terme $\lambda(f_0:A\Rightarrow B).\,
x_1(\lambda(x_0:A). k(f_0 x_0))$ est donc de type $\neg(A\Rightarrow
B)$, et la contradiction avec $\neg\neg(A\Rightarrow B)$ s'exprime en
lui appliquant $f_1$.  Finalement, le $\lambda$-terme tout entier (de
type $\neg\neg(A\Rightarrow B) \Rightarrow \neg\neg
A\Rightarrow\neg\neg B$) est :
\[
\lambda(f_1:\neg\neg(A\Rightarrow B)).\,
\lambda(x_1:\neg\neg A).\,\lambda(k:\neg B).\,
f_1(\lambda(f_0:A\Rightarrow B).\,
x_1(\lambda(x_0:A). k(f_0 x_0)))
\]
(ou en syntaxe Coq : \texttt{fun (f1 : \textasciitilde
  \textasciitilde(A->B)) (x1 : \textasciitilde\textasciitilde A) (k :
  \textasciitilde B) => f1 (fun f0 : A -> B => x1 (fun x0 : A => k (f0
  x0)))}).

\textbf{(b)} On commence par se placer avec $f_1:A\Rightarrow \neg\neg
B$ et $k:\neg(A\Rightarrow B)$ dans le contexte.  Appelons $y$
l'hypothèse $B$ : alors $\lambda(z:A).y$ est de type $A\Rightarrow B$,
donc $k(\lambda(z:A).y)$ est de type $\bot$, ainsi
$\lambda(y:B).k(\lambda(z:A).y)$ est de type $\neg B$.  Appelons $x$
l'hypothèse $A$ : alors $f_1 x$ montre $\neg\neg B$, donc en
l'appliquant au $\lambda$-terme $\lambda(y:B).k(\lambda(z:A).y)$
précédemment trouvé on trouve une contradiction, et
$\texttt{exfalso}^{(B)}(f_1 x \, \lambda(y:B).k(\lambda(z:A).y))$ est
de type $B$.  En abstrayant $x$ dans ce terme on a un terme de type
$A\Rightarrow B$, et en lui appliquant $k$ on a la contradiction
recherchée.  Finalement, le $\lambda$-terme tout entier (de type
$(A\Rightarrow\neg\neg B) \Rightarrow \neg\neg(A\Rightarrow B)$) est :
\[
\lambda(f_1:A\Rightarrow \neg\neg B).\,
\lambda(k:\neg(A\Rightarrow B)).\,
k(\lambda(x:A).\,\texttt{exfalso}^{(B)}(f_1 x \, \lambda(y:B).k(\lambda(z:A).y)))
\]
(ou en syntaxe Coq : \texttt{fun (f1 : A -> \textasciitilde
  \textasciitilde B) (k : \textasciitilde (A->B)) => k (fun x : A =>
  False\_ind B (f1 x (fun y : B => k (fun z : A => y))))}).

\textbf{(c)} La preuve de $A \Rightarrow \neg\neg A$ s'écrit
$\lambda(x:A).\lambda(k:\neg A).kx$.  La preuve de $A\Rightarrow C$ à
partir de $\neg\neg A \Rightarrow C$ s'obtient en composant avec cette
fonction, donc finalement le $\lambda$-terme tout entier (de type
$(\neg\neg A\Rightarrow\neg\neg B) \Rightarrow (A\Rightarrow\neg\neg
B)$) est :
\[
\lambda(f_1:\neg\neg A\Rightarrow \neg\neg B).\,
\lambda(x:A).f_1(\lambda(k:\neg A).kx)
\]
(ou en syntaxe Coq : \texttt{fun (f1 : \textasciitilde\textasciitilde
  A -> \textasciitilde\textasciitilde B) (x : A) => f1 (fun k :
  \textasciitilde A => k x)}).

Les trois raisonnements sont parfaitement valables intuitionnistement
(malgré les nombreux « supposons (...), contradiction », il s'agit à
chaque fois de \emph{preuves de négation} et pas de \emph{preuves par
l'absurde}) : le fait qu'on ait trouvé des $\lambda$-termes (sans
aucun call/cc dedans) de ces types le montre.

On peut retenir la conclusion que $\neg\neg(A\Rightarrow B)$,
$\neg\neg A\Rightarrow\neg\neg B$ et $A\Rightarrow\neg\neg B$ sont
tous les trois équivalents (en logique intuitionniste).
\end{corrige}

%

\exercice\ (${\star}{\star}{\star}$)\par\nobreak

En utilisant une fonction call/cc (typé comme la loi de Peirce),
construire un terme de type $(A\land B \Rightarrow C) \Rightarrow
(A\Rightarrow C) \lor (B\Rightarrow C)$ dont le comportement en tant
que programme est décrit informellement comme suit : donné $f$ de type
$A\land B \Rightarrow C$, pour construire une valeur de type
$(A\Rightarrow C) \lor (B\Rightarrow C)$, on capture une continuation
(disons $k$) et on renvoie « provisoirement » une fonction
$A\Rightarrow C$ qui attend un paramètre $x$ de type $A$ et qui, quand
elle le reçoit, invoque la contination $k$ pour « revenir en arrière
dans le temps » renvoyer finalement une fonction $B\Rightarrow C$ qui
prend en entrée $y$ de type $B$ et renvoie $f\langle x,y\rangle$.

\begin{corrige}
Posons $D := (A\Rightarrow C) \lor (B\Rightarrow C)$ pour abréger les
notations.

La fonction qu'on va finalement renvoyer une fois capturé le $x$ est
$\iota_2^{(A\Rightarrow C, B\Rightarrow C)}(\lambda (y:B).\,f\langle
x,y\rangle)$ (de type $D$) : on va donc invoquer la continuation
(appelons-la $k$) sur cette valeur.  La fonction provisoirement
renvoyée est donc $\iota_1^{(A\Rightarrow C, B\Rightarrow C)}(\lambda
(x:A).\,k(\cdots))$ où les points de suspension sont remplacés par la
valeur qu'on vient de dire (ce terme est de type $D$ et on voit que la
continuation fait semblant de renvoyer un type $C$ — sachant qu'en
fait elle ne renverra jamais rien puisque c'est une continuation).

Il n'y a donc plus qu'à appliquer call/cc de type $((D\Rightarrow
C)\Rightarrow D)\Rightarrow D$ à tout ça, ce qui donne :
\[
\lambda(f:A\land B \Rightarrow C).\;
\texttt{callcc}~(
\lambda(k:D \Rightarrow C).\;\iota_1^{(A\Rightarrow C, B\Rightarrow C)}(\lambda
(x:A).\,k(\iota_2^{(A\Rightarrow C, B\Rightarrow C)}(\lambda (y:B).\,f\langle
x,y\rangle))))
\]
terme de type $(A\land B \Rightarrow C) \Rightarrow D$ (où on rappelle
$D = (A\Rightarrow C) \lor (B\Rightarrow C)$).
\end{corrige}

%

\exercice\ (${\star}$)\par\nobreak

Montrer que la formule (de Gödel-Dummett)
\[
(A\Rightarrow B) \lor (B\Rightarrow A)
\]
est prouvable en calcul propositionnel classique et \emph{n'est pas}
prouvable en calcul propositionnel intuitionniste.

\begin{corrige}
Elle est prouvable en logique classique comme on le voit en
considérant son tableau de vérité : le seul cas où $A\Rightarrow B$
est faux est celui où $A$ est vrai et $B$ est faux, et le seul cas où
$B\Rightarrow A$ est faux est celui où $B$ est vrai et $A$ est faux :
on ne peut donc pas être dans ces deux cas à la fois.  (Si on préfère,
$A\Rightarrow B$ équivaut classiquement à $\neg A \lor B$, et
$B\Rightarrow A$ équivaut classiquement à $A \lor \neg B$, donc la
disjonction des deux équivaut classiquement à $\neg A \lor B \lor A
\lor \neg B$, qui est classiquement vrai.)

Elle n'est pas prouvable en logique intuitionniste à cause de la
propriété de la disjonction : si on avait $\vdash (A\Rightarrow B)
\lor (B\Rightarrow A)$, on aurait $\vdash A\Rightarrow B$ ou bien
$\vdash B\Rightarrow A$ ; mais ni $A\Rightarrow B$ ni $B\Rightarrow A$
seul n'est prouvable intuitionnistement car elles ne sont même pas
prouvables classiquement (leur tableau de vérité n'est pas entièrement
vrai).
\end{corrige}

%

\exercice\ (${\star}$)\par\nobreak

Montrer en calcul propositionnel intuitionniste que
\[
A\lor B \Rightarrow ((A\Rightarrow B) \Rightarrow B) \land
((B\Rightarrow A) \Rightarrow A)
\]

\begin{corrige}
Voici une preuve complète écrite dans le style « drapeau » :
\bgroup\normalsize
\begin{footnotesize}
\begin{flagderiv}[exercise-pseudodisjunction-proof]
\assume{mainhyp}{A\lor B}{}
\assume{parta}{A}{}
\assume{parta-left-subhyp}{A\Rightarrow B}{}
\step{parta-left-subconc}{B}{$\Rightarrow$Élim sur \ref{parta-left-subhyp} et \ref{parta}}
\conclude{parta-left}{(A\Rightarrow B) \Rightarrow B}{$\Rightarrow$Int de \ref{parta-left-subhyp} dans \ref{parta-left-subconc}}
\assume{parta-right-subhyp}{B\Rightarrow A}{}
\conclude{parta-right}{(B\Rightarrow A) \Rightarrow A}{$\Rightarrow$Int de \ref{parta-right-subhyp} dans \ref{parta}}
\step{parta-conc}{((A\Rightarrow B) \Rightarrow B) \land
((B\Rightarrow A) \Rightarrow A)}{$\land$Int sur \ref{parta-left}, \ref{parta-right}}
\done
\assume{partb}{B}{}
\assume{partb-right-subhyp}{B\Rightarrow A}{}
\step{partb-right-subconc}{A}{$\Rightarrow$Élim sur \ref{partb-right-subhyp} et \ref{partb}}
\conclude{partb-right}{(B\Rightarrow A) \Rightarrow A}{$\Rightarrow$Int de \ref{partb-right-subhyp} dans \ref{partb-right-subconc}}
\assume{partb-left-subhyp}{A\Rightarrow B}{}
\conclude{partb-left}{(A\Rightarrow B) \Rightarrow B}{$\Rightarrow$Int de \ref{partb-left-subhyp} dans \ref{partb}}
\step{partb-conc}{(((A\Rightarrow B) \Rightarrow B) \land
(B\Rightarrow A) \Rightarrow A)}{$\land$Int sur \ref{partb-left}, \ref{partb-right}}
\done
\step{mainconc}{(((A\Rightarrow B) \Rightarrow B) \land
(B\Rightarrow A) \Rightarrow A)}{$\lor$Élim sur \ref{mainhyp} de \ref{parta} dans \ref{parta-conc} et de \ref{partb} dans \ref{partb-conc}}
\conclude{}{A\lor B \Rightarrow ((A\Rightarrow B) \Rightarrow B) \land
((B\Rightarrow A) \Rightarrow A)}{$\Rightarrow$Int de \ref{mainhyp} dans \ref{mainconc}}
\end{flagderiv}
\end{footnotesize}
\egroup

La voici écrite informellement en langage naturel :

Supposons $A\lor B$.  Considérons d'abord le cas $A$.  Si on a
$A\Rightarrow B$ alors $B$, ce qui montre $(A\Rightarrow B)\Rightarrow
B$ ; par ailleurs, de $A$ on tire aussi $(B\Rightarrow A)\Rightarrow
A$ ; on a donc montré $((A\Rightarrow B) \Rightarrow B) \land
((B\Rightarrow A) \Rightarrow A)$.  Le cas $B$ est exactement analogue
par symétrie (à l'ordre des conclusions près dans la conjonction
finale).  Dans les deux cas de la disjonction on a montré
$(A\Rightarrow B)\Rightarrow B$.  Donc finalement $A\lor B \Rightarrow
((A\Rightarrow B) \Rightarrow B) \land ((B\Rightarrow A) \Rightarrow
A)$.

En Coq, cette preuve s'écrit :

{\tt\noindent
Parameter A B C : Prop.\\
Theorem thm : A\textbackslash /B -> ((A->B)->B) /\textbackslash\ ((B->A)->A).\\
Proof.  intro H.  destruct H.  split.  intro H1.  apply H1.  exact H.
intro H2.  exact H.  split.  intro H1.  exact H.  intro H2.  apply H2.
exact H.  Qed.
\par}

On peut aussi directement donner un $\lambda$-terme correspondant à la
preuve en question :
\[
\begin{aligned}
\lambda(h:A\lor B).\,
(\texttt{match~}h\texttt{~with~}&\iota_1 h_0
\mapsto \langle \lambda(h_1:A\Rightarrow B).\,h_1 h_0,\;
\lambda(h_2:B\Rightarrow A).\,h_0\rangle,\\
&\iota_2 h_0
\mapsto \langle \lambda(h_1:A\Rightarrow B).\,h_0,\;
\lambda(h_2:B\Rightarrow A).\,h_2 h_0\rangle)
\end{aligned}
\]
(ou en syntaxe Coq : \texttt{(fun H : A \textbackslash/ B =>
 match H with
   or\_introl H0 => conj (fun H1 : A -> B => H1 H0) (fun H2 : B -> A => H0)
 | or\_intror H0 => conj (fun H1 : A -> B => H0) (fun H2 : B -> A => H2 H0)
 end)})
\end{corrige}

%

\exercice\ (${\star}{\star}$)\par\nobreak

\textbf{(1)} Montrer en calcul propositionnel intuitionniste que $(A
\lor \neg A) \Rightarrow (\neg\neg A \Rightarrow A)$.

\textbf{(2)} Montrer que $(\neg\neg A \Rightarrow A) \Rightarrow (A
\lor \neg A)$ n'est pas démontrable en calcul propositionnel
intuitionniste.  \textit{(Cette question est sans doute plus facile à
  traiter en utilisant l'une quelconque des sémantiques vues en cours
  pour le calcul propositionnel intuitionniste.)}

\textbf{(3)} Montrer qu'il revient pourtant au même, en calcul
propositionnel intuitionniste, de postuler $P \lor \neg P$ pour
\emph{toute} proposition $P$, ou bien de postuler $\neg\neg Q
\Rightarrow Q$ pour \emph{toute} proposition $Q$.  (Pour le sens qui
ne découle pas immédiatement de (1), on pourra démontrer la
proposition $\neg\neg (P \lor \neg P)$ sans hypothèse.)

\begin{corrige}
\textbf{(1)} Voici une démonstration écrite en informellement en
langage naturel : « Supposons $A \lor \neg A$, et on veut montrer
$\neg\neg A \Rightarrow A)$.  Considérons d'abord le cas $A$ : alors
certainement $\neg\neg A \Rightarrow A$.  Considérons maintenant le
cas $\neg A$ : alors $\neg\neg A$ aboutit à une contradiction, d'où on
peut tirer n'importe quelle conclusion et notamment $A$, bref,
$\neg\neg A \Rightarrow A$ dans ce cas aussi.  Dans les deux cas de la
disjonction on a montré $\neg\neg A \Rightarrow A$.  Donc finalement
$(A \lor \neg A) \Rightarrow (\neg\neg A \Rightarrow A)$. »

En Coq, cette preuve s'écrit :

{\tt\noindent
Parameter A : Prop.\\
Theorem thm : (A\textbackslash /\textasciitilde A) -> (\textasciitilde\textasciitilde A->A).\\
Proof.  intro H.  destruct H.  split.  intro H2.  exact H.
intro H2.  exfalso.  apply H2.  exact H.  Qed.
\par}

Revoici la même démonstration écrite comme un $\lambda$-terme :
\[
\begin{aligned}
\lambda(h:A\lor \neg A).\,
(\texttt{match~}h\texttt{~with~}&\iota_1 h_0
\mapsto \lambda(h_2:\neg\neg A).\,h_0,\\
&\iota_2 h_1
\mapsto \lambda(h_2:\neg\neg A).\,\texttt{exfalso}^{(A)}(h_2 h_1))
\end{aligned}
\]
(ou en syntaxe Coq : \texttt{(fun H : A \textbackslash/ \textasciitilde\ A =>
 match H with
  or\_introl H0 => fun H2 : \textasciitilde\ \textasciitilde\ A => H0
| or\_intror H1 => fun H2 : \textasciitilde\ \textasciitilde\ A => False\_ind A (H2 H1)
 end)})

\textbf{(2)} Pour prouver l'indémontrabitilité en calcul
propositionnel intuitionniste de $(\neg\neg A \Rightarrow A)
\Rightarrow (A \lor \neg A)$, on peut utiliser soit une approche
sémantique, soit une approche syntaxique.  On va expliciter ces
différentes approches.

Commençons par l'approche sémantique.  Montrons ce résultat avec
chacune des sémantiques vues en cours (n'importe laquelle suffit à
établir le résultat !).

Une preuve a été donnée en cours basée sur la sémantique des ouverts :
si $U$ désigne l'ouvert $\mathopen{]}0,1\mathclose{[}$ dans $X =
    \mathbb{R}$, alors $\dottedneg U =
    \mathopen{]}-\infty,0\mathclose{[} \cup
    \mathopen{]}1,+\infty\mathclose{[}$ donc $\dottedneg\dottedneg U =
    \mathopen{]}0,1\mathclose{[} = U$ donc $(\dottedneg\dottedneg U
    \dottedlimp U) = \mathbb{R}$ tandis que $(U \dottedlor \dottedneg
    U) = \mathbb{R}\setminus\{0,1\}$, donc finalement
    $((\dottedneg\dottedneg U \dottedlimp U) \dottedlimp (U \dottedlor
    \dottedneg U)) = \mathbb{R}\setminus\{0,1\}$ ; or si $(\neg\neg A
    \Rightarrow A) \Rightarrow (A \lor \neg A)$ était démontrable en
    calcul propositionnel intuitionniste, on trouverait $X$ tout
    entier quel que soit l'ouvert $U$ utilisé pour $A$ (par
    \emph{correction} de la sémantique des ouverts).  Comme ce n'est
    pas le cas, c'est que la proposition en question n'est pas
    démontrable.

Une autre preuve est fournie par le cadre de Kripke à trois mondes
$\{u,v,w\}$ avec $u\leq v$ et $u\leq w$ (imaginer $v,w$ comme deux
futurs possibles de $u$) et $p$ l'affectation de vérité $(u \mapsto 0,
v\mapsto 0, w\mapsto 1)$ (imaginer une vérité encore indécidée et qui
pourrait devenir fausse ou vraie), si bien que $\dottedneg p$ est
l'affectation de vérité $(u \mapsto 0, v\mapsto 1, w\mapsto 0)$, donc
$\dottedneg\dottedneg p$ est l'affectation de vérité $(u \mapsto 0,
v\mapsto 0, w\mapsto 1)$ qui est la même que $p$ et
$(\dottedneg\dottedneg p \dottedlimp p)$ est l'affectation de vérité
constante $1$ (i.e., $\dottedtop$) ; en revanche, $(p \dottedlor
\dottedneg p)$ est l'affectation de vérité $(u \mapsto 0, v\mapsto 1,
w\mapsto 1)$, et $((\dottedneg\dottedneg p \dottedlimp p) \dottedlimp
(p \dottedlor \dottedneg p))$ est également l'affectation de vérité
$(u \mapsto 0, v\mapsto 1, w\mapsto 1)$.  Or si $(\neg\neg A
\Rightarrow A) \Rightarrow (A \lor \neg A)$ était démontrable en
calcul propositionnel intuitionniste, on trouverait constamment $1$
quelle que soit l'affectation $p$ utilisé pour $A$ (par
\emph{correction} de la sémantique de Kripke).  Comme ce n'est pas
le cas, c'est que la proposition en question n'est pas démontrable.

Une autre preuve est fournie par la sémantique de la réalisabilité
propositionnelle : dans cette sémantique si $P$ est une partie
quelconque de $\mathbb{N}$, alors $\dottedneg P$ est $\mathbb{N}$ si
$P$ est vide et $\varnothing$ sinon, et $\dottedneg\dottedneg P$ est
$\varnothing$ si $P$ est vide et $\mathbb{N}$ sinon.  Ainsi,
$\dottedneg\dottedneg P \dottedlimp P$ est l'ensemble des programmes
qui prennent un entier naturel quelconque en entrée et doivent
renvoyer \emph{un élément de $P$ s'il y en a un}.  Par contraste, $P
\dottedlor \dottedneg P$ est l'ensemble des couples $\langle
0,n\rangle$ avec $n\in P$ ou bien de la forme $\langle 1,n\rangle$
avec $n$ quelconque si $P=\varnothing$.  Un élément hypothétique de
$((\dottedneg\dottedneg P \dottedlimp P) \dottedlimp (P \dottedlor
\dottedneg P))$ pour tous les $P$ à la fois serait un programme qui
prend en entrée un élément de $\dottedneg\dottedneg P \dottedlimp P$
et renvoie un élément de $P \dottedlor \dottedneg P$.  Or pour $P =
\varnothing$ (pour lequel $\dottedneg\dottedneg P \dottedlimp P$ vaut
$\mathbb{N}$), ce programme doit renvoyer un couple de la forme
$\langle 1,n\rangle$ quelle que soit son entrée ; mais pour $P =
\varnothing$, ce même programme doit renvoyer un couple de la forme
$\langle 0,n\rangle$ quand on lui passe un élément de
$\dottedneg\dottedneg P \dottedlimp P$ (qui est l'ensemble des
programmes qui terminent, et en tout cas n'est pas vide) : ceci est
contradictoire.  Or si $(\neg\neg A \Rightarrow A) \Rightarrow (A \lor
\neg A)$ était démontrable en calcul propositionnel intuitionniste, il
devrait exister un programme appartenant à $((\dottedneg\dottedneg P
\dottedlimp P) \dottedlimp (P \dottedlor \dottedneg P))$ pour tous
les $P$ (par \emph{correction} de la sémantique de la réalisabilité
propositionnelle).  Comme ce n'est pas le cas, c'est que la
proposition en question n'est pas démontrable.

Une quatrième preuve est fournie par la sémantique des problèmes finis
de Medvedev : considérons les deux problèmes
$(\{\bullet\},\{\bullet\})$ et $(\{\bullet\},\varnothing)$ (qui ont le
même ensemble de candidats, et qui sont échangés par $\dottedneg$).
Sur le premier, $\dottedneg\dottedneg P \dottedlimp P$ vaut
$(\{\bullet\},\{\bullet\})$ (où on note abusivement $\bullet$ pour
l'unique fonction $\{\bullet\} \to \{\bullet\}$) tandis que $P
\dottedlor \dottedneg P$ vaut $(\{\bullet_1,\bullet_2\},
\{\bullet_1\})$, donc $((\dottedneg\dottedneg P \dottedlimp P)
\dottedlimp (P \dottedlor \dottedneg P))$ vaut
$(\{\bullet_1,\bullet_2\}, \{\bullet_1\})$ (où on note abusivement
$\bullet_i$ pour l'unique fonction $\{\bullet\} \to
\{\bullet_1,\bullet_2\}$ envoyant $\bullet$ sur $\bullet_i$).  Sur le
second problème, $\dottedneg\dottedneg P \dottedlimp P$ vaut
$(\{\bullet\},\{\bullet\})$ tandis que $P \dottedlor \dottedneg P$
vaut $(\{\bullet_1,\bullet_2\}, \{\bullet_2\})$, donc
$((\dottedneg\dottedneg P \dottedlimp P) \dottedlimp (P \dottedlor
\dottedneg P))$ vaut $(\{\bullet_1,\bullet_2\}, \{\bullet_2\})$.  Ces
deux ensembles de solutions sont disjoints, donc il n'y a pas de
solution commune.  Or si $(\neg\neg A \Rightarrow A) \Rightarrow (A
\lor \neg A)$ était démontrable en calcul propositionnel
intuitionniste, il devrait exister une solution appartenant à
$((\dottedneg\dottedneg P \dottedlimp P) \dottedlimp (P \dottedlor
\dottedneg P))$ pour tous les $P$ ayant un même ensemble de candidats
(par \emph{correction} de la sémantique de la sémantique de Medvedev).
Comme ce n'est pas le cas, c'est que la proposition en question n'est
pas démontrable.

Outre ces quatre preuves sémantiques, on peut aussi prouver
l'idémontrabilité de $(\neg\neg A \Rightarrow A) \Rightarrow (A \lor
\neg A)$ en calcul propositionnel intuitionniste de façon syntaxique,
par la recherche d'une démonstration sans coupure, comme suit.

Puisque $\vdash (\neg\neg A \Rightarrow A) \Rightarrow (A \lor \neg
A)$ si et seulement si $\neg\neg A \Rightarrow A \vdash A \lor \neg A$
(par les règles d'introduction et d'élimination du $\Rightarrow$ en
déduction naturelle), il suffit de montrer qu'on n'a pas $\neg\neg A
\Rightarrow A \vdash A \lor \neg A$.  Par l'existence d'une
démonstration sans coupure (ou plus précisément, la propriété de la
sous-formule), si on avait ce séquent, il y en aurait une
démonstration dans laquelle toute formule qui apparaît est l'une des
six suivantes : $\bot, A, \neg A, \neg\neg A, A\lor\neg A, \neg\neg A
\Rightarrow A$.  Il s'agit donc de considérer tous les séquents ayant
un sous-ensemble de ces six formules comme hypothèses et une de ces
six formules comme conclusion : cela fait $2^6\times 6 = 384$
possibilités, un peu fastidieux à lister complètement à la main, mais
on peut se simplifier la tâcher en considérant les ensembles suivants
(tous facilement démontrables) :
\begin{itemize}
\item ceux ayant la conclusion parmi les hypothèses,
\item ceux ayant $\bot$ dans les hypothèses, ou bien à la fois $A$ et
  $\neg A$, ou bien à la fois $\neg A$ et $\neg\neg A$, et une
  conclusion quelconque,
\item $\Gamma, A \vdash \neg\neg A$ (où $\Gamma$ est quelconque),
\item $\Gamma, \neg\neg A, (A\lor\neg A) \vdash A$,
\item $\Gamma, \neg\neg A, (\neg\neg A \Rightarrow A) \vdash A$,
\item $\Gamma, A \vdash (A\lor\neg A)$,
\item $\Gamma, \neg A \vdash (A\lor\neg A)$.
\item $\Gamma, \neg\neg A, (\neg\neg A \Rightarrow A) \vdash A\lor\neg
  A$,
\item $\Gamma, A \vdash (\neg\neg A \Rightarrow A)$.
\item $\Gamma, \neg A \vdash (\neg\neg A \Rightarrow A)$.
\end{itemize}
On peut ensuite se convaincre en examinant chaque règle de la logique
(en calcul des séquents) qu'aucune application d'aucune règle à un de
ces séquents ne donne de séquent nouveau (parmi ceux dont les
hypothèses et la conclusion sont dans les six formules listées !).
Les séquents qu'on vient de lister sont donc exactement tous les
séquents valables dont les hypothèses et la conclusion sont dans les
six formules listées.  Comme $\neg\neg A \Rightarrow A \vdash A \lor
\neg A$ n'en fait pas partie, il n'est pas valable, pas plus que
$\vdash (\neg\neg A \Rightarrow A) \Rightarrow (A \lor \neg A)$.

(Cett approche syntaxique est nettement plus pénible que les approches
sémantiques qu'on a vues.  Elle a cependant l'avantage de relever d'un
algorithme systématique.)

\textbf{(3)} Observons d'abord que si on postule $P \lor \neg P$ pour
toute proposition $P$, alors on peut en déduire $\neg\neg Q
\Rightarrow Q$ pour toute proposition $Q$ : ce sens-là est évident,
car $\neg\neg Q \Rightarrow Q$ découle de $Q \lor \neg Q$ comme
expliqué dans la question (1) (et en utilisant implicitement le fait
qu'on peut substituer une proposition quelconque à une variable
propositionnelle).

Reste à traiter l'autre sens, i.e., montrer que si on postule
$\neg\neg Q \Rightarrow Q$ pour toute proposition $Q$, alors on peut
en déduire $P \lor \neg P$ pour toute proposition $P$.  Soit donc $P$
une proposition quelconque (ou une variable propositionnelle, si on
préfère).  Or $\neg(P \lor \neg P)$ équivaut à $\neg P \land \neg\neg
P$ (ceci est une application de la tautologie $((A\lor B)\Rightarrow
C) \Leftrightarrow (A\Rightarrow C) \land (B\Rightarrow C)$ en
remplaçant $A$ par $P$, $B$ par $\neg P$ et $C$ par $\bot$) ; mais
clairement $\neg P \land \neg\neg P$ implique $\bot$ (ceci est une
application de la tautologie $A \land (A\Rightarrow B) \Rightarrow B$
en remplaçant $A$ par $\neg P$ et $B$ par $\bot$) : donc on a montré
$\neg\neg(P \lor \neg P)$ (sans hypothèse).  Si on postule $\neg\neg Q
\Rightarrow Q$, il n'y a qu'à appliquer ce fait avec $Q$ valant $P
\lor \neg P$ pour conclure $P \lor \neg P$.
\end{corrige}

%

\exercice\label{exercise-negneg-conjunction}\ (${\star}{\star}$)\par\nobreak

Montrer $\neg\neg A \land \neg\neg B \Rightarrow \neg\neg (A\land B)$
en calcul propositionnel intuitionniste.  On écrira la preuve très
soigneusement et on en donnera un $\lambda$-terme.

\begin{corrige}
Voici une preuve écrite informellement en langage naturel :

Supposons $\neg\neg A\land\neg\neg B$ (donc à la fois $\neg\neg A$ et
$\neg\neg B$), et on veut montrer $\neg\neg (A\land B)$.  Pour ça,
supposons $\neg (A\land B)$ et on veut arriver à une contradiction.
Supposons provisoirement qu'on ait $A$.  Si de plus on a $B$, alors on
a $A\land B$, ce qui contredit $\neg (A\land B)$ : ceci montre $\neg
B$ (toujours sous l'hypothèse $A$ !) ; mais commme on a $\neg\neg B$,
on a une contradiction.  On a donc prouvé $\neg A$ ; mais comme on a
$\neg\neg A$, on a une contradiction.

La revoici écrite dans le style « drapeau » :
\bgroup\normalsize
\begin{footnotesize}
\begin{flagderiv}[exercise-negneg-conjunction-proof]
\assume{mainhyp}{\neg\neg A\land\neg\neg B}{}
\assume{khyp}{\neg (A\land B)}{}
\step{nna}{\neg\neg A}{$\land$Élim$_1$ sur \ref{mainhyp}}
\assume{ahyp}{A}{}
\step{nnb}{\neg\neg B}{$\land$Élim$_2$ sur \ref{mainhyp}}
\assume{bhyp}{B}{}
\step{conj}{A\land B}{$\land$Int sur \ref{ahyp}, \ref{bhyp}}
\step{contrb}{\bot}{$\Rightarrow$Élim sur \ref{khyp} et \ref{conj}}
\conclude{negb}{\neg B}{$\Rightarrow$Int de \ref{bhyp} dans \ref{contrb}}
\step{contra}{\bot}{$\Rightarrow$Élim sur \ref{nnb} et \ref{negb}}
\conclude{nega}{\neg A}{$\Rightarrow$Int de \ref{ahyp} dans \ref{contra}}
\step{contr}{\bot}{$\Rightarrow$Élim sur \ref{nna} et \ref{nega}}
\conclude{negk}{\neg\neg (A\land B)}{$\Rightarrow$Int de \ref{khyp} dans \ref{contr}}
\conclude{mainconc}{\neg\neg A\land\neg\neg B \Rightarrow \neg\neg (A\land B)}{$\Rightarrow$Int de \ref{mainhyp} dans \ref{negk}}
\end{flagderiv}
\end{footnotesize}
\egroup

La voici réécrite en forme d'arbre de preuve (en omettant certaines
hypothèses superflues à gauche du symbole ‘$\vdash$’, ou, ce qui
revient au même, en faisant un usage tacite de la règle
d'affaiblissement) :
\begin{tiny}
\[
\inferrule*[left={$\Rightarrow$Int}]{
\inferrule*[Left={$\Rightarrow$Int}]{
\inferrule*[Left={$\Rightarrow$Élim}]{
\inferrule*[Left={$\land$Élim$_1$}]{
\inferrule*[Left={Ax}]{
}{\neg\neg A \land \neg\neg B \vdash \neg\neg A \land \neg\neg B}
}{\neg\neg A \land \neg\neg B \vdash \neg\neg A}\\
\inferrule*[Right={$\Rightarrow$Int}]{
\inferrule*[Right={$\Rightarrow$Élim}]{
\inferrule*[Left={$\land$Élim$_2$}]{
\inferrule*[Left={Ax}]{
}{\neg\neg A \land \neg\neg B \vdash \neg\neg A \land \neg\neg B}
}{\neg\neg A \land \neg\neg B \vdash \neg\neg B}\\
\inferrule*[Right={$\Rightarrow$Int}]{
\inferrule*[Right={$\Rightarrow$Élim}]{
\inferrule*[Left={Ax}]{ }{\neg (A\land B) \vdash \neg (A\land B)}\\
\inferrule*[Right={$\land$Int}]{
\inferrule*[Left={Ax}]{ }{A\vdash A}\\
\inferrule*[Right={Ax}]{ }{B\vdash B}
}{A, B \vdash A\land B}
}{\neg (A\land B), A, B \vdash \bot}
}{\neg (A\land B), A \vdash \neg B}
}{\neg\neg A \land \neg\neg B, \neg (A\land B), A \vdash \bot}
}{\neg\neg A \land \neg\neg B, \neg (A\land B) \vdash \neg A}
}{\neg\neg A \land \neg\neg B, \neg (A\land B) \vdash \bot}
}{\neg\neg A \land \neg\neg B \vdash \neg\neg (A\land B)}
}{\vdash \neg\neg A \land \neg\neg B \Rightarrow \neg\neg (A\land B)}
\]
\end{tiny}

Ou, avec uniquement les conclusions de chaque séquent (mais en
indiquant, en contrepartie, à chaque décharge d'hypothèse, où cette
hypothèse est déchargée : les lettres $h,k,h_1,h_2$ sont de tels
renvois, et on a choisi les mêmes lettres que dans le terme de preuve
écrit ci-dessous pour aider à la comparaison) :
\begin{footnotesize}
\[
\inferrule*[left={$\Rightarrow$Int($h$)}]{
\inferrule*[Left={$\Rightarrow$Int($k$)}]{
\inferrule*[Left={$\Rightarrow$Élim}]{
\inferrule*[Left={$\land$Élim$_1$}]{
\inferrule*[Left={$h$}]{
}{\neg\neg A \land \neg\neg B}
}{\neg\neg A}\\
\inferrule*[Right={$\Rightarrow$Int($h_1$)}]{
\inferrule*[Right={$\Rightarrow$Élim}]{
\inferrule*[Left={$\land$Élim$_2$}]{
\inferrule*[Left={$h$}]{
}{\neg\neg A \land \neg\neg B}
}{\neg\neg B}\\
\inferrule*[Right={$\Rightarrow$Int($h_2$)}]{
\inferrule*[Right={$\Rightarrow$Élim}]{
\inferrule*[Left={$k$}]{ }{\neg (A\land B)}\\
\inferrule*[Right={$\land$Int}]{
\inferrule*[Left={$h_1$}]{ }{A}\\
\inferrule*[Right={$h_2$}]{ }{B}
}{A\land B}
}{\bot}
}{\neg B}
}{\bot}
}{\neg A}
}{\bot}
}{\neg\neg (A\land B)}
}{\neg\neg A \land \neg\neg B \Rightarrow \neg\neg (A\land B)}
\]
\end{footnotesize}

En Coq, cette preuve s'écrit :

{\tt\noindent
Parameter A : Prop.  Parameter B : Prop.\\
Theorem thm : \textasciitilde\textasciitilde A /\textbackslash\ \textasciitilde\textasciitilde B -> \textasciitilde\textasciitilde(A/\textbackslash B).\\
Proof.  intro H.  intro K.  destruct H as [H1w H2w].
apply H1w.  intro H1.  apply H2w.  intro H2.
apply K.  split.  exact H1.  exact H2.  Qed.
\par}

En voici un $\lambda$-terme de preuve :
\[
\lambda(h:\neg\neg
A\land\neg\neg B). \penalty0\, \lambda(k:\neg(A\land B)). \penalty0\,
(\pi_1 h) \penalty0 (\lambda(h_1:a). \penalty500\, (\pi_2 h)
\penalty1000 (\lambda(h_2:B). \penalty1500\, k\langle h_1,
h_2\rangle))
\]
(ou en syntaxe Coq : \texttt{fun (H :
  \textasciitilde\ \textasciitilde\ A
  /\textbackslash\ \textasciitilde\ \textasciitilde\ B) (K :
  \textasciitilde\ (A /\textbackslash\ B)) => match H with conj H1w
  H2w => H1w (fun H1 : A => H2w (fun H2 : B => K (conj H1 H2))) end} ;
la forme est un peu différente parce que Coq utilise un \texttt{match}
pour déstructurer une conjonction alors que nous avons utilisé $\pi_1$
et $\pi_2$, mais une fois noté que $\pi_1 h$ et $\pi_2 h$
correspondent à \texttt{H1w} et \texttt{H2w} c'est bien
essentiellement le même terme).
\end{corrige}

%

\exercice\ (${\star}{\star}{\star}$)\par\nobreak

Montrer qu'il revient pourtant au même, en calcul propositionnel
intuitionniste, de postuler $\neg P \lor \neg\neg P$ pour \emph{toute}
proposition $P$ (« loi faible du tiers exclu »), ou bien de postuler
$\neg (Q\land R) \Rightarrow \neg Q \lor \neg R$ pour \emph{toutes}
propositions $Q,R$ (« quatrième loi de De Morgan »).  On pourra
prendre connaissance de la conclusion de
l'exercice \ref{exercise-negneg-conjunction}.

Expliquer pourquoi $\neg (A\land B) \Rightarrow \neg A \lor \neg B$
n'est pas démontrable en calcul propositionnel intuitionniste.

\begin{corrige}
Montrons d'abord que si on postule $\neg P \lor \neg\neg P$ pour toute
proposition $P$, on peut en déduire $\neg (Q\land R) \Rightarrow \neg
Q \lor \neg R$ pour toutes propositions $Q,R$.  Soient donc $Q,R$ deux
propositions quelconques (ou variables propositionnelles, si on
préfère), et on veut prouver $\neg (Q\land R) \Rightarrow \neg Q \lor
\neg R$.  Supposons $\neg (Q\land R)$ et on veut prouver $\neg Q \lor
\neg R$.  Appliquons notre postulat $\neg P \lor \neg\neg P$ avec $P$
valant $Q$ puis $R$ successivement : on a $\neg Q \lor \neg\neg Q$ et
$\neg R \lor \neg\neg R$ ; par la loi d'élimination du $\lor$, ceci
nous permet de raisonner par cas (dans chaque cas, on cherche à
prouver $\neg Q \lor \neg R$) : dans le cas où $\neg Q$, ainsi que
dans le cas où $\neg R$, on a évidemment $\neg Q \lor \neg R$ ; il
reste donc à traiter le cas où $\neg\neg Q$ et $\neg\neg R$.  Or
$\neg\neg A \land \neg\neg B \Rightarrow \neg\neg (A\land B)$ est
démontrable en calcul propositionnel intuitionniste
(exercice \ref{exercise-negneg-conjunction}).  Donc on peut affirmer
$\neg\neg (Q\land R)$, ce qui contredit $\neg (Q\land R)$, et une
contradiction permet de tirer n'importe quelle conclusion, notamment
$\neg Q \lor \neg R$.  Bref, dans chacun des cas on a bien prouvé
$\neg Q \lor \neg R$.  On peut maintenant décharger la supposition
$\neg (Q\land R)$ et affirmer que $\neg (Q\land R) \Rightarrow \neg Q
\lor \neg R$ comme on voulait.

Réciproquement, montrons que si on postule $\neg (Q\land R)
\Rightarrow \neg Q \lor \neg R$ pour toutes propositions $Q,R$, on
peut en déduire $\neg P \lor \neg\neg P$ pour toute proposition $P$.
Soit donc $P$ une proposition quelconque (ou variable
propositionnelle, si on préfère), et on veut prouver $\neg P \lor
\neg\neg P$.  Appliquons notre postulat $\neg (Q\land R) \Rightarrow
\neg Q \lor \neg R$ avec $Q$ valant $P$ et $R$ valant $\neg P$ : il
nous donne $\neg (P\land \neg P) \Rightarrow \neg P \lor \neg\neg P$.
Or $\neg (P\land \neg P)$ est évident (une preuve en est donnée par le
$\lambda$-terme $\lambda(h:P\land\neg P). \penalty0\, (\pi_2 h)(\pi_1
h)$), donc on a bien $\neg P \lor \neg\neg P$ comme ou voulait.

Si $\neg (A\land B) \Rightarrow \neg A \lor \neg B$ était démontrable
en calcul propositionnel intuitionniste, on pourrait remplacer les
variables propositionnelles $A,B$ par des propositions $Q,R$
quelconques.  Notamment, d'après ce qu'on vient de voir au paragraphe
précédent, on pourrait démontrer $\neg P \lor \neg\neg P$ pour $P$ une
proposition quelconque.  Mais $\neg C \lor \neg\neg C$ n'est
certainement pas prouvable en calcul propositionnel intuitionniste :
car s'il l'était, par la propriété de la disjonction, $\neg C$ ou bien
$\neg\neg C$ le serait, ce qui n'est pas le cas (même classiquement,
on ne peut pas prouver $\neg C$ seul, ni $\neg\neg C$ seul, qui peut
prendre l'une ou l'autre valeur de vérité).
\end{corrige}

%

\exercice\ (${\star}{\star}$)\par\nobreak

En appliquant l'algorithme de Hindley-Milner, trouver une annotation
de type (dans le $\lambda$-calcul simplement typé) au terme suivant du
$\lambda$-calcul non typé :
\[
\lambda k. \, k\,(k\,(\lambda x.x)\,(\lambda y.y))
\]
(autrement dit, \texttt{fun k -> k (k (fun x -> x) (fun y -> y))} en
syntaxe OCaml).

\begin{corrige}
On prend bien garde au fait qu'il faut parenthéser comme $\lambda
k. \, k\,((k\,(\lambda x.x))\,(\lambda y.y))$.

Dans une première phase, on collecte les types et contraintes
suivants : $k:\eta_1$, $x:\eta_2$ donc $\lambda x.x:\eta_2\to\eta_2$,
puis $k(\lambda x.x):\eta_3$ avec $\eta_1 =
((\eta_2\to\eta_2)\to\eta_3)$, puis $y:\eta_4$ donc $\lambda
y.y:\eta_4\to\eta_4$, puis $k(\lambda x.x)(\lambda y.y):\eta_5$ avec
$\eta_3 = ((\eta_4\to\eta_4)\to\eta_5)$, puis $k(k(\lambda
x.x)(\lambda y.y)) : \eta_6$ avec $\eta_1 = (\eta_5\to\eta_6)$ ; et le
type final est $\eta_1 \to \eta_6$.  On a donc les contraintes
suivantes :
\[
\begin{aligned}
\eta_1 &= ((\eta_2\to\eta_2)\to\eta_3)\\
\eta_3 &= ((\eta_4\to\eta_4)\to\eta_5)\\
\eta_1 &= (\eta_5\to\eta_6)\\
\end{aligned}
\]
On examine d'abord la contrainte $\eta_1 =
((\eta_2\to\eta_2)\to\eta_3)$ : ici, $\eta_1$ est une variable de
type, et elle n'apparaît pas dans l'autre membre de la contrainte ; il
n'y a rien à faire à part enregistrer $\eta_1 \mapsto
((\eta_2\to\eta_2)\to\eta_3)$ dans la substitution et l'appliquer aux
contraintes qui deviennent donc :
\[
\begin{aligned}
\eta_3 &= ((\eta_4\to\eta_4)\to\eta_5)\\
((\eta_2\to\eta_2)\to\eta_3) &= (\eta_5\to\eta_6)\\
\end{aligned}
\]
On examine ensuite la contrainte $\eta_3 =
((\eta_4\to\eta_4)\to\eta_5)$ : de nouveau, $\eta_3$ est une variable
de type, et elle n'apparaît pas dans l'autre membre de la contrainte :
on enregistre $\eta_3 \mapsto ((\eta_4\to\eta_4)\to\eta_5)$ et on
l'applique à la substitution déjà enregistrée qui devient donc $\eta_1
\mapsto ((\eta_2\to\eta_2)\to(\eta_4\to\eta_4)\to\eta_5)$, ainsi qu'à
la contrainte restante.  Cette dernière contrainte à examiner est :
\[
((\eta_2\to\eta_2)\to(\eta_4\to\eta_4)\to\eta_5) = (\eta_5\to\eta_6)
\]
Cette fois-ci les deux membres sont des types complexes, donc elle se
transforme en deux contraintes (attention au parenthésage !) :
\[
\begin{aligned}
(\eta_2\to\eta_2) &= \eta_5\\
((\eta_4\to\eta_4)\to\eta_5) &= \eta_6\\
\end{aligned}
\]
On examine à présent $(\eta_2\to\eta_2) = \eta_5$ : ici, $\eta_5$ est
une variable de type, et elle n'apparaît pas dans l'autre membre de la
contrainte : on enregistre $\eta_5 \mapsto (\eta_2\to\eta_2)$, et on
applique cette substitution à ce qui reste de contrainte (qui devient
donc $((\eta_4\to\eta_4)\to\eta_2\to\eta_2) = \eta_6$) et à ce qu'on a
déjà déterminé comme substitution (qui devient donc $\eta_1 \mapsto
((\eta_2\to\eta_2)\to(\eta_4\to\eta_4)\to\eta_2\to\eta_2)$ et $\eta_3
\mapsto ((\eta_4\to\eta_4)\to\eta_2\to\eta_2)$).  Il reste enfin la
contrainte $((\eta_4\to\eta_4)\to\eta_2\to\eta_2) = \eta_6$ : de
nouveau, $\eta_6$ est une variable de type, et elle n'apparaît pas
dans l'autre membre de la contrainte : on enregistre donc $\eta_6
\mapsto ((\eta_4\to\eta_4)\to\eta_2\to\eta_2)$, et on a trouvé la
substitution finale :
\[
\begin{aligned}
\eta_1 &\mapsto ((\eta_2\to\eta_2)\to(\eta_4\to\eta_4)\to\eta_2\to\eta_2)\\
\eta_3 &\mapsto ((\eta_4\to\eta_4)\to\eta_2\to\eta_2)\\
\eta_5 &\mapsto (\eta_2\to\eta_2)\\
\eta_6 &\mapsto ((\eta_4\to\eta_4)\to\eta_2\to\eta_2)\\
\end{aligned}
\]
Il ne reste plus qu'à l'appliquer à tous les types intermédiaires et
au type final, ce qui donne
\[
\begin{aligned}
&\lambda (k:(\eta_2\to\eta_2)\to(\eta_4\to\eta_4)\to\eta_2\to\eta_2). \, k\,(k\,(\lambda (x:\eta_2).x)\,(\lambda (y:\eta_4).y))\\
:\quad& ((\eta_2\to\eta_2)\to(\eta_4\to\eta_4)\to\eta_2\to\eta_2)
\to (\eta_4\to\eta_4)\to\eta_2\to\eta_2
\end{aligned}
\]
ou, si on veut renommer les variables restantes pour être un peu plus
lisible
\[
\begin{aligned}
&\lambda (k:(\alpha\to\alpha)\to(\beta\to\beta)\to\alpha\to\alpha). \, k\,(k\,(\lambda (x:\alpha).x)\,(\lambda (y:\beta).y))\\
:\quad& ((\alpha\to\alpha)\to(\beta\to\beta)\to\alpha\to\alpha)
\to (\beta\to\beta)\to\alpha\to\alpha
\end{aligned}
\]
comme $\lambda$-terme annoté et comme type final.
\end{corrige}


%
%
%

\section{Sémantique(s) du calcul propositionnel intuitionniste}


\exercice\ (${\star}$)\par\nobreak

On considère le cadre de Kripke dessiné ci-dessous, où une flèche $w
\to w'$ signifie que $w \leq w'$ (« le monde $w'$ est accessible
depuis le monde $w$ »), sachant que la relation $\leq$ doit bien sûr
être réflexive et transitive (c'est la clôture réflexive-transitive de
celle qui est représentée par les flèches : c'est-à-dire qu'on a bien
sûr $u_0\leq u_0$ et $u_2\leq u_0$ par exemple, malgré l'absence de
flèches explicites pour le rappeler).
\begin{center}
\begin{tikzpicture}[>=stealth']
\node (u0) at (0bp,0bp) {$u_0$};
\node (v0) at (35bp,-15bp) {$v_0$};
\node (u1) at (0bp,-30bp) {$u_1$};
\node (v1) at (35bp,-45bp) {$v_1$};
\node (u2) at (0bp,-60bp) {$u_2$};
\draw[->] (u1)--(u0);
\draw[->] (u2)--(u1);
\draw[->] (v1)--(v0);
\draw[->] (u2)--(v0);
\draw[->] (v1)--(u0);
\end{tikzpicture}
\end{center}
Soit $p$ l'affectation de vérité qui vaut $1$ en $u_0$ et $0$ en
chacun de $v_0,u_1,v_1,u_2$.  Pour ce $p$, calculer l'affectation de
vérité de $((\neg\neg p \Rightarrow p)\Rightarrow(p\lor\neg p))
\Rightarrow(\neg p\lor\neg\neg p)$ (c'est-à-dire plus exactement
$((\dottedneg\dottedneg p \dottedlimp
p)\dottedlimp(p\dottedlor\dottedneg p)) \dottedlimp(\dottedneg
p\dottedlor\dottedneg\dottedneg p)$, où les points rappellent qu'on
parle de l'interprétation des connecteurs données par la sémantique de
Kripke).  En déduire que la formule $((\neg\neg A \Rightarrow
A)\Rightarrow(A\lor\neg A)) \Rightarrow(\neg A\lor\neg\neg A)$ n'est
pas démontrable en calcul propositionnel intuitionniste.

\begin{corrige}
On calcule successivement l'affectation de vérité de chacune des
sous-formules de la formule proposée : on se rappelle que
$q\Rightarrow r$ vaut $1$ en un monde $w$ lorsque dans tout monde $w'$
accessible depuis $w$ pour lequel $q(w') = 1$ on a aussi $r(w') = 1$,
et notamment, $\neg q$ vaut $q$ en un monde $w$ lorsque dans tout
monde $w'$ accessible depuis $w$ on a $q(w') = 0$.  Par ailleurs, pour
éviter de se trompre, on vérifie à tout stade que les affectations de
vérité sont permanentes, c'est-à-dire que si $q(w) = 1$ on a $q(w') =
1$ pour tout monde $w'$ accessible depuis $w$.  On obtient les
résultats tabulés ci-dessous :

\begin{center}
\begin{tabular}{r|ccccc}
Formule&$u_2$&$v_1$&$u_1$&$v_0$&$u_0$\\\hline
$p$&$0$&$0$&$0$&$0$&$1$\\
$\neg p$&$0$&$0$&$0$&$1$&$0$\\
$p\lor \neg p$&$0$&$0$&$0$&$1$&$1$\\
$\neg\neg p$&$0$&$0$&$1$&$0$&$1$\\
$\neg p\lor \neg\neg p$&$0$&$0$&$1$&$1$&$1$\\
$\neg\neg p \Rightarrow p$&$0$&$1$&$0$&$1$&$1$\\
$(\neg\neg p \Rightarrow p)\Rightarrow (p\lor \neg p)$&$1$&$0$&$1$&$1$&$1$\\
$((\neg\neg p \Rightarrow p)\Rightarrow (p\lor \neg p))
\Rightarrow(\neg p\lor\neg\neg p)$&$0$&$1$&$1$&$1$&$1$\\
\end{tabular}
\end{center}

On constate que l'affecatation de vérité de $((\neg\neg p \Rightarrow
p)\Rightarrow (p\lor \neg p)) \Rightarrow(\neg p\lor\neg\neg p)$ n'est
pas identiquement $1$ sur le cadre.  Or si $((\neg\neg A \Rightarrow
A)\Rightarrow(A\lor\neg A)) \Rightarrow(\neg A\lor\neg\neg A)$ était
démontrable en calcul propositionnel intuitionniste, on trouverait
constamment $1$ dans tout cadre et en remplaçant $p$ par n'importe
quelle affectation de vérité sur le cadre (par \emph{correction} de la
sémantique de la sémantique de Kripke) : c'est donc que cette formule
n'est pas démontrable.
\end{corrige}

%

\exercice\ (${\star}{\star}$)\par\nobreak

Pour $n\in\mathbb{N}$, on considère le cadre de Kripke
$\{w_0,\ldots,w_{n-1}\}$ formé de $n$ mondes totalement ordonnés par
$w_i \leq w_j$ lorsque $i\geq j$ (le fait d'inverser l'ordre s'avérera
plus commode pour l'écriture de la suite) :
\begin{center}
\begin{tikzpicture}[>=stealth']
\node (wmax) at (-180bp,0bp) {$w_{n-1}$};
\node (dots) at (-140bp,0bp) {$\cdots$};
\node (w2) at (-100bp,0bp) {$w_2$};
\node (w1) at (-50bp,0bp) {$w_1$};
\node (w0) at (0bp,0bp) {$w_0$};
\draw (wmax)--(dots);
\draw[->] (dots)--(w2);
\draw[->] (w2)--(w1);
\draw[->] (w1)--(w0);
\end{tikzpicture}
\end{center}

On définit aussi $n+1$ affectations de vérité $p_0,\ldots,p_n$ par
$p_k(w_i) = 1$ lorsque $i<k$ et $p_k(w_i) = 0$ lorsque $i\geq k$
(notamment, $p_0$ est l'affectation $\dottedbot$ et $p_n$ est
l'affectation $\dottedtop$).  Pourquoi sont-ce les seules affectations
de vérité sur ce cadre ?  Calculer les tableaux de $\dottedland,
\dottedlor, \dottedlimp, \dottedneg$ sur $p_0,\ldots,p_n$.

Donner un exemple de formule propositionnelle classiquement
démontrable et qui n'est pas validée par ce cadre (si $n\geq 2$), et
un exemple de formule propositionnelle validée par ce cadre et qui pas
démontrable intuitionnistement.

\begin{corrige}
La contrainte sur une affectation de vérité sur un cadre de Kripke $W$
est d'être permanente, i.e., si $p(w)=1$ alors $p(w')=1$ pour tout
monde $w'$ accessible depuis $w$ (c'est-à-dire $w\leq w'$), autrement
dit d'etre une fonction croissante $W \to \{0,1\}$.  Comme ici $W$ est
simplement $\{0,\ldots,n-1\}$ avec l'ordre inversé, les seules
fonctions décroissantes $\{0,\ldots,n-1\} \to \{0,1\}$ étant celles
qui valent $1$ jusqu'à un certain point et $0$ ensuite, on obtient les
$p_k$ qu'on a décrits.

Le calcul de $\dottedland$ et $\dottedlor$ est facile puisqu'il se
fait monde par monde : on a $(p_k \dottedland p_\ell)(w_i) = 1$
lorsque $p_k(w_i) = 1$ et $p_\ell(w_i) = 1$, c'est-à-dire $i<k$ et
$i<\ell$, ce qui équivaut à $i<\min(k,\ell)$, ce qui montre $p_k
\dottedland p_\ell = p_{\min(k,\ell)}$ ; de même, on a $(p_k
\dottedlor p_\ell)(w_i) = 1$ lorsque $p_k(w_i) = 1$ ou $p_\ell(w_i) =
1$, c'est-à-dire $i<k$ ou $i<\ell$, ce qui équivaut à
$i<\max(k,\ell)$, ce qui montre $p_k \dottedlor p_\ell =
p_{\max(k,\ell)}$.

Reste à traiter le cas de $\dottedlimp$.  Si $k\leq\ell$, alors
$p_k(w_j) = 1$ implique $p_\ell(w_j) = 1$ (car $j<k$ implique
$j<\ell$), et ce, pour tout $j$, ce qui montre $(p_k\dottedlimp
p_\ell) = \dottedtop = p_n$ dans ce cas.  En revanche, si $k>\ell$, il
faut distinguer deux cas : lorsque $j<\ell$, on a $p_k(w_j) = 1$ et
$p_\ell(w_j) = 1$, ce qui montre que $(p_k\dottedlimp p_\ell)(w_i) =
1$ pour $i<\ell$ ; mais comme $p_k(w_\ell) = 1$ et cependant
$p_\ell(w_\ell) = 0$, on a $(p_k\dottedlimp p_\ell)(w_i) = 0$ pour
$i\geq \ell$ ; c'est-à-dire que $(p_k\dottedlimp p_\ell) = p_\ell$
dans ce cas.  Enfin, $\dottedneg q$ est simplement $q\dottedlimp
\dottedbot$, c'est-à-dire $q\dottedlimp p_0$.

Pour résumer :
\[
\begin{aligned}
p_k \dottedland p_\ell &= p_{\min(k,\ell)} \\
p_k \dottedlor p_\ell &= p_{\max(k,\ell)} \\
(p_k \dottedlimp p_\ell) &=
\left\{
\begin{array}{ll}
p_n&\text{~si $k\leq\ell$}\\
p_\ell&\text{~si $k>\ell$}\\
\end{array}
\right.
\\
\dottedneg p_k &=
\left\{
\begin{array}{ll}
p_n&\text{~si $k=0$}\\
p_0&\text{~si $k>0$}\\
\end{array}
\right.
\end{aligned}
\]

La formule $A\lor\neg A$, classiquement démontrable, n'est pas validée
par ce cadre lorsque $n\geq 2$ (prendre $p_1$ pour $A$ : on trouve
$p_1 \dottedlor \dottedneg p_1 = p_1$).  La formule $\neg
A\lor\neg\neg A$, en revanche, est validée par ce cadre (car
$\dottedneg\dottedneg p_k = p_n$ si $k>0$ et $p_0$ si $k=0$) mais
n'est pas prouvable intuitionnistement (par la propriété de la
disjonction : si elle était prouvable, soit $\neg A$ soit $\neg\neg A$
le serait, or elles ne sont même pas prouvables classiquement).
\end{corrige}

%

\exercice\ (${\star}{\star}{\star}{\star}$)\par\nobreak

En prenant connaissance du résultat de
l'exercice \ref{exercise-computably-inseparable-sets}, montrer que la
formule suivante (« axiome de Kreisel-Putnam ») n'est pas réalisable :
\[
(\neg A \Rightarrow B\lor C) \Rightarrow
(\neg A\Rightarrow B) \lor (\neg A\Rightarrow C)
\]
(On pourra supposer par l'absurde qu'il y a un programme $r$ qui
réalise cette formule, et chercher à s'en servir pour séparer les
ensembles $L$ et $M$ définis dans
l'exercice \ref{exercise-computably-inseparable-sets}.
\emph{Indication :} plus précisément, on pourra poser $B_z$ comme
valant $\mathbb{N}$ si $z\in L$ et $\varnothing$ sinon ; $C_z$ comme
valant $\mathbb{N}$ si $z\in M$ et $\varnothing$ sinon ; et $A_z$
comme valant $\varnothing$ si $z\in L\cup M$ et $\mathbb{N}$ sinon ;
et chercher à définir un élément de $(\dottedneg A_z \dottedlimp
B_z\dottedlor C_z)$ auquel appliquer $r$.)

\begin{corrige}
Supposons par l'absurde qu'il existe (un \emph{même}) $r$ qui réalise
$(\neg A \Rightarrow B\lor C) \Rightarrow (\neg A\Rightarrow B) \lor
(\neg A\Rightarrow C)$, c'est-à-dire qui appartienne à $(\dottedneg A
\dottedlimp B\dottedlor C) \dottedlimp (\dottedneg A\dottedlimp B)
\dottedlor (\dottedneg A\dottedlimp C)$ quels que soient
$A,B,C\subseteq \mathbb{N}$ (où $P \dottedlor Q = \{\langle 0,m\rangle
: m\in P\} \cup \{\langle 1,n\rangle : n\in Q\}$ et $(P \dottedlimp Q)
= \{e \in \mathbb{N} : \varphi_e(P){\downarrow} \subseteq Q\}$, et on
se rappelle bien sûr que $\dottedneg P = (P\dottedlimp \varnothing)$
vaut $\mathbb{N}$ si $P$ est vide et $\varnothing$ si $P$ n'est pas
vide).

On pose $L := \{\langle e,x\rangle : \varphi_e(x){\downarrow} = 1\}$
et $M := \{\langle e,x\rangle : \varphi_e(x){\downarrow} = 2\}$.  On a
vu dans l'exercice \ref{exercise-computably-inseparable-sets} qu'il
n'existe aucun programme $s$ qui, prenant en entrée un
$z\in\mathbb{N}$, termine toujours en temps fini, et renvoie « vrai »
lorsque $z \in L$ et « faux » lorsque $z \in M$ (et une réponse
quelconque, lorsque $z \not\in L\cup M$, mais le programme doit quand
même terminer).  Or on va utiliser $r$ pour faire précisément une
telle chose, ce qui constituera une contradiction.

Définissons les ensembles $A_z,B_z,C_z$ pour $z = \langle e,x\rangle$
comme suit :
\[
\begin{aligned}
A_z &= \varnothing \text{~si~$\varphi_e(x){\downarrow} = 1$ ou
  $\varphi_e(x){\downarrow} = 2$,}\\
&= \mathbb{N} \text{~sinon}\\
\text{(donc~}\dottedneg A_z &= \mathbb{N} \text{~si~$\varphi_e(x){\downarrow} = 1$ ou
  $\varphi_e(x){\downarrow} = 2$,}\\
&= \varnothing \text{~sinon)}\\
B_z &= \mathbb{N} \text{~si~$\varphi_e(x){\downarrow} = 1$,}\\
&= \varnothing \text{~sinon}\\
C_z &= \mathbb{N} \text{~si~$\varphi_e(x){\downarrow} = 2$,}\\
&= \varnothing \text{~sinon}\\
\end{aligned}
\]

Donné $z\in \mathbb{N}$, considérons le programme $p_z$ défini comme
suit.  Il prend un unique argument, qu'il ignore.  Il décode le couple
$z = \langle e,x\rangle$, puis il exécute $\varphi_e(x)$ (au moyen
d'un interpréteur universel).  Si $\varphi_e(x){\downarrow} = 1$, il
renvoie $\langle 0,0\rangle$, et si $\varphi_e(x){\downarrow} = 2$, il
renvoie $\langle 1,0\rangle$ ; dans tout autre cas, il fait une boucle
infinie (y compris bien sûr si $\varphi_e(x){\uparrow}$, c'est-à-dire
si l'exécution ne termine jamais : dans ce cas forcément $p_z$ ne
termine pas non plus).

On remarque que $p_z$ est construit algorithmiquement en fonction
de $z$ (par le théorème s-m-n si on veut).

Par construction, si on passe à $p_z$ un argument de $\dottedneg A_z$,
ce qui implique que $\varphi_e(x){\downarrow} = 1$ ou
$\varphi_e(x){\downarrow} = 2$ (l'argument est une \emph{promesse} de
ce fait), alors il renvoie un élément de $B_z \dottedlor C_z$ (à
savoir un couple $\langle 0,n\rangle$ avec $n\in B_z$ ou $\langle
1,n\rangle$ avec $n\in C_z$).  On a donc $p_z \in (\dottedneg A_z
\dottedlimp B_z\dottedlor C_z)$.

Par hypothèse sur $r$, on doit donc avoir $\varphi_r(p_z)$ défini et
appartenant à $(\dottedneg A_z\dottedlimp B_z) \dottedlor (\dottedneg
A_z\dottedlimp C_z)$.  Notamment, $\varphi_r(p_z)$ est (le codage
d')un couple dont la première coordonnée vaut $0$ ou $1$ et indique si
la seconde est dans $\dottedneg A_z\dottedlimp B_z$ ou dans
$\dottedneg A_z\dottedlimp C_z$.  On considère le programme $s$ qui
prend un $z$ en entrée, calcule $\varphi_r(p_z)$ (noter que ce calcul
est bien algorithmique puisque $p_z$ est construit algorithmiquement
en fonction de $z$, et il termine toujours d'après ce qu'on vient de
dire), et renvoie « vrai » si la première coordonnée du résultat
vaut $0$ et « faux » si la seconde coordonnée du résultat vaut $1$.

Si $z = \langle e,x\rangle$ avec $\varphi_e(x){\downarrow} = 1$, alors
$C_z = \varnothing$ et $\dottedneg A_z = \mathbb{N}$, donc
$(\dottedneg A_z\dottedlimp C_z) = \varnothing$, donc $\varphi_r(p_z)$
doit forcément être de la forme $\langle 0,\cdots\rangle$, et $s$
renvoie « vrai » sur l'entrée $z$.  Symétriquement, si $z = \langle
e,x\rangle$ avec $\varphi_e(x){\downarrow} = 2$, alors $(\dottedneg
A_z\dottedlimp B_z) = \varnothing$ donc $\varphi_r(p_z)$ doit
forcément être de la forme $\langle 1,\cdots\rangle$, et $s$ renvoie
« faux » sur l'entrée $z$.

Donc $s$ termine toujours, et sépare $L$ et $M$ puisqu'il renvoie
« vrai » sur le premier et « faux » sur le second.  Ceci contredit le
fait que $L$ et $M$ sont récursivement inséparables.  On a abouti à
une contradiction : c'est qu'en fait $r$ n'existait pas.

(La formule de Kreisel-Putnam n'est donc pas réalisable.  En
particulier, par \emph{correction} de la sémantique de la
réalisabilité, elle n'est pas démontrable dans le calcul
propositionnel intuitionniste, c'est-à-dire, par Curry-Howard, qu'il
n'y a pas de terme du $\lambda$-calcul simplement typé ayant pour type
$(\neg A \Rightarrow B\lor C) \Rightarrow (\neg A\Rightarrow B) \lor
(\neg A\Rightarrow C)$.)
\end{corrige}

%

\exercice\ (${\star}{\star}{\star}$)\par\nobreak

\textbf{(1)} Donné un problème fini $(X,S)$, décrire soigneusement le
problème $\dottedneg(X,S)$ (combien de candidats il a, et combien de
solutions : on distinguera $S = \varnothing$ ou $S \neq\varnothing$).

\textbf{(2)} Expliquer pourquoi la formule suivante (« axiome de
Kreisel-Putnam ») est valable dans la sémantique de Medvedev des
problèmes finis :
\[
(\neg A \Rightarrow B\lor C) \Rightarrow
(\neg A\Rightarrow B) \lor (\neg A\Rightarrow C)
\]
(On décrira explicitement le problème fini décrit par la partie gauche
et par la partie droite de l'implication centrale de cette formule, et
leurs solutions.)

\begin{corrige}
\textbf{(1)} Soit $(X,S)$ un problème fini.  On rappelle que
$\dottedbot = (\{\bullet\},\varnothing)$.  Le problème
$\dottedneg(X,S)$, c'est-à-dire $(X,S) \dottedlimp
(\{\bullet\},\varnothing)$ a une seul candidat, à savoir l'unique
fonction $X \to \{\bullet\}$, qu'on notera abusivement aussi
$\bullet$ ; les solutions de $\dottedneg(X,S)$ sont les fonctions $X
\to \{\bullet\}$ qui envoient chaque élément de $S$ dans
$\varnothing$ : or si $S = \varnothing$ c'est le cas de l'unique
fonction $\bullet$, tandis que si $S \neq \varnothing$, elle n'envoie
pas $S$ dans $\varnothing$.  Pour résumer, $\dottedneg(X,S)$ a
toujours exactement un candidat $\bullet$, et ce candidat est solution
exactement lorsque le problème $(X,S)$ de départ a n'a pas de
solution.

\textbf{(2)} Considérons d'abord $\dottedneg A \dottedlimp B\dottedlor
C$ en remplaçant $A,B,C$ par trois problèmes finis $(X,S)$, $(Y,T)$ et
$(Z,U)$ respectivement.  On a vu que $\dottedneg A$ a toujours
exactement un candidat $\bullet$, et qu'il est solution précisément
lorsque $S = \varnothing$, et sinon aucune.  Les candidats de
$\dottedneg A \dottedlimp B\dottedlor C$ sont donc des fonctions de
$\{\bullet\}$ vers l'ensemble des candidats $Y \uplus Z$ de
$B\dottedlor C$, qu'on peut donc identifier à $Y \uplus Z$ (en
identifiant une fonction $\{\bullet\} \to V$, pour $V$ un ensemble
quelconque, avec l'image de $\bullet$ par cette fonction).  Parmi ces
candidats, si $S \neq\varnothing$, ils sont tous solutions (car
$\dottedneg A$ n'a pas de solution donc il n'y a pas de contrainte),
tandis que si $S = \varnothing$, ceux qui sont solution sont ceux
solution de $B$ et de $C$, i.e., c'est $T\uplus U$.

Or la même description vaut pour $(\dottedneg A\dottedlimp B)
\dottedlor (\dottedneg A\dottedlimp C)$ : le problème $(\dottedneg
A\dottedlimp B)$ a un ensemble de candidats qui s'identifie à celui
$Y$ de $B$, et si $S \neq\varnothing$ ils sont tous solutions tandis
que si $S = \varnothing$ ceux qui sont solutions sont les solutions
de $B$, c'est-à-dire $T$ ; et le problème $(\dottedneg A\dottedlimp
C)$ a un ensemble de candidats qui s'identifie à celui $Z$ de $C$, et
si $S \neq\varnothing$ ils sont tous solutions tandis que si $S =
\varnothing$ ceux qui sont solutions sont les solutions de $C$,
c'est-à-dire $V$.  Bref, $(\dottedneg A\dottedlimp B) \dottedlor
(\dottedneg A\dottedlimp C)$ a pour ensemble de candidats $Y \uplus Z$
et pour ensemble de solutions $Y\uplus Z$ lorsque $S \neq\varnothing$
et $T\uplus U$ lorsque $S = \varnothing$.  C'est exactement pareil que
$\dottedneg A \dottedlimp B\dottedlor C$.

On a donc un candidat évident dans $(\dottedneg A \dottedlimp
B\dottedlor C) \dottedlimp (\dottedneg A\dottedlimp B) \dottedlor
(\dottedneg A\dottedlimp C)$, c'est l'identité (une fois faites nos
identifications des candidats de $\dottedneg A \dottedlimp B\dottedlor
C$ comme $(\dottedneg A\dottedlimp B) \dottedlor (\dottedneg
A\dottedlimp C)$ avec $Y\uplus Z$), et par la description faite des
solutions, ce candidat est toujours solution.  C'est précisément ce
qui montre que $(\neg A \Rightarrow B\lor C) \Rightarrow (\neg
A\Rightarrow B) \lor (\neg A\Rightarrow C)$ est valide dans la
sémantique des problèmes finis de Medvedev.
\end{corrige}

%

\exercice\ (${\star}$)\par\nobreak

En considérant les parties suivantes de $\mathbb{R}^2$ :
\[
\begin{array}{c@{\hskip 3em}c}
U_1 = \{(x,y)\in\mathbb{R}^2 : x < 0\}
& U_2 = \{(x,y)\in\mathbb{R}^2 : x > 0\}
\\
V_1 = \{(x,y)\in\mathbb{R}^2 : y > -x^2\}
& V_2 = \{(x,y)\in\mathbb{R}^2 : y < x^2\}
\end{array}
\]
(faire un dessin !) montrer que la formule de Tseitin
\[
\begin{aligned}
&(\neg (A_1 \land A_2) \land (\neg A_1 \Rightarrow (B_1 \lor B_2))
  \land (\neg A_2 \Rightarrow (B_1 \lor B_2)))\\
\Rightarrow\,& ((\neg A_1 \Rightarrow B_1)\lor(\neg A_2 \Rightarrow B_1)\lor(\neg
  A_1 \Rightarrow B_2)\lor(\neg A_2 \Rightarrow B_2))
\end{aligned}
\]
n'est pas démontrable en calcul propositionnel intuitionniste.

\begin{corrige}
On va interpréter la formule de Tseitin dans la sémantique des ouverts
de $\mathbb{R}^2$ en prenant $U_1,U_2,V_1,V_2$ pour $A_1,A_2,B_1,B_2$
respectivement, c'est-à-dire qu'on va calculer :
\[
\begin{aligned}
&(\dottedneg (U_1 \dottedland U_2) \dottedland (\dottedneg U_1 \dottedlimp (V_1 \dottedlor V_2))
  \dottedland (\dottedneg U_2 \dottedlimp (V_1 \dottedlor V_2)))\\
\dottedlimp\,& ((\dottedneg U_1 \dottedlimp V_1)\dottedlor(\dottedneg U_2 \dottedlimp V_1)\dottedlor(\dottedneg
  U_1 \dottedlimp V_2)\dottedlor(\dottedneg U_2 \dottedlimp V_2))
\end{aligned}
\]
On a représenté ci-dessous en grisé les parties \emph{ouvertes}
décrites par différentes sous-expressions de cette formule :
\begin{center}
\lineskip=1explus5ex
\begin{tabular}{c}
\begin{tikzpicture}
\fill[gray] (-1.5,-1.5) -- (0,-1.5) -- (0,1.5) -- (-1.5,1.5) -- cycle;
\draw[->] (-1.55,0)--(1.55,0) node[right]{$x$};
\draw[->] (0,-1.55)--(0,1.55) node[above]{$y$};
\end{tikzpicture}
\\ $U_1 \;=\; \dottedneg U_2$
\end{tabular}
\penalty0
\begin{tabular}{c}
\begin{tikzpicture}
\fill[gray] (1.5,-1.5) -- (0,-1.5) -- (0,1.5) -- (1.5,1.5) -- cycle;
\draw[->] (-1.55,0)--(1.55,0) node[right]{$x$};
\draw[->] (0,-1.55)--(0,1.55) node[above]{$y$};
\end{tikzpicture}
\\ $U_2 \;=\; \dottedneg U_1$
\end{tabular}
\penalty0
\begin{tabular}{c}
\begin{tikzpicture}
\begin{scope}
\clip (-1.5,-1.5) -- (1.5,-1.5) -- (1.5,1.5) -- (-1.5,1.5) -- cycle;
\fill[gray] (-1.5,-2.25) .. controls (-1,-0.75) and (-0.5,0) .. (0,0) .. controls (0.5,0) and (1,-0.75) .. (1.5,-2.25) -- (1.5,1.5) -- (-1.5,1.5) -- cycle;
\end{scope}
\draw[->] (-1.55,0)--(1.55,0) node[right]{$x$};
\draw[->] (0,-1.55)--(0,1.55) node[above]{$y$};
\end{tikzpicture}
\\ $V_1$
\end{tabular}
\penalty0
\begin{tabular}{c}
\begin{tikzpicture}
\begin{scope}
\clip (-1.5,-1.5) -- (1.5,-1.5) -- (1.5,1.5) -- (-1.5,1.5) -- cycle;
\fill[gray] (-1.5,2.25) .. controls (-1,0.75) and (-0.5,0) .. (0,0) .. controls (0.5,0) and (1,0.75) .. (1.5,2.25) -- (1.5,-1.5) -- (-1.5,-1.5) -- cycle;
\end{scope}
\draw[->] (-1.55,0)--(1.55,0) node[right]{$x$};
\draw[->] (0,-1.55)--(0,1.55) node[above]{$y$};
\end{tikzpicture}
\\ $V_2$
\end{tabular}
\penalty0
\begin{tabular}{c}
\begin{tikzpicture}
\begin{scope}
\clip (-1.5,-1.5) -- (1.5,-1.5) -- (1.5,1.5) -- (-1.5,1.5) -- cycle;
\fill[gray] (-1.5,-1.5) -- (0,-1.5) -- (0,0) .. controls (0.5,0) and (1,-0.75) .. (1.5,-2.25) -- (1.5,1.5) -- (-1.5,1.5) -- cycle;
\end{scope}
\draw[->] (-1.55,0)--(1.55,0) node[right]{$x$};
\draw[->] (0,-1.55)--(0,1.55) node[above]{$y$};
\end{tikzpicture}
\\ $\dottedneg U_1 \dottedlimp V_1$
\end{tabular}
\penalty0
\begin{tabular}{c}
\begin{tikzpicture}
\begin{scope}
\clip (-1.5,-1.5) -- (1.5,-1.5) -- (1.5,1.5) -- (-1.5,1.5) -- cycle;
\fill[gray] (-1.5,-2.25) .. controls (-1,-0.75) and (-0.5,0) .. (0,0) -- (0,-1.5) -- (1.5,-1.5) -- (1.5,1.5) -- (-1.5,1.5) -- cycle;
\end{scope}
\draw[->] (-1.55,0)--(1.55,0) node[right]{$x$};
\draw[->] (0,-1.55)--(0,1.55) node[above]{$y$};
\end{tikzpicture}
\\ $\dottedneg U_2 \dottedlimp V_1$
\end{tabular}
\penalty0
\begin{tabular}{c}
\begin{tikzpicture}
\begin{scope}
\clip (-1.5,-1.5) -- (1.5,-1.5) -- (1.5,1.5) -- (-1.5,1.5) -- cycle;
\fill[gray] (-1.5,1.5) -- (0.1,5) -- (0,0) .. controls (0.5,0) and (1,0.75) .. (1.5,2.25) -- (1.5,-1.5) -- (-1.5,-1.5) -- cycle;
\end{scope}
\draw[->] (-1.55,0)--(1.55,0) node[right]{$x$};
\draw[->] (0,-1.55)--(0,1.55) node[above]{$y$};
\end{tikzpicture}
\\ $\dottedneg U_1 \dottedlimp V_2$
\end{tabular}
\penalty0
\begin{tabular}{c}
\begin{tikzpicture}
\begin{scope}
\clip (-1.5,-1.5) -- (1.5,-1.5) -- (1.5,1.5) -- (-1.5,1.5) -- cycle;
\fill[gray] (-1.5,2.25) .. controls (-1,0.75) and (-0.5,0) .. (0,0) -- (0,1.5) -- (1.5,1.5) -- (1.5,-1.5) -- (-1.5,-1.5) -- cycle;
\end{scope}
\draw[->] (-1.55,0)--(1.55,0) node[right]{$x$};
\draw[->] (0,-1.55)--(0,1.55) node[above]{$y$};
\end{tikzpicture}
\\ $\dottedneg U_2 \dottedlimp V_2$
\end{tabular}
\end{center}
Le membre de droite $((\dottedneg U_1 \dottedlimp
V_1)\dottedlor(\dottedneg U_2 \dottedlimp V_1)\dottedlor(\dottedneg
U_1 \dottedlimp V_2)\dottedlor(\dottedneg U_2 \dottedlimp V_2))$ du
$\dottedlimp$ externe de la formule est la réunion des quatre
dernières parties dessinées ci-dessus, c'est-à-dire le complémentaire
de l'origine (noter qu'on n'attrape jamais l'origine puisqu'elle est
au bord de chacune des parties dessinées, donc jamais dedans).  En
revanche, chacun des facteurs, $\dottedneg (U_1 \dottedland U_2)$,
$\dottedneg U_1 \dottedlimp (V_1 \dottedlor V_2)$ et $\dottedneg U_2
\dottedlimp (V_1 \dottedlor V_2)$ du membre de gauche est
$\mathbb{R}^2$ tout entier, puisque $U_1 \dottedland U_2$ est vide, et
que $V_1 \dottedlor V_2$ est le complémentaire de l'origine donc
contient $\dottedneg U_1$ comme $\dottedneg U_2$.  Finalement, pour
l'expression tout entière, on trouve le complémentaire de l'origine,
qui n'est pas $\mathbb{R}^2$, donc la formule de Tseitin n'est pas
validée par ce choix d'ouverts, et (par \emph{correction} de la
sémantique des ouverts) elle ne peut pas être démontrable.
\end{corrige}


%
%
%

\section{Quantificateurs}

\exercice\label{exercise-system-f}\ (${\star}{\star}{\star}$)\par\nobreak

On appelle \textbf{système F} de Girard (dit aussi $\lambda 2$), ou
plus exactement le système F étendu par types produits, sommes, 1, 0
et quantificateur existentiel, le système de logique et/ou typage de
la manière décrite ci-dessous.  L'idée générale est que le système F
ajoute au calcul propositionnel intuitionniste la capacité à
quantifier sur des propositions ; ou, si on voit ça comme un système
de typage, un polymorphisme explicite.

On part des règles du $\lambda$-calcul simplement typé étendu par
types produits, sommes, 1, 0, qu'on notera avec les notations logiques
($\land,\lor,\top,\bot$), ou, ce qui revient au même, du calcul
propositionnel intuitionniste.  On ajoute un symbole spécial $*$ pour
représenter le « type des propositions » (ou types) avec les règles de
typage suivantes définissant la construction des propositions (ou
types) :
\begin{center}
$\inferrule{\Gamma\vdash P:*\\\Gamma\vdash Q:*}{\Gamma\vdash P\Rightarrow Q:*}$
\penalty0\quad
$\inferrule{\Gamma\vdash P:*\\\Gamma\vdash Q:*}{\Gamma\vdash P\land Q:*}$
\penalty0\quad
$\inferrule{\Gamma\vdash P:*\\\Gamma\vdash Q:*}{\Gamma\vdash P\lor Q:*}$
\penalty-100\quad
$\inferrule{ }{\Gamma\vdash \top:*}$
\penalty0\quad
$\inferrule{ }{\Gamma\vdash \bot:*}$
\end{center}
On ajoute des règles permettant de quantifier sur $*$ :
\begin{center}
$\inferrule{\Gamma,\;T:*\;\vdash\; Q:*}{\Gamma\;\vdash\; (\forall(T:*).\,Q):*}$
\penalty0\quad
$\inferrule{\Gamma,\;T:*\;\vdash\; Q:*}{\Gamma\;\vdash\; (\exists(T:*).\,Q):*}$
\end{center}
Et on ajoute les règles d'introduction et d'élimination de des
quantificateurs :
\begin{center}
$\inferrule{\Gamma,\; T:*\;\vdash\; s:Q}{\Gamma\vdash\; (\lambda(T:*).\,s)\;:\;(\forall (T:*).\, Q)}$
\penalty0\quad
$\inferrule{\Gamma\vdash\; f \;:\; (\forall (T:*).\, Q)\\\Gamma\vdash P:*}{\Gamma\vdash\; fP \;:\; Q[T\backslash P]}$
\penalty0\quad
$\inferrule{\Gamma\vdash P:*\\\Gamma\vdash\; z \;:\; Q[T\backslash P]}{\Gamma\vdash\; \langle P,z\rangle \;:\; (\exists (T:*).\, Q)}$
\penalty0\quad
$\inferrule{\Gamma\vdash\; z \;:\; (\exists (T:*).\, P)\\\Gamma,\, T:*,\, h:P\;\vdash\; Q}{\Gamma\vdash (\texttt{match~}z\texttt{~with~}\langle T,h\rangle \mapsto s) \;:\; Q}$
\end{center}

Pour abréger les notations, on écrit souvent, et on le fera dans la
suite, « $\Lambda T$ » plutôt que « $\lambda(T:*)$ » (abstraction sur
les propositions/types), et « $\forall T$ » / « $\exists T$ » plutôt
que « $\forall(T:*)$ » / « $\exists(T:*)$ » (puisque c'est la seule
forme de quantification autorisée par le système F).

À titre d'exemple, dans le système F, on peut écrire le terme $\Lambda
A.\, \Lambda B.\, \lambda(x:A).\, \lambda(y:B).\, x$, qui prouve (ou a
pour type) $\forall A.\, \forall B.\, (A\Rightarrow B\Rightarrow A)$,
ou encore $\Lambda A.\, \lambda(f:\forall B.B).\, fA$ qui prouve (ou a
pour type) $\forall A.\, ((\forall B.\, B) \Rightarrow A)$.

\smallskip

\textbf{(1)} Montrer (en donnant des $\lambda$-termes) dans le
système F que $\forall A.\,(A \Rightarrow \forall C.((A\Rightarrow
C)\Rightarrow C))$ et que $\forall A.\,((\forall C.((A\Rightarrow
C)\Rightarrow C)) \Rightarrow A)$.  Si on voit ces termes comme des
programmes ayant les types en question, décrire brièvement ce que font
ces programmes.

\textbf{(2)} Montrer de même que $\forall A.\,\forall B.\,(A\land B
\Rightarrow \forall C.((A\Rightarrow B\Rightarrow C)\Rightarrow C))$
et que $\forall A.\,\forall B.\,((\forall C.((A\Rightarrow
B\Rightarrow C)\Rightarrow C)) \Rightarrow A\land B)$.  Quel rapport
avec l'exercice \ref{exercise-coding-of-pairs-untyped} ?

\textbf{(3)} La question précédente indique que dans le système F,
$P\land Q$ peut être considéré comme un synonyme de $\forall C.\,
(P\Rightarrow Q\Rightarrow C)\Rightarrow C$ (ces deux propositions
sont équivalentes ; ceci ne montre pas tout à fait qu'on peut se
dispenser entièrement du signe $\land$ car il n'est pas évident qu'on
n'en ait pas besoin dans les démonstrations, mais il s'avère que c'est
le cas).  Proposer de façon semblable un synonyme de $P\lor Q$ ne
faisant intervenir que $\Rightarrow$ et $\forall$, et montrer
l'équivalence.

\textbf{(4)} Proposer une piste pour réécrire $\exists T.\,Q$ (où $T$
peut apparaître libre dans $Q$) en ne faisant intervenir que
$\Rightarrow$ et $\forall$.  (Comme le système F ne permet pas de
quantifier sur les propositions dépendant d'une autre
proposition\footnote{Il faudrait pour cela faire apparaître le type
$*\to *$ et quantifier dessus, ce qui dépasse le système F (le système
appelé « F$\omega$ » le permet).}, on se contentera d'écrire les
équivalences sur un $Q$ quelconque.)

\begin{corrige}
\textbf{(1)} La formule $\forall A.\,(A \Rightarrow \forall
C.((A\Rightarrow C)\Rightarrow C))$ est prouvée par le $\lambda$-terme
$\Lambda A.\, \lambda(x:A).\, \Lambda C.\, \lambda(k:A\Rightarrow
C).\, kx$.  Dans l'autre sens, la formule $\forall A.\,((\forall
C.((A\Rightarrow C)\Rightarrow C)) \Rightarrow A)$ est prouvée par le
$\lambda$-terme $\Lambda A.\, \lambda(h:\forall C.((A\Rightarrow
C)\Rightarrow C)).\, hA(\lambda(x:A).x)$.

Si on voit ces termes comme des programmes, le premier qu'on vient de
décrire prend un type $A$ et une valeur $x$ de ce type et convertit
cette valeur en un « passage par continuation » (ou « cachée dans une
clôture »), c'est-à-dire en une fonction qui prend un type $C$ et une
fonction $k$ (qu'il faut imaginer comme une sorte continuation) et
passe $x$ comme argument à $k$.  Le second fait la conversion
inverse : il prend un type $A$ et une valeur $h$ passée par
continuation et extrait la valeur en appliquant $h$ au type $A$
lui-même et à la continuation-de-fait donnée par la fonction identité.

\textbf{(2)} La formule $\forall A.\,\forall B.\,(A\land B \Rightarrow
\forall C.((A\Rightarrow B\Rightarrow C)\Rightarrow C))$ est prouvée
par $\Lambda A.\, \Lambda B.\, \lambda(z:A\land B).\, \Lambda C.\,
\lambda(k:A\Rightarrow B\Rightarrow C).\, k(\pi_1 z)(\pi_2 z)$.  Dans
l'autre sens, la formule $\forall A.\,\forall B.\,((\forall
C.((A\Rightarrow B\Rightarrow C)\Rightarrow C)) \Rightarrow A\land B)$
est prouvée par $\Lambda A.\, \Lambda B.\, \lambda(h:\forall
C.((A\Rightarrow B\Rightarrow C)\Rightarrow C)).\, h(A\land
B)(\lambda(x:A).\lambda(y:B).\langle x,y\rangle)$.

Si on voit ces termes comme des programmes, le premier prend en entrée
deux types et un couple dont les coordonnées ont ces deux types, et le
transforme en la fonction qui attend une fonction de deux arguments et
l'appelle avec les deux coordonnées du couple comme arguments
successifs : c'est la fonction notée \texttt{fromnative} dans la
correction de l'exercice \ref{exercise-coding-of-pairs-untyped} (i.e.,
elle prend un couple « natif » et le convertit en sa représentation
sous forme de clôture).  Le second programme fait la conversion
inverse et correspond à la fonction notée \texttt{tonative} dans la
correction de l'exercice \ref{exercise-coding-of-pairs-untyped}.  La
subtilité évoquée dans la correction dudit exercice revient à observer
que le typage de OCaml n'est pas aussi riche que le système F.

\textbf{(3)} L'idée est que pour passer un type somme $P\lor Q$ par
continuation on va recevoir deux fonctions, l'une invoquée dans le
cas $P$ et l'autre dans le cas $Q$.  On peut aussi prendre son
inspiration d'un des axiomes du système de Hilbert $(A\Rightarrow C)
\Rightarrow (B\Rightarrow C)\Rightarrow A\lor B\Rightarrow C$.  Bref,
on va représenter $P\lor Q$ par $\forall C.\,(P\Rightarrow C)
\Rightarrow (Q\Rightarrow C)\Rightarrow C$.

La formule $\forall A.\,\forall B.\,(A\lor B \Rightarrow
\forall C.((A\Rightarrow C) \Rightarrow (B\Rightarrow C)\Rightarrow
C))$ est prouvée par $\Lambda A.\, \Lambda B.\, \lambda(z:A\lor B).\,
\Lambda C.\, \lambda(k:A\Rightarrow C).\, \lambda(\ell:B\Rightarrow
C).\, (\texttt{match~}z\texttt{~with~}\iota_1 x \mapsto kx,\; \iota_2
y \mapsto \ell y)$.  Dans l'autre sens, la formule $\forall
A.\,\forall B.\,((\forall C.((A\Rightarrow C)\Rightarrow (B\Rightarrow
C)\Rightarrow C)) \Rightarrow A\lor B)$ est prouvée par $\Lambda A.\,
\Lambda B.\, \lambda(h:\forall C.((A\Rightarrow C)\Rightarrow
(B\Rightarrow C)\Rightarrow C)).\, h(A\lor
B)(\lambda(x:A).\iota_1^{(A,B)} x) (\lambda(y:B).\iota_2^{(A,B)} y)$.

\textbf{(4)} Enfin, la quantification existentielle $\exists T.\,Q(T)$
va être représentée par $\forall C.\,((\forall U. (Q(U)\Rightarrow
C))\Rightarrow C)$.  La formule $(\exists T. Q(T)) \Rightarrow \forall
C.((\forall U. (Q(U)\Rightarrow C))\Rightarrow C)$ est prouvée par
$\lambda(z:\exists T. Q(T)).\, \Lambda C.\, \lambda(k:\forall
U. (Q(U)\Rightarrow C)).\, (\texttt{match~}z\texttt{~with~}\langle
T,t\rangle \mapsto kTt)$.  Dans l'autre sens, la formule $(\forall
C.((\forall U. (Q(U)\Rightarrow C))\Rightarrow C)) \Rightarrow
(\exists T. Q(T))$ est prouvée par $\lambda(h:\forall C.((\forall
U. (Q(U)\Rightarrow C))\Rightarrow C)).\, h(\exists T. Q(T)) (\Lambda
U. \lambda(u:U). \langle U,u\rangle)$.

\emph{Remarque :} grâce à ces manipulations, on peut réécrire tout
terme du système F pour ne faire intervenir que $\Rightarrow$
et $\forall$ (il est facile de voir que $\bot$ se réécrit comme
$\forall C.C$, et $\top$ comme $\forall C.(C\Rightarrow C)$).  On peut
même montrer que le système F limité à ces seuls connecteurs est
équivalent à celui que nous avons défini, et c'est la définition la
plus habituelle de ce lui-ci.  Le système ainsi obtenu est plus
économique et plus élégant car il ne comporte comme seule
constructions des termes que l'application et les deux niveaux
d'abstraction, $\lambda$ et $\Lambda$.

\emph{Remarque 2 :} les règles de Coq incorporent notamment celles du
système F (même F$\omega$).  Par exemple, les termes que nous avons
construits à la question (2) peuvent se décrire ainsi :

{\tt\noindent
Theorem fromnative : forall (A B:Prop), (A/\textbackslash B) -> forall (C:Prop), ((A->B->C)->C).\\
Proof.  intros A B z C k.  destruct z as [x y].  exact (k x y).  Qed.\\
Theorem tonative : forall (A B:Prop), (forall (C:Prop), ((A->B->C)->C)) -> A/\textbackslash B.\\
Proof.  intros A B h.  apply (h (A/\textbackslash B)).  intros x y.  split.  exact x.  exact y.  Qed.
\par}

Les $\lambda$-termes s'écrivent alors, en syntaxe
Coq :\penalty0\hskip1emplus5em \texttt{fun (A B : Prop) (z : A
  /\textbackslash\ B) (C : Prop) (k : A -> B -> C) => match z with
  conj x y => k x y end}\penalty0\hskip1emplus5em pour le premier
et\penalty0\hskip1emplus5em \texttt{fun (A B : Prop) (h : forall C :
  Prop, (A -> B -> C) -> C) => h (A /\textbackslash\ B) (fun (x : A)
  (y : B) => conj x y)}\penalty0\hskip1emplus5em pour le second, ce
qui correspond bien à ce qui a été proposé ci-dessus.
\end{corrige}

%

\exercice\ (${\star}$)\par\nobreak

\textbf{(1)} Montrer chacune des propositions suivantes en pure
logique du premier ordre (où $A,B$ désignent des relations unaires) en
donnant un $\lambda$-terme de preuve : \textbf{(a)} $(\forall
x. A(x)\Rightarrow B(x)) \Rightarrow (\forall x. A(x)) \Rightarrow
(\forall x. B(x))$\hskip 1emplus5em \textbf{(b)} $(\forall
x. A(x)\Rightarrow B(x)) \Rightarrow (\exists x. A(x)) \Rightarrow
(\exists x. B(x))$\hskip 1emplus5em \textbf{(c)} $(\exists
x. A(x)\Rightarrow B(x)) \Rightarrow (\forall x. A(x)) \Rightarrow
(\exists x. B(x))$

\textbf{(2)} En interprétant chacun de ces termes comme un programme,
en en oubliant les types tels qu'on vient de les écrire, qu'obtient-on
si on demande à OCaml (c'est-à-dire en fait à l'algorithme de
Hindley-Milner étendu avec des types produits) de leur reconstruire
des types ?  Commenter brièvement.

\begin{corrige}
\textbf{(1)} \textbf{(a)} $\lambda(h:\forall x. A(x)\Rightarrow
B(x)).\, \lambda(u:\forall x. A(x)).\, \lambda(x:I).\, hx(ux)$\hskip
1emplus5em \textbf{(b)} $\lambda(h:\forall x. A(x)\Rightarrow B(x)).\,
\lambda(v:\exists x. A(x)).\, (\texttt{match~}v\texttt{~with~}\langle
z,w\rangle \mapsto \langle z,hzw\rangle)$\hskip 1emplus5em
\textbf{(c)} $\lambda(p:\exists x. A(x)\Rightarrow B(x)).\,
\lambda(u:\forall x. A(x)).\, (\texttt{match~}p\texttt{~with~}\langle
z,q\rangle \mapsto \langle z,q(uz)\rangle)$

\textbf{(2)} On demande à OCaml de typer les trois expressions
\texttt{fun h -> fun u -> fun x -> h x (u x)} et \texttt{fun h -> fun
  v -> match v with (z,w) -> (z, h z w)} et \texttt{fun p -> fun u ->
  match p with (z,q) -> (z, q(u z))} : les réponses, réécrites avec
des variables et des lettres collant mieux avec ce qu'on a trouvé
ci-desus, sont : $(I \to A \to B) \to (I\to A) \to (I\to B)$ pour le
premier, $(I\to A\to B) \to (I\times A) \to (I\times B)$ pour le
second, et $I\times(A\to B) \to (I\to A) \to (I\times C)$ pour le
troisième.  Ces types sont, forcément, moins précis que les types
$(\prod_{x:I} A(x) \to B(x)) \to (\prod_{x:I} A(x)) \to (\prod_{x:I}
B(x))$ et $(\prod_{x:I} A(x) \to B(x)) \to (\sum_{x:I} A(x)) \to
(\sum_{x:I} B(x))$ et $(\sum_{x:I} A(x) \to B(x)) \to (\prod_{x:I}
A(x)) \to (\sum_{x:I} B(x))$ des types que nous avons écrits à la
question (1), puisque OCaml ne permet pas les types dépendants.

(On peut aussi remarquer au passage que le premier terme $\lambda
hux.hx(ux)$, détypé, est le même que le combinateur $\mathsf{S}$ : le
combinateur $\mathsf{S}$ est l'axiome permettant d'appliquer le
\textit{modus ponens} sous une implication, de même que notre premier
terme permet d'applique le \textit{modus ponens} sous un
quantificateur universel.)
\end{corrige}


%
%
%

\section{Arithmétique du premier ordre et théorème de Gödel}

\exercice\label{exercise-zero-leftneutral}\ (${\star}$)\par\nobreak

Démontrer dans l'arithmétique de Heyting que $\forall n.(0+n=n)$.

\begin{corrige}
Voici une preuve complète écrite dans le style « drapeau » :
\bgroup\normalsize
\begin{footnotesize}
\begin{flagderiv}[exercise-zero-leftneutral-proof]
\step{recurr}{(0+0=0) \Rightarrow\\\quad (\forall n.((0+n=n)\Rightarrow (0+Sn=Sn)))
\Rightarrow\\\quad (\forall n.(0+n=n))}{Instance du schéma de récurrence}
\step{zero-rightneutral}{\forall m.(m+0 = m)}{Un des axiomes de Peano}
\step{recurr-start}{0+0 = 0}{$\forall$Élim sur \ref{zero-rightneutral} et $0$}
\step{recurr-zero}{(\forall n.((0+n=n)\Rightarrow (0+Sn=Sn)))
\Rightarrow\\\quad (\forall n.(0+n=n))}{$\Rightarrow$Élim sur \ref{recurr} et \ref{recurr-start}}
\assume{varn}{n}{}
\assume{rhyp}{0+n=n}{}
\step{subst}{\forall m.\forall m'.((m=m') \Rightarrow\\\quad (0+Sn=Sm) \Rightarrow (0+Sn=Sm'))}{Instance du schéma de substitution de l'égalité}
\step{subst-1}{\forall m'.((0+n=m') \Rightarrow\\\quad (0+Sn=S(0+n)) \Rightarrow (0+Sn=Sm'))}{$\forall$Élim sur \ref{subst} et $0+n$}
\step{subst-2}{(0+n=n) \Rightarrow\\\quad (0+Sn=S(0+n)) \Rightarrow (0+Sn=Sn)}{$\forall$Élim sur \ref{subst-1} et $n$}
\step{subst-3}{(0+Sn=S(0+n)) \Rightarrow (0+Sn=Sn)}{$\Rightarrow$Élim sur \ref{subst-2} et \ref{rhyp}}
\step{plus-succ}{\forall m.\forall m'.(m+Sm' = S(m+m'))}{Un des axiomes de Peano}
\step{plus-succ-zero}{\forall m'.(0+Sm' = S(0+m'))}{$\forall$Élim sur \ref{plus-succ} et $0$}
\step{plus-succ-zero-n}{0+Sn = S(0+n)}{$\forall$Élim sur \ref{plus-succ-zero} et $n$}
\step{subst-4}{0+Sn=Sn}{$\Rightarrow$Élim sur \ref{subst-3} et \ref{plus-succ-zero-n}}
\conclude{recurr-induct}{(0+n=n)\Rightarrow (0+Sn=Sn)}{$\Rightarrow$Int de \ref{rhyp} dans \ref{subst-4}}
\conclude{recurr-induct-univ}{\forall n.((0+n=n)\Rightarrow (0+Sn=Sn))}{$\forall$Int de \ref{varn} dans \ref{recurr-induct}}
\step{mainconc}{\forall n.(0+n=n)}{$\Rightarrow$Élim sur \ref{recurr-zero} et \ref{recurr-induct-univ}}
\end{flagderiv}
\end{footnotesize}
\egroup

(On comprend pourquoi on écrit rarement de telles choses
complètement.)  Le $\lambda$-terme correspondant est le suivant :
$\mathsf{recurr}^{(\lambda k.(0+k=k))}\penalty500\,
(\mathsf{defplusz}\,0)\penalty0\, (\lambda(n:\mathsf{nat}).\,
\lambda(h:0+n=n).\penalty500\, \mathsf{subst}^{(\lambda
  k.(0+Sn=Sk))}(0+n)\,n\,h\penalty500\,(\mathsf{defplusn}\,0\,n))$ où
on a donné les noms aux axiomes $\mathsf{defplusz} \; : \; \forall
m.(m+0 = m)$ et $\mathsf{defplusn} \; : \; \forall m.\forall m'.(m+Sm'
= S(m+m'))$ et pour les instances de schémas
$\mathsf{recurr}^{(\lambda k. P(k))}$ pour $P(0) \Rightarrow (\forall
n.(P(n)\Rightarrow P(Sn))) \Rightarrow (\forall n.P(n))$ et
$\mathsf{subst}^{(\lambda k. P(k))}$ pour $\forall m.\forall n.((m=n)
\Rightarrow P(m) \Rightarrow P(n))$ (le $\lambda$ en exposant est ici
complètement conventionnel), et enfin $\mathsf{nat}$ pour marquer le
type des individus.  Si tant est que ce programme fasse quoi que ce
soit, c'est une boucle prenant $n$ fois le successeur de $0$ jusqu'à
obtenir un témoignage d'égalité de $0+n$ à $n$.
\end{corrige}

%

\exercice\label{exercise-loebs-theorem}\ (${\star}{\star}{\star}{\star}$)\par\nobreak

Soit $T$ une théorie logique telle
que\footnote{\label{footnote-goedelian-theory}Comme pour le théorème
de Gödel, il s'agit essentiellement que $T$ soit codable par des
entiers naturels, permette de formaliser l'arrêt des machines de
Turing, et ait des démonstrations algorithmiquement testables.  On ne
donnera pas de conditions exactes (ce serait trop fastidieux) mais les
théories proposées en exemple les vérifient.}  l'arithmétique de
Heyting $\mathsf{HA}$, l'arithmétique de Peano $\mathsf{PA}$, Coq ou
$\mathsf{ZFC}$.  On se propose ici de démontrer le \textbf{théorème de
  Löb} : si $A$ est un énoncé de $T$ et si $T$ prouve « si $T$
prouve $A$, alors $A$ », alors $T$ prouve $A$.

On construit pour cela le programme $g_A$ suivant : $g_A$ cherche une
preuve dans $T$ de l'énoncé « si $g_A$ termine, alors $A$ », et s'il
en trouve une, il termine immédiatement.

\textbf{(1)} Expliquer pourquoi $g_A$ est bien défini.  A-t-on besoin
d'arriver à tester la véracité ou la démontrabilité de $A$ pour le
construire ?

\textbf{(2)} Supposons que $g_A$ termine : montrer que $A$ est
démontrable dans $T$.

\textbf{(3)} Remarquer que le point précédent est lui-même démontrable
dans $T$.

\textbf{(4)} Déduire de (3) que si $T$ prouve « si $T$ prouve $A$,
alors $A$ », alors $g_A$ termine.

\textbf{(5)} Déduire de (2) et (4) que si $T$ prouve « si $T$
prouve $A$, alors $A$ », alors $T$ provue $A$.

\begin{corrige}
\textbf{(1)} Le programme $g_A$ est, comme d'habitude, bien défini
grâce à l'astuce de Quine : il s'agit de construire le programme
$h(e)$ qui cherche une preuve dans $T$ de l'énoncé « si $e$ termine,
alors $A$ » et, s'il en trouve une, termine immédiatement, et de lui
appliquer le théorème de récursion de Kleene.  Or l'énoncé entre
guillemets peut être formalisé de façon algorithmique, et la
contruction qui suit, notamment la recherche de démonstrations
dans $T$, est bien algorithmique, donc $e \mapsto h(e)$ est bien
calculable, et le théorème de récursion de Kleene s'applique.  On n'a
pas besoin de tester $A$, juste de savoir algorithmiquement
reconnaître une démonstration de « si $e$ termine, alors $A$ ».

\textbf{(2)} Si $g_A$ termine, c'est qu'il a trouvé une démonstration
dans $T$ de « si $g_A$ termine, alors $A$ ».  Mais le fait que $g_A$
termine implique qu'il y a une démonstration dans $T$ de
« $g_A$ termine » (en recopiant pas à pas la trace d'exécution, ou
l'arbre de calcul, de $g_A$).  Comme $T$ permet le \textit{modus
  ponens}, on en déduit qu'il y a une démonstration dans $T$ de $A$.

\textbf{(3)} Tout ce qui a été utilisé dans la démonstration de (2)
est de l'arithmétique élémentaire, et par ailleurs valable en logique
intuitionniste (aucun raisonnement par l'absurde n'a été tenu), donc
l'arithmétique de Heyting comme tous les systèmes proposés dans
l'énoncé prouve le point du (2).  Bref, $T$ prouve « si $g_A$ termine,
alors $T$ prouve $A$ ».

\textbf{(4)} On vient de voir au (3) que $T$ prouve « si $g_A$
termine, alors $T$ prouve $A$ » ; par conséquent, s'il prouve « si $T$
prouve $A$, alors $A$ », par \textit{modus ponens}, il prouve « si
$g_A$ termine, alors $A$ ».  Mais c'est exactement la preuve que $g_A$
cherche, donc $g_A$ termine.

\textbf{(5)} On a vu au (2) que si $g_A$ termine alors $T$ prouve $A$,
et au (4) que si $T$ prouve « si $T$ prouve $A$, alors $A$ » alors
$g_A$ termine.  On en déduit si $T$ prouve « si $T$ prouve $A$,
alors $A$ » alors $T$ prouve $A$.  C'est le théorème de Löb.
\end{corrige}

%

\exercice\ (${\star}{\star}{\star}{\star}$)\par\nobreak

Soit $T$ une théorie logique telle que\footnote{Voir la
note \ref{footnote-goedelian-theory}.}  l'arithmétique de Heyting
$\mathsf{HA}$, l'arithmétique de Peano $\mathsf{PA}$, Coq ou
$\mathsf{ZFC}$.

On considère le programme $h$ suivant : $h$ cherche une démonstration
dans $T$ du fait que $h$ termine, et s'il en trouve une, il termine
immédiatement.

\textbf{(1)} Expliquer pourquoi $h$ est bien défini.

(Remarquez que ce programme est « conciliant » : alors que dans la
démonstration du théorème de Gödel on a considéré un programme
« contrariant » qui fait le contraire de la démonstration qu'il a
trouvée, ici, si le programme trouve une preuve de son arrêt, il obéit
sagement.)

On va prouver que $h$ termine effectivement.

\textbf{(2)} Montrer que si $T$ prouve que $h$ termine, alors $h$
termine.

\textbf{(3)} Remarquer que le point précédent est lui-même démontrable
dans $T$.

\textbf{(4)} Considérer l'énoncé « le programme $h$ termine »,
appliquer le théorème de Löb énoncé à
l'exercice \ref{exercise-loebs-theorem}, et conclure.

\begin{corrige}
\textbf{(1)} Comme dans le (1) de
l'exercice \ref{exercise-loebs-theorem}, le programme est bien défini
par l'astuce de Quine et par le fait que vérifier une preuve dans $T$
est algorithmique.

\textbf{(2)} Par construction de $h$, si $T$ prouve que $h$ termine
alors $h$ finira par trouver cette preuve, et terminera.
(\emph{Attention :} ici il ne faut pas dire « si $T$ prouve que $h$
termine alors $h$ termine puisqu'on l'a prouvé » : on n'a pas fait
l'hypothèse que $T$ dit la vérité sur les arrêts de machines de
Turing.)

\textbf{(3)} Tout ce qui a été utilisé dans la démonstration de (2)
est de l'arithmétique élémentaire, et par ailleurs valable en logique
intuitionniste (aucun raisonnement par l'absurde n'a été tenu), donc
l'arithmétique de Heyting comme tous les systèmes proposés dans
l'énoncé prouve le point du (2).  Bref, $T$ prouve « si $T$ prouve que
$h$ termine, alors $h$ termine ».

\textbf{(4)} On a montré question (3) que $T$ prouve « si $T$ prouve
que $h$ termine, alors $h$ termine ».  Mais on a montré à
l'exercice \ref{exercise-loebs-theorem} que si $T$ prouve « si $T$
prouve que $h$ termine, alors $h$ termine », alors $T$ prouve que $h$
termine (c'est le théorème de Löb avec pour $A$ l'énoncé
« $h$ termine »).  Par conséquent, $T$ prouve que $h$ termine.  Donc,
par la question (1), $h$ termine effectivement.
\end{corrige}

\emph{Remarque :} on peut résumer la conclusion de cet exercice en
disant que « ceci est un théorème » est un théorème (mais la preuve,
comme on vient de le voir, n'est pas vraiment évidente).

%

\exercice\ (${\star}{\star}{\star}{\star}$)\par\nobreak

Dans cet exercice, on s'intéresse au programme $p$ suivant\footnote{On
pourra supposer ici les programmes écrits dans un langage de
raisonnablement haut niveau.} : il prend en entrée un entier qu'il
ignore ; il part de $N=42$, puis il énumère toutes les démonstrations
dans l'arithmétique de Peano dont la longueur (en nombre de symboles
sur un alphabet fini fixé) vaut au plus $10^{10^{100}}$, et, pour
chacune, si la \emph{conclusion} de la démonstration est une
affirmation du type « le $e$-ième programme termine sur l'entrée $i$ »
(i.e., $\varphi_e(i){\downarrow}$) avec $e$ et $i$ des entiers
naturels explicites, alors il exécute $\varphi_e(i)$ et si ce
programme termine, il ajoute le résultat (considéré comme un entier
naturel) à $N$.  Une fois que tout ceci est fait, il renvoie $N$.

\textbf{(1)} Commenter brièvement les aspects suivants du
programme $p$ : est-il très long à écrire ? est-il plus ou moins
évident qu'il termine ?

On désigne par « $\Sigma_1\mathsf{Sound}(\mathsf{PA})$ » (peu importe
ce que signifie « $\Sigma_1\mathsf{Sound}$ » ici) l'affirmation
suivante : « si l'arithmétique de Peano prouve
$\varphi_e(i){\downarrow}$, alors effectivement
$\varphi_e(i){\downarrow}$ ».  On ne cherchera pas (pour l'instant) à
se demander si $\Sigma_1\mathsf{Sound}(\mathsf{PA})$ est vrai, ou bien
démontrable.

\textbf{(2)} En \emph{admettant} l'affirmation
$\Sigma_1\mathsf{Sound}(\mathsf{PA})$ qu'on vient de définir, montrer
que le programme $p$ termine en temps fini.  Expliquer, de plus,
pourquoi l'arithmétique de Peano $\mathsf{PA}$ prouve ce fait.

\textbf{(3)} Sans entrer dans les détails, et toujours en admettant
$\Sigma_1\mathsf{Sound}(\mathsf{PA})$, expliquer pourquoi la valeur
renvoyée par $p$ (i.e., $\varphi_p(0)$) est gigantesque, par exemple
beaucoup \emph{beaucoup} plus grande que $10^{10^{1000}}$ itérations
de la fonction $n \mapsto A(n,n,n)$, où $A$ est la fonction
d'Ackermann, en partant de $10^{10^{1000}}$ (\emph{indication :}
esquisser comment on pourrait écrire un programme $q$ qui calcule
cette dernière valeur, et comment on prouverait que ce programme $q$
termine).  Quelle est la partie la plus longue dans l'exécution de
$p$ : l'énumération des démonstrations ou l'exécution des programmes
dont on a prouvé l'arrêt ?

\textbf{(4)} Toujours en admettant
$\Sigma_1\mathsf{Sound}(\mathsf{PA})$, montrer que la démonstration la
plus courte de l'arrêt de $p$ dans l'arithmétique de Peano est de
longueur supérieure à $10^{10^{100}}$.

En fait, $\mathsf{ZFC}$ prouve l'affirmation
$\Sigma_1\mathsf{Sound}(\mathsf{PA})$ énoncée plus haut, et la
démonstration n'est pas terriblement longue (on ne rentre pas dans les
détails ici).

\textbf{(5)} Conclure que la démonstration du fait que $p$ termine,
bien que démontrable dans l'arithmétique de Peano, est beaucoup plus
long à démontrer que dans $\mathsf{ZFC}$.

\begin{corrige}
\textbf{(1)} Le programme $p$ n'est pas très long à écrire : il s'agit
de calculer $10^{10^{100}}$, ce qui n'est pas difficile, puis
d'effectuer une boucle finie sur tous les entiers pouvant coder une
chaîne de longueur au plus $10^{10^{100}}$, de tester si chacun est
une démonstration valable dans $\mathsf{PA}$ (ce qui est algorithmique
et pas terriblement difficile) avec pour conclusion
$\varphi_e(i){\downarrow}$ et, le cas échéant, d'utiliser
l'interpréteur universel pour simuler l'exécution de $\varphi_e(i)$ et
l'ajouter à $N$, puis finalement renvoyer $N$.  Aucune de ces
étapes n'est très difficile, et on peut certainement dire que le
programme est de longueur bien inférieure à $10^{100}$, par exemple.

Quant à sa terminaison, la seule partie qui n'est pas évidente est le
fait que $\varphi_e(i){\downarrow}$.  On en a une preuve dans
$\mathsf{PA}$ (par construction du programme !), mais encore faut-il
que $\mathsf{PA}$ dise la vérité, ce qui constitue précisément
l'hypothèse $\Sigma_1\mathsf{Sound}(\mathsf{PA})$ introduite par
l'énoncé.

\textbf{(2)} Une fois admise l'hypothèes
$\Sigma_1\mathsf{Sound}(\mathsf{PA})$, il est évident que $p$ termine,
puisqu'il exécute $\varphi_e(i)$ une fois qu'il a une preuve dans
$\mathsf{PA}$ que $\varphi_e(i){\downarrow}$, donc toutes ces
exécutions terminent bien par $\Sigma_1\mathsf{Sound}(\mathsf{PA})$.
Tout le reste est une boucle finie, donc termine certainement.

Comme on a vu en cours que si un programme s'arrête, ce fait est
prouvable dans $\mathsf{PA}$ (en transformant une trace d'exécution
complète en démonstration), c'est le cas pour notre programme $p$.

\textbf{(3)} Soit $q$ le programme qui implémente la fonction
d'Ackermann et s'en sert pour calculer $10^{10^{1000}}$ itérations de
la fonction $n \mapsto A(n,n,n)$ en partant de $10^{10^{1000}}$ : ce
programme termine car la fonction d'Ackermann est bien définie.  Le
\emph{fait} que ce programme termine est facile à prouver, et cette
preuve, qui repose sur des considérations arithmétiques (des
récurrences imbriquées) est menable dans $\mathsf{PA}$ ; de plus, elle
n'est pas terriblement longue, certainement bien plus courte que
$10^{100}$.  Par conséquent, $p$ va, entre autres, tomber sur cette
preuve au cours de son énumération, donc il va exécuter $p$, donc
ajouter son résultat à $N$.  Donc la valeur de retour de $p$ est
supérieure à celle de $q$, et même considérablement plus grand puisque
$q$ va ajouter le résultat de tous les calculs de ce genre dont on
peut prouver la terminaison en moins de $10^{10^{100}}$ symboles.

Au cours de l'exécution de $p$, on doit parcourir toutes les
démonstrations de longueur au plus $10^{10^{100}}$ : cette énumération
est certes très longue (en gros une exponentielle de plus), mais
complètement minuscule par rapport au temps d'exécution des programmes
car on vient de voir qu'il y en a qui calculent des nombres
considérablement plus gigantesques (et parmi les programmes il y en a
qui calculent ce nombre puis attendent ce nombre d'étapes).

\textbf{(4)} On a vu en (2) que $\mathsf{PA}$ prouve que $p$ termine,
mais il n'est pas possible qu'il le prouve en moins de $10^{10^{100}}$
symboles, car si tel était le cas, $p$ ferait partie des programmes
exécutés lors de la boucle de $p$, donc le résultat ferait partie de
ceux ajoutés à $N$, donc on aurait $N>N$.

\textbf{(5)} Le fait que $p$ termine est prouvable dans $\mathsf{ZFC}$
grâce au raisonnement expliqué en (2), et la preuve n'est pas très
longue, certainement bien plus courte que $10^{100}$ symboles, même
quand on ajoute la preuve de $\Sigma_1\mathsf{Sound}(\mathsf{PA})$ qui
n'est elle-même pas bien longue d'après l'énoncé.  Comme on vient de
voir que la plus courte preuve de la terminaison de $p$ dans
$\mathsf{PA}$ est, pour sa part, plus longue que $10^{10^{100}}$
symboles, la preuve dans $\mathsf{ZFC}$ est beaucoup plus courte.
\end{corrige}


%
%
%

\section{Divers}

\exercice\label{exercise-dragon-riddle}\ (${\star}{\star}{\star}{\star}{\star}$)\par\nobreak

Vous êtes dans un donjon.  Devant vous se trouvent trois portes,
étiquetées $A,B,C$.  Derrière l'une de ces trois portes, mais vous ne
savez pas laquelle, se trouve un dragon, qui vous dévorera si vous
l'ouvrez ; les deux autres portes, qu'on appellera « sûres » dans la
suite mènent à la sortie du donjon (avec un trésor à la clé).  Votre
but est d'ouvrir une porte sûre (peu importe laquelle).

Sur chaque porte est affiché un programme ($p_A,p_B,p_C$) qui
constitue un indice pour trouver une porte sûre.  Plus précisément, le
programme affiché sur chaque porte prend une entrée $q$ et va,
\emph{sous certaines conditions} terminer et renvoyer l'étiquette
d'une des deux autres portes qui soit sûre.  Plus exactement, si on
suppose que $X$ est l'étiquette de la porte avec le dragon :
\begin{itemize}
\item Le programme $p_X$ qui est affiché sur la porte au dragon va
  \emph{toujours} terminer (quelle que soit l'entrée $q$ qu'on lui a
  passée) et renvoyer l'étiquette d'une des deux autres portes $Y,Z$
  (qui sont toutes les deux sûres puisqu'il n'y a qu'un seul dragon) :
  $\varphi_{p_X}(q) {\downarrow} \in \{Y,Z\}$ quel que soit $q$ (mais
  noter que le résultat peut dépendre de $q$).
\item Le programme $p_Y$ qui est affiché sur une porte sûre $Y$ va
  terminer et renvoyer l'étiquette de l'autre porte sûre $Z$, mais
  \emph{à condition} qu'on lui ait passé comme argument un programme
  $q$ (sans argument) qui fasse exactement ça (i.e., à condition que
  $q$ lui-même termine et renvoie $Z$) : si
  $\varphi_{q}(0){\downarrow} = Z$ alors\footnote{Le $0$ est mis pour
  une absence d'argument.}  $\varphi_{p_Y}(q) {\downarrow} = Z$ (dans
  tout autre cas, on n'a aucune garantie sur $\varphi_{p_Y}(q)$ : il
  pourrait ne pas terminer, renvoyer une étiquette fausse, ou renvoyer
  complètement autre chose).
\end{itemize}

Comment pouvez-vous utiliser ces indications pour ouvrir (de façon
certaine, et même algorithmique d'après $p_A,p_B,p_C$) une porte
sûre ?

(\emph{Indication :} on pourra utiliser
l'exercice \ref{exercise-good-bad-choice-lemma}(1).)

Question subsidiaire (indépendante) : montrer que l'énigme ci-dessus
serait insoluble si au lieu d'avoir des \emph{programmes} affichés sur
les portes on avait des tableaux de valeurs (tableaux donnant un
résultat « $A$ », « $B$ », « $C$ » ou n'importe quoi en fonction d'une
entrée « $A$ », « $B$ » ou « $C$ ») avec les mêmes contraintes (i.e. :
le tableau sur la porte du dragon donne l'étiquette d'une des deux
autres portes quelle que soit l'entrée, et le tableau sur une porte
sûre donne l'étiquette de l'autre porte sûre si l'entrée est justement
l'étiquette de l'autre porte sûre).

\begin{corrige}
Traitons d'abord la question subsidiaire.  Si les tableaux sont les
suivants :
\[
\begin{array}{c@{\hskip 2em}c@{\hskip 2em}c}
p_A[A] = B & p_A[B] = B & p_A[C] = C\\
p_B[A] = A & p_B[B] = C & p_B[C] = C\\
p_C[A] = A & p_C[B] = B & p_C[C] = A\\
\end{array}
\]
alors chacun des tableaux vérifie les deux contraintes (il renvoie
toujours l'étiquette d'une des deux autres portes, et si on le
consulte sur l'étiquette d'une des deux autres portes, il renvoie
celle-ci) ; or ces tableaux sont complètement symétriques (ils sont
invariants par les permutations cycliques de $A,B,C$) donc ils ne
peuvent pas permettre de déduire l'information de l'emplacement du
dragon qui pourrait être derrière n'importe quelle porte.

Il s'agit donc de faire quelque chose avec les programmes qu'on ne
peut pas faire avec de simples tableaux.  Ceci suggère d'utiliser le
théorème de récursion de Kleene.

En s'inspirant de l'exercice \ref{exercise-good-bad-choice-lemma}(1),
si $Z$ est une des trois portes et $X,Y$ les deux autres, on va
définir le programme $q_Z$ suivant : il invoque le programme $p_Z$ sur
$q_Z$ lui-même et ensuite, si $p_Z$ termine et renvoie $X$ ou $Y$, il
échange $X$ et $Y$ (autrement dit, si $p_Z$ appelé sur $q_Z$ renvoie
$X$, alors il renvoie $Y$, et si $p_Z$ appelé sur $q_Z$ renvoie $Y$,
alors il renvoie $X$), tandis que dans tout autre cas il fait une
boucle infinie.  La construction de ce $q_Z$, et notamment le fait que
$q_Z$ fasse appel à lui-même dans sa définition, est justifiée par
l'astuce de Quine.  Si $Z$ est la porte au dragon, alors, d'après les
garanties qu'on a reçues, $p_Z$ invoqué sur $q_Z$ va forcément
terminer et renvoyer une des étiquettes $X,Y$ (les deux sont sûres).
Si $Z$ est une porte sûre, mettons pour fixer les idées que $X$ soit
la porte au dragon : alors $p_Z$ invoqué sur $q_Z$ ne peut pas
renvoyer $X$, car si tel était le cas, $q_Z$ renverrait $Y$ (puisqu'il
échange les deux réponses $X,Y$), et par les garanties qu'on a reçues,
$p_Z$ invoqué sur $q_Z$ doit renvoyer $Y$, contradiction.

Maintenant, on lance en parallèle $p_A$ sur $q_A$, $p_B$ sur $q_B$, et
$p_C$ sur $q_C$. D'après ce qui a été dit au paragraphe précédent,
aucun des trois ne peut renvoyer l'étiquette de la porte du dragon, et
l'un des trois (celui de la porte du dragon) devra finir par renvoyer
l'étiquette d'une porte sûre.  Il suffit donc d'attendre que l'un des
trois termine et renvoie une étiquette dans $\{A,B,C\}$, et, quand ça
se produit, ouvrir la porte en question.
\end{corrige}

%

\exercice\ (${\star}{\star}{\star}{\star}{\star}$)\par\nobreak

On considère la formule propositionnelle suivante (due à V. Plisko) :
\[
\begin{aligned}
&\Big((\neg\neg A \Leftrightarrow \neg B\land \neg C)
\land (\neg\neg B \Leftrightarrow \neg C\land \neg A)
\land (\neg\neg C \Leftrightarrow \neg A\land \neg B)
\\
& \quad \land \big((\neg A \Rightarrow (\neg B\lor \neg C)) \Rightarrow (\neg B\lor \neg C)\big) \\
& \quad \land \big((\neg B \Rightarrow (\neg C\lor \neg A)) \Rightarrow (\neg C\lor \neg A)\big) \\
& \quad \land \big((\neg C \Rightarrow (\neg A\lor \neg B)) \Rightarrow (\neg A\lor \neg B)\big) \Big) \\
\Rightarrow\;& \big(\neg A \lor \neg B \lor \neg C\big)
\end{aligned}
\]
(on notera la symétrie entre $A,B,C$ qui rend la formule moins
complexe qu'elle n'en a l'air).

Déduire de l'exercice \ref{exercise-dragon-riddle} que cette formule
est réalisable mais (d'après la question subsidiaire de cet exercice)
pas Medvedev-valide.

(\emph{Indication :} comme $A,B,C$ n'apparaissent qu'avec une négation
partout dans cette formule, leur contenu n'a pas d'importance, la
seule chose qui importe est qu'ils aient ou non des éléments — ou dans
le cas des problèmes finis, des solutions ; il s'agit alors de
reconnaître que la formule reflète exactement les termes de l'énigme.)

\begin{corrige}
Rappelons que, pour la réalisabilité propositionnelle, $\dottedneg P$
vaut $\mathbb{N}$ lorsque $P$ est vide, et $\varnothing$ lorsque $P$
n'est pas vide (et symétriquement, $\dottedneg\dottedneg P$ vaut
$\varnothing$ lorsque $P$ est vide, et $\mathbb{N}$ lorsque $P$ n'est
pas vide).

Montrons que la formule est réalisable : pour ça, on va chercher à en
produire un réalisateur $r$ (qui doit être le même quelles que soient
$A,B,C \subseteq \mathbb{N}$).  Ce réalisateur est un programme
prenant six entrées.  Les trois premiers sont des réalisateurs de
$(\dottedneg\dottedneg A \dottedlimp \dottedneg B\dottedland
\dottedneg C)$ et symétriquement, et ne nous intéresseront pas pour
leurs valeurs, ce sont simplement des promesses que si $A$ est
non-vide alors $B$ et $C$ sont vides et symétriquement : on retient
simplement de ces trois entrées la garantie qu'exactement une des
trois parties $A,B,C$ est non-vide (« a un dragon »).  On va appeler
$p_A,p_B,p_C$ les trois autres entrées (quitte à les modifier un tout
petit peu comme on va l'expliquer).  Par exemple, $p_A$ doit réaliser
$((\dottedneg A \dottedlimp (\dottedneg B\dottedlor \dottedneg C))
\dottedlimp (\dottedneg B\dottedlor \dottedneg C)$, c'est-à-dire que
si on lui passe en argument un réalisateur $q$ de $\dottedneg A
\dottedlimp (\dottedneg B\dottedlor \dottedneg C)$ il doit terminer et
renvoyer un réalisateur de $\dottedneg B\dottedlor \dottedneg C$ ; or
un réalisateur de $\dottedneg B\dottedlor \dottedneg C$ est la donnée
d'un élément de $\dottedneg B$ ou de $\dottedneg C$ ainsi que
l'information duquel on a pris : mais cela revient simplement à donner
l'étiquette d'un des deux ensembles $B$ ou $C$ qui est vide (« n'a pas
de dragon »).  Donc la contrainte sur $q$ est que si $A$ est vide il
renvoie l'étiquette d'un des deux ensembles $B$ ou $C$ qui est vide
(et si $A$ n'est pas vide, n'y a aucune contrainte sur $q$) ; et la
contrainte sur $p_A$ est que si on lui passe un tel $q$, il renvoie
lui-même l'étiquette d'un des deux ensembles $B$ ou $C$ qui est vide.
Et au final, notre programme $r$ doit renvoyer l'étiquette d'un des
trois ensembles $A,B,C$ qui est vide.

Ce sont donc exactement les termes de l'énigme de
l'exercice \ref{exercise-dragon-riddle} : le réalisateur $r$ recherché
pour la formule est le programme qui prend $p_A,p_B,p_C$ comme on l'a
dit, qui effectue le calcul décrit dans la solution
de \ref{exercise-dragon-riddle}, et renvoie l'étiquette qu'il a
trouvé.  La formule de Plisko est bien réalisable.

Néanmoins, cette formule n'est pas Medvedev-valide : en effet,
appelons $A,B,C$ trois problèmes finis ayant chacun un candidat et
dont un seul a une solution (« un dragon »).  Le problème décrit par
l'ensemble de la formule consiste toujours à trouver une solution de
l'énigme mais cette fois en interprétant les programmes comme des
tableaux finis (fonctions entre ensembles finis).  Or on a expliqué
dans la question subsidiaire de
l'exercice \ref{exercise-dragon-riddle} que l'énigme n'avait pas de
solution dans ces conditions : c'est dire que la formule n'est pas
Medvedev-valide.

(Notamment, la formule de Plisko n'est pas démontrable en calcul
propositionnel intuitionniste, par \emph{correction} de la sémantique
de Medvedev.)
\end{corrige}


%
%
%
\end{document}
